% !TeX program = lualatex
\documentclass[11pt,oneside,openany]{book}
\usepackage[paperwidth=178mm,paperheight=246mm,inner=22mm,outer=21mm,top=21mm,bottom=23mm,headheight=15pt,footskip=11mm]{geometry}
\usepackage{amsmath,amsthm,mathtools}
\usepackage{libertinus}
\AtBeginDocument{\renewcommand{\Box}{\mathord{\mdlgwhtsquare}}}
\usepackage[final]{microtype}
\usepackage{xcolor,graphicx,tikz,booktabs,array,tabularx,enumitem}
\usetikzlibrary{arrows.meta,positioning,calc}
\usepackage{titlesec,fancyhdr}
\definecolor{ink}{HTML}{142B3D}
\definecolor{marine}{HTML}{146464}
\definecolor{ochre}{HTML}{A34C28}
\definecolor{quiet}{HTML}{52636F}
\definecolor{pale}{HTML}{E4F2EF}
\usepackage[unicode=true,colorlinks=true,linkcolor=marine,citecolor=marine,urlcolor=marine,bookmarksnumbered,bookmarksopen=false,pdfdisplaydoctitle]{hyperref}
\usepackage{bookmark}
\hypersetup{pdftitle={Lorentzian Sewing: Geometry, Dynamics, and Observables},pdfsubject={An English graduate text on assembling classical and perturbative quantum field theories},bookmarksdepth=2}
\setcounter{tocdepth}{1}
\setcounter{secnumdepth}{1}
\raggedbottom
\setlength{\parindent}{1.15em}
\setlength{\emergencystretch}{2em}
\setlist[itemize]{topsep=.4em,itemsep=.18em,leftmargin=1.5em}
\setlist[enumerate]{topsep=.4em,itemsep=.24em,leftmargin=1.8em}
\titleformat{\chapter}[display]{\normalfont\huge\bfseries\color{ink}}{\normalfont\sffamily\small\bfseries\color{marine}\MakeUppercase{\chaptertitlename}\enspace\thechapter}{6pt}{}
\titleformat{\section}{\normalfont\Large\bfseries\color{ink}}{\thesection}{.65em}{}
\titleformat{\subsection}{\normalfont\large\bfseries\color{marine}}{\thesubsection}{.55em}{}
\titlespacing*{\chapter}{0pt}{12pt}{17pt}
\titlespacing*{\section}{0pt}{2ex plus .6ex minus .25ex}{.7ex}
\pagestyle{fancy}\fancyhf{}
\fancyhead[L]{\small\sffamily\color{quiet}Lorentzian Sewing}
\fancyhead[R]{\small\sffamily\color{quiet}\nouppercase{\leftmark}}
\renewcommand{\chaptermark}[1]{\markboth{\chaptername\ \thechapter}{}}
\renewcommand{\headrulewidth}{.3pt}
\fancyfoot[R]{\small\thepage}
\fancypagestyle{plain}{\fancyhf{}\fancyfoot[R]{\small\thepage}\renewcommand{\headrulewidth}{0pt}}
\theoremstyle{plain}\newtheorem{theorem}{Theorem}[chapter]\newtheorem{proposition}[theorem]{Proposition}\newtheorem{lemma}[theorem]{Lemma}\newtheorem{corollary}[theorem]{Corollary}
\theoremstyle{definition}\newtheorem{definition}[theorem]{Definition}\newtheorem{example}[theorem]{Example}\newtheorem{exercise}{Exercise}[chapter]
\theoremstyle{remark}\newtheorem*{remark}{Remark}
\begin{document}
\frontmatter
\begin{titlepage}\thispagestyle{empty}\setlength{\parindent}{0pt}
\vspace*{6mm}
{\color{marine}\sffamily\bfseries\small GRADUATE LECTURE NOTES}\par
\vspace{11mm}
{\fontsize{39}{41}\selectfont\bfseries\color{ink}Lorentzian\\[2pt] Sewing\par}
\vspace{4mm}
{\Large Geometry, Dynamics, and Observables\par}
\vspace{12mm}
\begin{center}
\begin{tikzpicture}[x=1cm,y=1cm,>=Stealth]
\fill[ink!7] (0,0) rectangle (4,4.0);\fill[pale] (4,0) rectangle (8,4.0);
\draw[thin,ink!30] (0,0) rectangle (8,4.0);
\draw[marine,very thick] (4,0)--(4,4.0);
\draw[ochre,very thick,->] (1.95,.35)--(4,2.4)--(5.32,3.72);
\fill[ochre] (1.95,.35) circle (2.3pt);\fill[ink] (4,2.4) circle (2.2pt);
\draw[ink!38,dashed] (4,2.4)--(2.68,3.72);
\node[ink] at (.72,3.48) {$R_-$};\node[marine] at (7.28,3.48) {$R_+$};
\node[below,marine] at (4,0) {$\Gamma$};
\end{tikzpicture}
\end{center}
\vfill
{\small\sffamily\color{quiet}Classical fields and perturbative quantum theory\par}
\end{titlepage}
\tableofcontents
\mainmatter
\chapter*{Introduction: Is there a wall here?}\label{ch:intro}
\addcontentsline{toc}{chapter}{Introduction: Is there a wall here?}\markboth{Introduction}{}
A pulse is emitted to the left of a plane and detected to its right. Draw a line on the spacetime diagram at the plane: nothing happens to the pulse. Now give the left and right halves of the diagram to two different physicists. Ask each for a retarded solution with the field held at zero on the line. Each calculation is legitimate, but the detector on the right reads zero. A line in a drawing has become a reflecting wall in an initial-boundary value problem. How can the two physicists use their answers to remove the wall without looking up the solution on the uncut spacetime?

They need one more experiment that each can perform independently. With no bulk source on their own side, prescribe a small history of the field at the cut and record its outward response. The two response functions say how to vary that common history until no uncompensated force acts at the plane. The value required on the plane at noon must not depend on what the source will do tomorrow. This last condition makes the experiment Lorentzian. We will calculate this example, including the first arrival at the detector, in Chapter~\ref{ch:08}.

\section*{Two experiments that run in opposite directions}
In a \emph{cutting test} the whole theory is known. Restrict its fields, equations and observables to chosen pieces; sew them back and compare. In a \emph{construction} the pieces are independently given. One must first find out whether their faces, bundles and fixed backgrounds can be identified. There might be no possible whole, or several different wholes. In particular, two local descriptions in different gauges need not have literally equal values on their common face. They need an identification of their fibers. Holding such an identification fixed is a choice of geometric problem, not the insertion of an interaction.

Once a choice has been made, the two experiments are related by a restriction map. On the left of the map put smooth solutions of an independently specified whole problem. On the right put regional solutions satisfying transmission equations obtained from the regional actions. Restriction is useful when it goes both ways: every such pair comes from one whole solution, with the corresponding statement about changes of gauge. Requiring equality of \emph{all} derivatives on the right would often make the assertion immediate. The physical question is whether the much smaller amount of information exchanged by the equations is sufficient. Chapter~\ref{ch:07} proves that it is for a class of nonlinear equations and constructs a counterexample when a branch assumption is dropped.

A detector asks a different question of the map. It measures the change of a functional of the field when a source is switched on. Its answer involves the retarded inverse of the linearized equations, not merely the set of solutions of the homogeneous equations. A loop that crosses the cut asks a third question. Even if the two parties agree on every bulk equation and every response, what must each record at its two exposed endpoints so that the loop can be measured at all? These are reasons to keep \emph{fields}, \emph{causal response} and \emph{observable algebras} separate long enough to construct a map between them.

\section*{What a region carries}
A region carries a spacetime with its labeled faces, the bundles on which the fields live, fixed couplings and backgrounds, and a local action including the terms that specify the physical boundary problem. The word \emph{open} will mean that the data on an artificial cut may vary. It will never mean that boundary conditions at infinity, or at a real wall, have been discarded. A \emph{transparent seam} means that, after matching, there is no physical sheet left there. A source confined to a sheet or a field propagating along it is allowed, but it specifies a different target theory. A jump in a derivative can be precisely what that theory demands.

We use signature $(-,+,\ldots,+)$. A coorientation chooses which side of a hypersurface is positive. A collar is a thin strip of coordinates next to a face. A null face has no unit metric normal; none is needed to integrate an oriented flux. On a space of solutions $\delta$ denotes variation of the fields, rather than spacetime differentiation. The sign convention for the integrated presymplectic form is $\omega=\delta\theta$, where $\theta$ includes its corner contribution. The sign of a causal propagator belongs to its operator: throughout, the difference for a specified Hessian is its retarded inverse minus its advanced inverse. Multiplying the wave equation by $-1$ also multiplies that inverse by $-1$.

\section*{How to use the book}
If you know the action principle and wave equations from QFT and GR, you can start at Chapter~\ref{ch:01}. No previous use of jet bundles, gauge groupoids or microlocal analysis is assumed. For a short first reading, follow one scalar field through the signed collar in Chapter~\ref{ch:03}, the open variation in Chapter~\ref{ch:04}, and the transmission and response problems in Chapters~\ref{ch:07} and~\ref{ch:08}. Then return to the gauge theory chapters: what was an equality of field components becomes an equality after transport, and some transformations that look like redundancies on a closed slice carry a measurable charge at its exposed end. Exercises throughout have short checks or hints. The appendices develop three tools at the points where precision matters most: orientation, trace and support, and products of singular quantum fields.

The quantum question has a useful control experiment: sew the classical linearized theory first and quantize its observables. This produces a comparison algebra without multiplying two Hilbert spaces. A successful regional quantization should reproduce it within its stated scope. Reproducing its commutators does not choose a state; reproducing it at free order does not renormalize a pair of vertices that approach the seam from opposite sides. The last part starts with these distinctions. No Euclidean continuation is required.

\part{Spacetime, bundles, and smooth fields}
\chapter{Sewing spacetime regions}\label{ch:01}

Cut a strip out of Minkowski space and identify its two spatial edges. A particle that leaves on the right can now return on the left. But \emph{when} does it return? An identification can include a time translation even when each edge carries the same induced metric. To decide whether the result is a spacetime, we first construct its points and smooth charts, then ask whether the metric and time orientation fit those charts. Only after these steps can we ask a global question about a particle's worldline.

\section{Faces and their orientations}

A smooth region $R$ of dimension $d\geq 2$ is a Hausdorff, second-countable manifold with corners. Near a point of \emph{depth} $k$ it has coordinates $(r_1,\ldots,r_k,y)$ in $[0,\epsilon)^k\times\mathbb R^{d-k}$, with all $r_i=0$ at the point. Smooth up to a corner means locally the restriction of a smooth function on an open Euclidean neighborhood. The connected pieces of the depth-$k$ locus are strata. We use regions with embedded faces: each local wall $r_i=0$ belongs to a specified closed boundary hypersurface, and distinct incident walls belong to distinct faces. Closed here means a closed subset of $R$, not necessarily compact or without boundary. For instance, the end of a timelike wall on a final spacelike cap belongs to both; it is not a third hypersurface. Labelling those incidences tells us which face is artificial and will disappear in a sewing, and which exterior face must survive.

Two strata closures $A,B$ meet \emph{cleanly} if their intersection is a submanifold with corners and $T_p(A\cap B)=T_pA\cap T_pB$ there. In our product charts, distinct incident faces are coordinate walls with independent defining covectors $dr_i$. This is stronger than cleanness alone. It specifies the sectors on which fields are initially defined. We distinguish face \emph{occurrences}: a single region can contribute both members of a pair, as the strip does.

Give $R$ an orientation. At the interior of a face $F$, an ordered basis $e_1,\dots,e_{d-1}$ is positive on $F$ if $(N_{\rm out},e_1,\dots,e_{d-1})$ is positive on $R$, where $dr(N_{\rm out})<0$ for an inward defining coordinate $r$. Changing this transverse vector without changing its outward sense does not change the orientation. Taking the boundary of one face and then of another interchanges the two outward directions if the order is reversed. This is the sign behind opposite incidences at a corner. It is not a criterion that the sectors there fill a neighborhood.

A Lorentzian region carries a smooth nondegenerate metric $g$ of signature $(-,+,\ldots,+)$ up to its faces. A time orientation selects one of the two components of the timelike cone continuously as future. At a hypersurface with nonzero conormal $\nu$, the face is timelike, spacelike, or null according as $g^{-1}(\nu,\nu)$ is positive, negative, or zero \cite{ON}. In the null case $\nu^\sharp=g^{-1}(\nu,\cdot)$ is itself tangent: $\nu(\nu^\sharp)=0$. No rescaling makes it a unit normal, or even transverse. A defining coordinate and the outward-first orientation still work. Neither depends on the existence of Gaussian normal coordinates.

\section{From paired collars to a quotient atlas}

Let $f:F_-\to F_+$ be a diffeomorphism of closed face occurrences, preserving the labels and incidences that are to be retained. For oriented sewing it must reverse the induced boundary orientation. Choose paired collars $c_\pm:[0,\epsilon)\times F_\pm\to R_\pm$, each a diffeomorphism onto a neighborhood of its face, with $c_\pm(0,p)=p$. In these coordinates each incident wall is the product of its intersection with $F_\pm$ and $[0,\epsilon)$. Widths may be chosen locally. The collars specify how to approach a face; $f$ does not specify that by itself. Form the set $M$ by identifying $p$ with $f(p)$ in the disjoint union of the regions. Give it the quotient topology and write $q$ for the quotient map. We identify \emph{only} the faces, not points at positive collar distance. If $y$ is a chart on $F_-$, the map
\[
 q(c_-(r,p))\longmapsto(-r,y(p)),\qquad
 q(c_+(r,f(p)))\longmapsto(r,y(p))
\]
has one value on the seam. Its inverse uses the minus collar for $s\leq 0$ and the plus collar for $s\geq 0$, where $s$ denotes the first target coordinate. The two outward directions are now $+\partial_s$ and $-\partial_s$, respectively. This explains the required reversal without comparing metric normals. Different choices of transverse parametrization need not give the same smooth coordinates by the piecewise identity.

Here is a useful explicit regularity test for finite diagrams. At a class with $k$ crossing internal walls and $l$ surviving exterior walls, take a small box
\[
 B=(-\epsilon,\epsilon)^k\times[0,\epsilon)^l\times V,
 \qquad V\subset\mathbb R^{d-k-l}\text{ open}.
\]
Require one disjoint regional chart for each sign vector $\sigma\in\{-1,1\}^k$, with possible occurrences in the same region. In each chart use $r_i,\rho_a\geq0$ and $y\in V$. The proposed map is
\[
 (r,\rho,y)_\sigma\longmapsto
 (\sigma_1r_1,\ldots,\sigma_kr_k,\rho,y)\in B.
\]
The boxes must contain all representatives over this neighborhood: the prescribed pairings join exactly neighboring sectors on their common walls, fixing the displayed remaining coordinates, and no pairing leaves the collection. Thus $k=1,l=1$ at a cut ending on an exterior wall gives two quadrants filling a half-plane; $k=2,l=0$ requires four quadrants filling a plane, as in Figure~\ref{fig:01-sectors}.

\begin{figure}[htbp]
\centering
\begin{tikzpicture}[x=.85cm,y=.75cm,>=Stealth,font=\small]
  \fill[ink!7] (-2.7,0) rectangle (0,2.15);
  \fill[pale] (0,0) rectangle (2.7,2.15);
  \draw[quiet!65] (-2.7,0) rectangle (2.7,2.15);
  \draw[marine,very thick] (0,0)--(0,2.15);
  \node at (-1.35,1.72) {$R_-$};
  \node at (1.35,1.72) {$R_+$};
  \draw[ochre,thick,->] (-.09,.6)--(-2.10,.6) node[midway,above] {$r_-$};
  \draw[ochre,thick,->] (.09,.6)--(2.10,.6) node[midway,above] {$r_+$};
  \draw[ink,->] (-2.7,-.35)--(2.85,-.35) node[right] {$s$};
  \draw[ink] (0,-.43)--(0,-.27);
  \node[below] at (0,-.43) {$0$};
  \node[below] at (-1.50,-.43) {$s=-r_-$};
  \node[below] at (1.50,-.43) {$s=r_+$};
  \begin{scope}[shift={(6.25,.9)}]
    \fill[pale] (-1.48,-1.16) rectangle (0,0);
    \fill[pale] (0,0) rectangle (1.48,1.16);
    \fill[ink!7] (-1.48,0) rectangle (0,1.16);
    \fill[ink!7] (0,-1.16) rectangle (1.48,0);
    \draw[quiet!65] (-1.48,-1.16) rectangle (1.48,1.16);
    \draw[marine,thick,->] (-1.48,0)--(1.77,0) node[right] {$s_1$};
    \draw[marine,thick,->] (0,-1.16)--(0,1.49) node[right] {$s_2$};
    \fill[ink] (0,0) circle (1.5pt);
    \node at (-.76,.59) {$(-,+)$}; \node at (.76,.59) {$(+,+)$};
    \node at (-.76,-.59) {$(-,-)$}; \node at (.76,-.59) {$(+,-)$};
  \end{scope}
\end{tikzpicture}
\caption{Left: the inward collar parameters point away from the seam, but $s$ increases in just one direction. Right: the four choices of $(\operatorname{sign}s_1,\operatorname{sign}s_2)$ fill a neighborhood of the central point. Along each open ray there are exactly two sides; attaching a third sheet there would produce a branch rather than a coordinate chart.}
\label{fig:01-sectors}
\end{figure}

The collars in different boxes must also be compatible: the sectorwise coordinate changes, written with signed internal variables, must be restrictions of one smooth change with smooth inverse, up to the exterior walls. Product collars using the same normal variables and the same tangential reparametrization on both sides are a direct way to meet this test. For more general changes the condition has to be checked; smoothness on each closed sector alone does not imply it across their junction. These tests involve only the given charts and pairings, not an ambient spacetime in which the regions have already been embedded.

\begin{proposition}\label{prop:sewn-atlas}
For finitely many regions and finitely many closed-face pairings satisfying the sector and overlap tests, the quotient is a Hausdorff, second-countable smooth manifold with corners. Its local depth at the displayed point is $l$.
\end{proposition}
\begin{proof}
Each collection of boxes is open and saturated (a union of entire equivalence classes). Its signed map is continuous by the quotient topology. The inverse is continuous by pasting the inverse coordinate maps on the finitely many closed sectors of $B$. These homeomorphisms, including the ordinary charts off the seams, give the stipulated smooth atlas.

For separation let $X$ be the disjoint union before sewing. A class has at most $2^d$ representatives by the sector test. Any two of them are connected by a word in the pairings and their inverses of length at most $2^d-1$: erase repeated vertices from a path. There are finitely many such words. Each is a homeomorphism between closed subsets of $X$, since compositions of the partial homeomorphisms restrict to closed domains and ranges. For closed $C\subset X$, its saturation is therefore closed: it is the finite union of the images of $C$ intersected with these domains, including the identity word. Distinct finite classes in the Hausdorff space $X$ have disjoint open neighborhoods $U,V$. Replacing $U$ by $X\setminus\operatorname{Sat}(X\setminus U)$, and likewise for $V$, gives disjoint saturated open neighborhoods, which descend to $M$. Finally $X$ is second-countable and hence Lindel\"of (every open cover has a countable subcover). Its continuous image $M$ is Lindel\"of as well. Select countably many of the open chart images covering $M$; the union of their countable bases is a basis for $M$.
\end{proof}

Closedness matters. If two copies of a face segment are paired except at a common limiting endpoint, the two unpaired endpoints cannot be separated by disjoint open sets in the quotient. The atlas construction is in the smooth category of \cite{Lee}; it does not smooth discontinuous physical data. A fixed metric must have coefficients with matching derivatives of every order in the signed charts, not merely matching pullbacks to the face. For example, $dr^2-(1+r)^2dt^2$ in both inward collars gives the common induced metric $-dt^2$, but the assembled $g_{tt}=-(1+|s|)^2$ is not differentiable. Future cones must also agree under the identification if a time orientation is specified. For a dynamical metric the manifold is constructed first; it is each proposed metric history that must meet these requirements.

The product-box test is sufficient, not necessary. For example three planar wedges of opening $2\pi/3$ can tile a normal disk around their common vertex. Each open ray still has just two sides, unlike the branched example in Figure~\ref{fig:01-sectors}. One useful extension is to specify charts from the incident pieces onto closed convex sectors of a regular fan in a ball (or an exterior corner model). The sectors must tile the model, be identified exactly on their common faces, and their changes of coordinates must extend to smooth changes on the assembled neighborhoods. Assume finitely many closed pairings and a uniform bound on the number of representatives in a class. The proof just given works with that bound in place of $2^d$. To test a smooth field in such a chart, express derivatives in the common Cartesian coordinates, not in polar coordinates singular at the vertex. Chapter~\ref{ch:03}'s pasting proof applies because each coordinate segment splits into finitely many sector segments. If the link of the junction does not bound the required local model, the proposed output is a different, singular space; the smooth-manifold conclusions do not apply.

\section{Worked check: a time shift on a spatial seam}

On $\mathbb R_t\times[0,L]_x$ take $L>0$, $g=-dt^2+dx^2$ and orientation $dt\wedge dx$. Identify $(t,0)$ with $(t+\tau,L)$. Contracting the orientation form with $-\partial_x$ at $0$ gives $dt$, while contracting it with $+\partial_x$ at $L$ gives $-dt$. Hence the shift has the correct boundary sign. It is a translation of the flat collar charts; the metric and future direction $\partial_t$ sew for \emph{every} real $\tau$.

On the covering plane the translation vector is $(\tau,L)$, of squared length $L^2-\tau^2$. The straight segment joining translated points projects to a closed timelike curve if $|\tau|>L$, and to a closed null curve if $|\tau|=L$. If $|\tau|<L$, put $b=\tau/L$ and boost to $(T,X)=(t-bx,x-bt)/\sqrt{1-b^2}$. The identification is now $(T,X)\sim(T,X+\sqrt{L^2-\tau^2})$. The circles $T=\mathrm{constant}$ are Cauchy surfaces: every inextendible causal curve meets each exactly once. Indeed $T$ increases along a future causal curve and $|dX/dT|\leq1$, so a finite endpoint of its $T$ range would give a limit on the spatial circle and permit extension. Thus this last quotient is globally hyperbolic, whereas the first two are not even causal. Smooth local sewing has not decided between them.

\begin{exercise}
Let $R=\{u\geq0\}$ in two dimensions with $g=-du\otimes dv-dv\otimes du$ and orientation $du\wedge dv$. Find the metric dual of $du$ and the induced orientation at $u=0$. Why are the two answers obtained with different vectors?

\emph{Answer.} The inverse metric gives $(du)^\sharp=-\partial_v$, a null tangent generator. An outward transverse vector is $-\partial_u$; its contraction with the orientation form is $-dv$. The first calculation uses metric orthogonality, the second uses only transversality and the side occupied by $R$.
\end{exercise}
\chapter{Bundles and connections at a cut}\label{ch:02}

Identifying two spacetime points does not tell us which vectors in their two fibers represent the same field value. For a charged field, even the statement that two complex numbers agree presupposes a comparison of frames. We will construct the fiber identification first. It will tell us how to compare matter fields and connections, and it will leave a measurable choice on a spatial circle.

\section{Lifting a face map}

Let $G$ be a Lie group, possibly disconnected. A principal $G$-bundle $P\to R$ has a right $G$-action, free and transitive on each fiber, and smooth local product charts. A local section $e_i$ chooses a frame: each point in the fiber has a unique expression $e_i(x)g$. An associated field with typical fiber $W$ and left $G$-action takes values in $P\times_G W$, with $[p g,w]=[p,g w]$. This construction works for a vector representation, but does not require $W$ to be linear.

For the face map $f:F_-\to F_+$ of Chapter~\ref{ch:01}, seek an equivariant isomorphism $D:P_-|_{F_-}\to P_+|_{F_+}$ covering $f$. Thus the test is whether $P_-|_{F_-}$ and $f^*(P_+|_{F_+})$ are isomorphic, not whether each can be trivialized on a small patch. To carry it out, choose frames on corresponding patches and write
\[
 e_j^\pm=e_i^\pm g_{ij}^\pm,\qquad
 D(e_i^-(y))=e_i^+(f(y))h_i(y).
\]
Evaluating $D(e_j^-)$ first in the $i$ frame and then in the $j$ frame gives
\begin{equation}
 h_i g_{ij}^-=(g_{ij}^+\circ f)h_j .\label{eq:02-clutch}
\end{equation}
Conversely, smooth $G$-valued $h_i$ satisfying \eqref{eq:02-clutch} define $D$ by $e_i^- g\mapsto e_i^+ h_i g$ at corresponding points. The equation is exactly what makes the definition independent of $i$ \cite{Steenrod}. Under a passive frame change $e_i^\pm\mapsto e_i^\pm k_i^\pm$ the same $D$ has matrix $(k_i^+\circ f)^{-1}h_i k_i^-$. No bundle has changed in this calculation.

Several bundles can be present. If a map $b_\pm:E_\pm\to H_\pm$ is part of the structure, the two proposed lifts must satisfy $D_H b_-=b_+D_E$. For example, an identification of spin bundles must project to the specified identification of oriented orthonormal frames; an abstract isomorphism of the two total spaces is not enough. The relevant frame comparison acts on the full spacetime tangent space, including the transverse direction identified by the signed collars. In contrast, an associated bundle has an induced lift $[p,w]\mapsto[D(p),w]$ by equivariance. There is no second independent twist for each associated species unless the theory actually supplies such structure.

\section{Smooth clutching and the connection law}

Construct $P$ over the sewn $M$ by identifying $p$ on the seam with $D(p)$. To see smoothness, choose product descriptions of the two collar bundles as pullbacks of their face bundles. Parallel transport along the collar lines with any auxiliary connection gives such descriptions for an ordinary principal bundle. This is a choice of charts, not of the physical connection. In extended frames let $g_-,g_+$ be the group coordinates on each side. A common fiber coordinate over the signed base box is $h_i(y)g_-$ on $s\leq0$ and $g_+$ on $s\geq0$. The inverse uses $h_i(y)^{-1}$ on the minus sector. Equation~\eqref{eq:02-clutch} makes the transition to the next product frame a single smooth map. The right action is ordinary right multiplication throughout this box. Thus the quotient is locally $B\times G$ with principal fibers, not merely a set of identified field components.

At a joint choose compatible product descriptions for the whole diagram, or give actual smooth signed transitions realizing the chosen lifts. Write $D_{ij}$ for transport from frame $j$ to frame $i$. On a common triple overlap one needs $D_{ij}D_{jk}=D_{ik}$ and $D_{ji}=D_{ij}^{-1}$. When only adjacent sectors have been paired at a four-sector vertex, the two routes around the square from one sector to the opposite one must agree there, with compatible smooth transitions. Otherwise going around a single point could identify $p$ with $p g$ for $g\neq 1$, destroying freeness. These are local conditions at the \emph{same} point, not an instruction to set the transport around a noncontractible spatial loop to one. Rigid structures, such as a prescribed framing or a fixed connection, require their own compatibility in the collar; an auxiliary product description need not preserve them.

A principal connection specifies which directions in $P$ are horizontal, so a base path can be lifted by parallel transport. Its Lie-algebra-valued form $\omega$ vanishes on horizontal vectors, returns $X$ on the vertical vector generated by $X$, and obeys $R_g^*\omega=\operatorname{Ad}_{g^{-1}}\omega$ under right multiplication. Its local potential is $A=e^*\omega$. Differentiating a change of section $e\mapsto eh$ gives $(eh)^*\omega=\operatorname{Ad}_{h^{-1}}(e^*\omega)+h^{-1}dh$: the first term differentiates $e$, the second is the vertical motion of $h$. In our common signed chart $e_-=e_+h$, with base pullbacks understood. Solving for $A_+$ gives
\[
 A_+=\operatorname{Ad}_h A_--(dh)h^{-1} .
\]
In a vector representation this is just $(d+A_+)(h z_-)=h(d+A_-)z_-$ for every matter section with components $z_+=h z_-$. The connection law is affine, not tensorial: only the difference of two connections transforms homogeneously. In particular a $U(1)$ change $h=e^{i c s}$ gives $A_+=-i c\,ds$ if $A_-=0$, even though $h=1$ on the face. The displayed law relates two descriptions of \emph{one} connection. To test two independently given regional connections, transform their one-sided coefficients into one signed frame, including the $ds$ component and the derivatives of the realized frame change, and compare the seam jets there. Equality of tangential pullbacks alone is weaker. We make the smoothness criterion precise in Chapter~\ref{ch:03}.

\section{Which changes identify the same sewing?}

A marked sewing remembers the labelled regions, the chosen cut and the prescribed exterior data. An isomorphism may reparametrize a labelled region by an allowed automorphism of the \emph{whole} region; it does not erase the decomposition. Specify which exterior fields, sources and frames are fixed, and which may be transported. Let $A_a$ consist of the actual automorphisms of region $a$ respecting these choices and the fixed structure. A symmetry defined only on a cut need not extend into the region subject to those requirements.

For a realized smooth sewing, write $D$ also for its full seam transition data (base and bundle maps with their signed Taylor coefficients). Write $\widehat u_a$ for the corresponding restriction of $u_a\in A_a$. An isomorphism from data $D$ to data $D'$ satisfies
\begin{equation}
 \widehat u_+ D=D'\widehat u_- .\label{eq:02-mark}
\end{equation}
Indeed a point transported across the seam and then reparametrized must agree with one reparametrized first and transported across the new seam. The same equality on derivatives is what makes the piecewise map smooth; the descent argument of the next chapter proves this for maps as well as fields. Conversely a smooth isomorphism preserving the marks restricts to actual regional maps satisfying \eqref{eq:02-mark}. Thus the classes are the orbits of admissible $D$ under $D\mapsto\widehat u_+D\widehat u_-^{-1}$. The stabilizer of $D$ consists of the \emph{pairs} $(u_-,u_+)$ satisfying \eqref{eq:02-mark} with $D'=D$. For instance, allowed maps supported away from the cut can be invisible in $\widehat u_a$ and still be distinct automorphisms. Several faces of one region use restrictions of the \emph{same} map. This describes marked sewings; an isomorphism allowed to move or forget the cut need not restrict to such maps.

\section{Worked check: two signs on a cylinder}

Build $\mathbb R_t\times S^1$ from two connected spatial strips, each carrying the trivial principal $G=\{+1,-1\}$ bundle. Fix the base and impose no exterior framing. In order around the circle let the seam lifts be $a$ from strip 1 to strip 2 and $b$ from strip 2 back to strip 1. A smooth vertical automorphism on a connected strip is a \emph{single constant} $u_i=\pm1$ because the group is discrete. Thus
\[
 a'=u_2 a u_1^{-1},\qquad b'=u_1 b u_2^{-1},\qquad b'a'=ba .
\]
Lift a spatial circuit at fixed $t$, starting with fiber element $1$. It stays constant in each strip frame, jumps to $a$ at the first seam and returns as $ba$ at the second. Its holonomy $H=ba$ therefore has a direct path-lifting meaning. Choosing $u_1=1,u_2=a^{-1}$ puts the pair in the form $(1,H)$. There are exactly two classes, each with diagonal stabilizer $u_1=u_2$. In the associated sign representation a field on the covering line obeys $\phi(x+L)=H\phi(x)$ for circumference $L$. Allowing an unrelated sign at each end of a strip would incorrectly erase the antiperiodic possibility. This is the holonomy of a discrete covering bundle; in a connected group a connection can have nontrivial holonomy even on a topologically trivial bundle.

\begin{exercise}
On the boundary of a disk prescribe a $U(1)$ gauge transformation $v(\theta)=e^{im\theta}$, $m\in\mathbb Z$. For which $m$ does it extend to a smooth gauge transformation on the disk? What does this imply for the winding of a transition between two trivial disk bundles?

\emph{Hint.} For an extension $u$, the form $u^{-1}du$ on the disk is closed. Stokes' theorem gives $0=\oint v^{-1}dv=2\pi i m$. The zero case extends by the constant map. Frame changes from either disk have zero boundary winding and cannot change the winding of the transition.
\end{exercise}
\chapter{Smooth fields and gauge transformations across a seam}\label{ch:03}

A Lagrangian can depend on just a field and its first derivatives without making every once-differentiable field a smooth configuration. These are two distinct choices: the \emph{order of a local rule} and the \emph{regularity of its inputs}. The distinction matters at a cut. A finite list of measurements can suffice to evaluate a density there, yet fail to decide whether a pair of independently prepared fields is the restriction of one smooth field.

\section{What a finite jet measures}

Let $E\to M$ be a field bundle on the spacetime constructed in Chapters~\ref{ch:01} and~\ref{ch:02}. A configuration is a smooth section, with any stated exterior conditions. \emph{Off shell} means that no equations of motion have been imposed. Two local sections have the same $k$-jet at $p$ if, in one bundle chart, their values and all their partial derivatives through total order $k$ agree at $p$. The chain rule makes the definition independent of the chart. The set of such jets as $p$ varies is a bundle $J^k E\to M$; $j^k_p\phi$ records derivatives of an actual section $\phi$. A local density of order $k$ is a smooth map on these jets to densities of the base, possibly using specified background jets as well \cite{GMS}.

For example, if $\phi$ takes values in a curved target with metric $H$, the sigma-model density in coordinates is
\[
 -\left[\tfrac12 g^{\mu\nu}(x) H_{AB}(\phi(x))\,
       \partial_\mu\phi^A\partial_\nu\phi^B+V(\phi(x))\right]
       \,d\mathrm{vol}_g .
\]
Here $d\mathrm{vol}_g$ is the metric volume density. The expression depends on $j^1\phi$ even if $H$ and $V$ depend nonlinearly on $\phi$. Under a change of fiber coordinates $v=F(s,y,u)$ in a common base chart, for a scalar fiber coordinate $u$ one has
\[
 v_s=F_s+F_u u_s,\qquad
 v_{ss}=F_{ss}+2F_{su}u_s+F_{uu}u_s^2+F_u u_{ss} .
\]
Thus transforming derivatives is not just multiplying each one by the linearized field map. When the base coordinates change as well, one first reexpresses the section over the new base using the inverse coordinate map. For independent regional rules, the one-sided coefficient functions on the jet bundle must fit smoothly at the seam for all values of the jet variables. This includes their signed base derivatives, not only their values on one favored field. For instance, two massive scalar Lagrangians with different $m^2$ agree on the zero field but have incompatible coefficients of $\phi^2$.

\section{Signed traces and mixed derivatives}

Write fields in a single realized seam frame, doing any tensor, affine or nonlinear coordinate change \emph{before} comparing components. Let $u_-(r_-,y)$ and $u_+(r_+,y)$ denote the resulting coefficients, with $s=-r_-$ on the minus side and $s=r_+$ on the plus side. For a one-form, for example, the coefficient of $dr_-$ acquires a minus sign when rewritten as a coefficient of $ds$. The smooth trace test uses the derivatives
\begin{equation}
 (-1)^m\partial_y^\alpha\partial_{r_-}^m u_-(0,y)
   =\partial_y^\alpha\partial_{r_+}^m u_+(0,y),
 \qquad m\geq0,\quad\text{all multi-indices }\alpha .\label{eq:03-jet}
\end{equation}
If the normal traces agree as smooth \emph{functions} of $y$, their tangential derivatives agree by differentiation. If $s$ is already used on both sides, its derivatives are equal with no further $(-1)^m$. There is no demand that the two fields be analytic, or that they agree on collars of positive width. A scalar equal to zero for $s\leq0$ and to $e^{-1/s^2}$ for $s>0$ passes every test: each right derivative is a polynomial in $1/s$ times the exponential and tends to zero.

At a crossing $s_i=\sigma_i r_i$ with $\sigma_i=\pm1$, put all incident components $u_\sigma$ in the same fiber chart by the coherent lifts. At every shared point compare
\[
 \left(\prod_{i=1}^k \sigma_i^{\beta_i}\right)
 \partial_{r_1}^{\beta_1}\cdots\partial_{r_k}^{\beta_k}
 \partial_w^\gamma u_\sigma
\]
for every pair of incident sectors and all multi-indices $\beta,\gamma$. Here $w$ comprises the remaining coordinates, including any exterior ones. This compares the actual derivatives, including mixed ones, independently of the order in which faces were reached. For sector fields smooth up to each closed wall, full normal trace equalities along the adjacent walls give the mixed equalities at their ends by tangential differentiation and continuity in these product charts. Bare corner values or separate pure derivatives at the vertex would not: $c_\sigma s_1s_2$ on each quadrant has vanishing pure derivatives at the vertex but mixed derivative $c_\sigma$.

Where exterior coordinates satisfy $\rho_a\geq0$, the definition of smoothness uses an ambient extension. The following local test makes this convention explicit: in each sector choose extensions across the \emph{exterior} walls to a slightly larger box and require that the extensions satisfy the same internal jet equalities there. It suffices that one compatible choice exists, not that arbitrary extensions work. This is necessary for a smooth sewn field, since a local ambient extension restricts to such a choice. It will be sufficient by the proof below. All tests are on actual regional fields or on such local extensions. Arbitrary assignments of Taylor coefficients on strata, without a common coordinate realization and extension conditions, are not their substitute.

\section{Descent of objects and arrows}

Fix the sewn bundles and the allowed local field and gauge operations; in particular these operations must be the same in corresponding frames. A vertical gauge transformation is an equivariant bundle automorphism over the identity. On $P$ it is a section of the group bundle $\operatorname{Ad}P=P\times_{\rm conj}G$: writing $U(e_i)=e_i u_i$, a change of frame conjugates $u_i$. It is an \emph{arrow} from $\phi$ to $\psi$ when $U$ takes $\phi$ to $\psi$. The collection of configurations with these invertible arrows is called a groupoid. Remembering the arrows retains, in particular, the transformations that fix a field.

\begin{theorem}[Off-shell descent]\label{thm:off-shell-descent}
In the finite product-sector atlas of Proposition~\ref{prop:sewn-atlas}, restriction identifies smooth configurations with tuples of regional configurations satisfying the signed, mixed and exterior tests above. It also identifies smooth vertical gauge arrows between two configurations with the regional arrows satisfying the tests in their group bundle. Impose the same exterior requirements in both descriptions. If a further nonlocal condition is specified, impose it on the joint tuple as well.
\end{theorem}
\begin{proof}
For fields, order zero gives a well-defined continuous piecewise section; this is finite closed-sector pasting in a bundle chart. For each multi-index $I$ paste the signed derivatives to a continuous coefficient $v_I$, with $v_0$ the original component. In a small open box a coordinate segment crosses finitely many walls. Apply the fundamental theorem of calculus on its pieces and add, obtaining, in direction $j$ with coordinate vector $e_j$,
\[
 v_I(x+h e_j)-v_I(x)=\int_0^h v_{I+e_j}(x+t e_j)\,dt .
\]
Here $I+e_j$ means incrementing the $j$th derivative count. A segment lying in a wall uses either incident smooth restriction; the traces coincide. Continuity gives $\partial_jv_I=v_{I+e_j}$, and induction proves smoothness of every order, including at an internal crossing. At a surviving exterior wall apply this argument to the extensions and then restrict. For a nonlinear fiber choose a chart about the common value and shrink the box by continuity. The bundle transition laws identify the resulting local sections on overlaps. Restriction is inverse to piecewise evaluation; a smooth section plainly has all the required traces.

For arrows use the same argument in a chart of $G$. In the frames of Chapter~\ref{ch:02} the seam condition at order zero is $u_+=h u_-h^{-1}$, since $U_+D(e_-)=e_+u_+ h$ and $DU_-(e_-)=e_+h u_-$. At higher orders one differentiates the \emph{transported} group-valued map. Thus the regional arrows give a single smooth section of $\operatorname{Ad}P$. Its inverse and its products are smooth by the group operations. Acting on the assembled field agrees with the assembled target because the two sections agree on each regional interior and are continuous. Conversely, restriction gives the original arrows with their composition, identities and stabilizers intact. This proves descent on arrows, not merely a bijection after taking gauge orbits.
\end{proof}

A theory with independent wall or corner fields applies the same argument to the bundles on those strata, with the given incidence conditions. Such a field need not be the trace of a bulk field. If gauge arrows can change the base, the analogous assertion concerns arrows preserving the labelled cut. Write the regional maps and their inverses in corresponding source and target charts and apply the same component proof to each; their composites restrict to the identity everywhere. A diffeomorphism moving the cut does not restrict to this kind of regional arrow. A smooth fixed metric or coupling must pass the test before it can be used as common background input. A dynamical metric, by contrast, is a section to which the test is applied along with its pointwise signature and any stipulated time orientation. A global causal restriction is a further condition on the assembled history, not a jet equality.

\section{Worked check: a nonlinear chart can hide a derivative jump}

Use $u>0$ in the minus chart and $v=\log u$ in the plus chart. Near the face take $u_-(r_-)=(1-r_-)^2$ and $v_+(r_+)=2\log(1+r_+)+a r_+^3$. In signed $s$ and common $v$ coordinates the left field is $2\log(1+s)$. Its first derivative at zero is $2$ and its second is $-2$: from the $u$ chart the latter is $u_{ss}/u-(u_s/u)^2=2-4$, not $u_{ss}=2$. Both right derivatives agree. The third derivative is $4$ on the left and $4+6a$ on the right. For $a\ne0$ the piecewise field is $C^2$ but not $C^3$, even though its first jet, and hence every common first-order density evaluated there, is unambiguous. For $a=0$ the two formulas give one smooth section. Replacing the cubic by $a s^{m+1}$ on the positive side defeats any proposed finite order $m$ off shell.

\begin{exercise}
Take a charged scalar identically zero on both sides in a trivial $U(1)$ bundle with identity seam transition. Consider regional gauge maps $u_-=1$ and $u_+=e^{i c s}$, $c\ne0$. They fix the scalar and agree on the face. Are they the restrictions of a smooth stabilizing arrow? What would they do to the zero connection?

\emph{Answer.} The signed first derivatives at the seam are $0$ and $ic$. So there is no such smooth arrow, although both regional scalars remain zero. With the $d+A$ convention their actions on the zero connection give $A'_-=0$ and $A'_+=-ic\,ds$, exposing the jump already in the values of the transformed one-forms.
\end{exercise}

\part{Classical field theory on an open region}
\chapter{Varying a region before and after opening a face}\label{ch:04}
Cut a vibrating string at a point and hold the displacement there equal to a
specified function of time. Each half now needs a force at the cut to follow
that history. When the two halves are rejoined, it is the \emph{sum} of the
required forces that must vanish. Setting each force to zero would describe
two free ends instead. We will obtain the corresponding statement for a field
from the first variation of its action. The calculation also tells us which
variables play the role of the displacement and where the forces are measured.

\section{The variation at prescribed data}
Let $R$ be an oriented Lorentzian slab of signature $(-,+,\ldots,+)$.
Use the same future slice orientation on the two spacelike ends, so that the
oriented boundary is $\Sigma_f-\Sigma_i+B+\Gamma$. The timelike face $B$
is a physical wall, possibly a regulator for a boundary at infinity; its
boundary condition and action are part of the theory. The timelike face
$\Gamma$ is an artificial cut. For simplicity the side faces meet the end
slices but not one another. Write $\partial_F\Sigma=\Sigma\cap F$ with the
orientation inherited as a boundary of $\Sigma$. We use $d$ for the spacetime
differential and $\delta$ for variation of fields at fixed spacetime point;
the latter commutes with spacetime differentiation in our formulas.

\begin{definition}
A regional problem at prescribed cut datum $j$ consists of a configuration
space with the physical condition on $B$ and a specified trace condition
$b_\Gamma(\phi)=j$ on $\Gamma$, together with a differentiable action on
that space. A trace assigns to a bulk field its selected boundary values;
these can include derivatives. A variation in this space satisfies
$\delta b_\Gamma=0$, not necessarily $\delta\phi|_\Gamma=0$.
\end{definition}

Dirichlet data fix the boundary field itself. Fixing a normal momentum
instead usually requires a different face term. To see why that term belongs
in the variation, write
\[
 S_R[j]=\int_R L+\int_B\ell_B+\int_\Gamma\ell_\Gamma[j],
 \qquad \delta L={\cal E}_a\,\delta\phi^a+d\Theta.
\]
Here $L$ is a spacetime top form, ${\cal E}_a$ are the Euler--Lagrange
densities and $\Theta$ is a spacetime form of one lower degree, linear in a
field variation. On each side face $F$ differentiability requires the
remainder $\Theta|_F+\delta\ell_F$ to vanish or to be a tangential derivative
$d_F C_F$ on the allowed variations. The corner potential $C_F$ has
spacetime degree $\dim R-2$ and is linear in the variation. If the wall
supports its own equation,
its coefficient is kept in the variation until that equation is imposed.
At fixed $j$, imposing the boundary conditions but not the bulk equations,
Stokes' theorem gives
\begin{align}
 \delta S_R[j]&=\int_R {\cal E}_a\,\delta\phi^a
              +\theta_{\Sigma_f}-\theta_{\Sigma_i},\nonumber\\
 \theta_\Sigma&=\int_\Sigma\Theta
         -\sum_{F=B,\Gamma}\int_{\partial_F\Sigma} C_F.
 \label{eq:04-endpoint}
\end{align}
The minus sign can be found without a metric: at the final corner,
$\partial F$ induces the opposite orientation to $\partial_F\Sigma_f$,
since the boundary of the oriented boundary vanishes. A tangential derivative
on a wall therefore contributes to the endpoint potential; it is not zero.

For a concrete corner, take a flat strip with a real scalar and add a kinetic
term $\tfrac12\kappa\int_{t_i}^{t_f}\dot q_B^2\,dt$ on one physical end of
the string, where $q_B=\phi|_B$. Its variation is
\[
 -\kappa\int_{t_i}^{t_f}\ddot q_B\,\delta q_B\,dt
      +[\kappa\dot q_B\,\delta q_B]_{t_i}^{t_f}.
\]
For the bulk kinetic term used below, the wall equation with free $q_B$ is
$-\partial_n\phi-\kappa\ddot q_B=0$, where $n$ points out of the strip.
Even on a solution of this equation, the endpoint term is still present: the
wall carries momentum $\kappa\dot q_B$. It must not be discarded with the
wall equation.

\section{Releasing a common trace}
To open $\Gamma$, allow $j$ to range over the admissible histories and keep
the condition on $B$ fixed. If $\mathcal C_R[j]$ is the configuration space
at prescribed $j$, this gives $\mathcal C_R^{\mathrm{open}}
=\bigcup_{j\;\mathrm{admissible}}\mathcal C_R[j]$; a history must still
respect any compatibility condition where the cut meets an end slice.
Remove the face action used \emph{only} to impose
the auxiliary closure, or explicitly undo its boundary Legendre transform.
Do not remove a term that defines the physical wall. Opening is not a
stationarity condition on one side of the cut. It makes the cut part of the
first variation visible as a linear functional of varying trace data $q$:
\[
 \delta S_R^{\mathrm{open}}=\int_R {\cal E}_a\delta\phi^a
 +\theta^{\mathrm{open}}_{\Sigma_f}-\theta^{\mathrm{open}}_{\Sigma_i}
 +\langle\Pi,\delta q\rangle_\Gamma.
\]
Here tangential derivatives of $\delta q$ have been integrated by parts;
the resulting corner contributions belong to $\theta^{\mathrm{open}}$.
For example, subtracting $\langle\Pi,q\rangle_\Gamma$ would turn the
face variation $\langle\Pi,\delta q\rangle_\Gamma$ into
$-\langle\delta\Pi,q\rangle_\Gamma$, suitable for fixing $j=\Pi$. Undoing
this subtraction restores the original face one-form; it does not set the
momentum to zero.
For a first-order scalar Lagrangian and Dirichlet closure, $q=j=\phi|_\Gamma$
and no artificial face action is needed. With fixed metric and volume form
$\epsilon$, the action density and its variation are
\begin{align*}
 L&=-\tfrac12(\nabla_\mu\phi\nabla^\mu\phi+m^2\phi^2)\epsilon,\\
 \delta L&=(\nabla_\mu\nabla^\mu\phi-m^2\phi)\delta\phi\,\epsilon
       -\nabla_\mu(\delta\phi\nabla^\mu\phi)\epsilon.
\end{align*}
On a timelike face the last term gives
$\langle\Pi,v\rangle_\Gamma=-\int_\Gamma (n^\mu\nabla_\mu\phi)v\,d\mu_\Gamma$,
where $v=\delta q$, $n$ is unit and outward, and $d\mu_\Gamma$ is the positive
induced measure. Thus $\Pi$ is a dual trace \emph{density} with the outward
sign already included. In arbitrary face coordinates $y$ it contains
$\sqrt{|\det\gamma|}\,|d^{\dim\Gamma}y|$, not just a component of the
normal derivative. Equivalently one can integrate the form
$-(\partial_n\phi)v\,\iota_n\epsilon$ with the outward face orientation.
Do not insert a second orientation sign when using the density formula.
One can even avoid normalizing a metric normal. For an outward conormal
$\nu$ and any transverse $N$ with $\nu(N)=1$, let
$\epsilon_\nu=(\iota_N\epsilon)|_\Gamma$. The outward surface density is
$\nu_\mu\,|\epsilon_\nu|$. Rescaling $\nu$ by a positive
function rescales $N$ inversely, and adding a tangential vector to $N$ has
zero effect on the pullback. This explains why the pairing, rather than a
bare coordinate component, is the invariant object.

Suppose a fixed smooth identification $D$ carries trace data on side 2 to
side 1 with smooth inverse. Continuity $q_1=D(q_2)$ is a restriction on
configurations, made before varying them. Differentiating $q_2=D^{-1}(q_1)$
gives $\delta q_2=T_{q_1}D^{-1}\,\delta q_1$; only for a linear map may one
use $D^{-1}$ itself on variations. At the matched data the dual transport is
defined by $\langle D^\vee\Pi_2,v\rangle_{\Gamma_1}
=\langle\Pi_2,T_{q_1}D^{-1}\,v\rangle_{\Gamma_2}$, including the change of
measure. Consequently, the sum of the two opened actions has the cut variation
\[
 \langle\Pi_1,\delta q_1\rangle_{\Gamma_1}
 +\langle\Pi_2,T_{q_1}D^{-1}\,\delta q_1\rangle_{\Gamma_2}
 =\langle\Pi_1+D^\vee\Pi_2,\delta q_1\rangle_{\Gamma_1}.
\]
Stationarity gives $\Pi_1+D^\vee\Pi_2=0$ if all compactly supported common
trace variations are available. With a restricted trace space it only says
that the sum annihilates that space. At the flat cut $x=0$, with region 1 at
$x<0$ and region 2 at $x>0$, the sum per unit measure is
$-\partial_x\phi_1+\partial_x\phi_2$: opposite outward normals give
continuity of the derivative in a \emph{single} signed coordinate.

\section{Which traces and which corner representative?}
The first variation, not a count of field components, decides how many
response densities are present. Consider a term proportional to
$\tfrac12(\partial_x^2\phi)^2$ in a flat collar, suppressing spectator
coordinates. At fixed tangential coordinates, two integrations by parts give
\[
 \int_a^b \phi''\,\delta\phi''\,dx
 =\int_a^b \phi''''\,\delta\phi\,dx
    +[\phi''\,\delta\phi'-\phi'''\,\delta\phi]_a^b.
\]
Now $q=(\phi,\partial_x\phi)$, and both coefficients must be paired in the
common signed $x$ coordinate. A smooth bulk field restricted to the two
regions matches \emph{all} normal derivatives in that coordinate. Conversely,
conditions from the displayed first variation alone guarantee neither
infinite-order matching nor the regularity of a weak solution. Such a
conclusion needs the equations and a regularity argument.

There is also a genuine freedom of presentation \cite{HW}. For a local
spacetime form $Y$ linear in the variation, $\Theta\mapsto\Theta+dY$ and
$C_F\mapsto C_F+Y|_F$ leave \eqref{eq:04-endpoint} unchanged by Stokes'
theorem. Moving a total derivative $d\mu$ into $L$ and subtracting $\mu|_F$
from each side action shifts $\Theta$ by $\delta\mu$; it changes
$\theta_\Sigma$ only by $\delta\int_\Sigma\mu$. Its field-space exterior
derivative is unchanged. An unmatched wall interaction is not such an
ambiguity: it changes what the common trace must do. The following test
detects the difference.

\begin{exercise}
On the scalar cut at $x=0$ leave the single surface action
$-\tfrac12\lambda\int j(t)^2\,dt$ in the joined system. Does stationarity
still give continuity of the normal derivative in the common coordinate?
\end{exercise}
\noindent\emph{Answer.} Its variation is $-\lambda\int j\,\delta j\,dt$.
Adding the two outward responses gives
$\partial_x\phi_2-\partial_x\phi_1=\lambda j$. The result describes a
physical defect in the string, not an invisible cut.
\chapter{The covariant phase space and its boundary charges}\label{ch:05}
The endpoint term in an action variation does more than supply canonical
momenta. Its change under a second variation tells us whether a proposed
gauge transformation is a redundancy of a region or changes a measurable
flux at its boundary. No canonical coordinates on the full space of fields
are needed to make the test.

\section{From endpoints to a two-form}
At prescribed cut datum $j$, the \emph{pre-phase space}
$\widetilde{\mathcal P}_R[j]$ is the space of solutions satisfying the physical
and prescribed cut conditions, before any identification by gauge. At a
regular solution its tangent vectors solve the linearized equations and the
linearized boundary conditions. A spacetime integral can still be a form on
this space: for instance
\[
 \theta_\Sigma(Y)=\int_\Sigma\Theta(\phi;Y)
      -\sum_F\int_{\partial_F\Sigma}C_F(\phi;Y)
\]
is a number linear in the variation $Y$. The sum includes each side face of
Chapter~\ref{ch:04}. The exterior derivative in \emph{field} space is
\begin{align}
 \omega_\Sigma&=\delta\theta_\Sigma,\nonumber\\
 \omega_\Sigma(Y,Z)&=Y[\theta_\Sigma(Z)]-Z[\theta_\Sigma(Y)]
                                      -\theta_\Sigma([Y,Z]).
 \label{eq:05-omega}
\end{align}
In particular, $\theta=p\,\delta q$ gives $\omega=\delta p\wedge\delta q$,
not its negative. The commutator term allows $Y$ and $Z$ to vary from one
solution to the next. The integrals of $\delta C_F$ enter $\omega_\Sigma$
with exactly the same minus signs as the $C_F$ enter $\theta_\Sigma$.

Pull the first-variation identity back to the space of solutions. It becomes
$\delta S_R=\theta_{\Sigma_f}-\theta_{\Sigma_i}$. Taking one more field-space derivative gives
$0=\delta^2S_R=\omega_{\Sigma_f}-\omega_{\Sigma_i}$. This proof of slice
independence would have failed had we dropped a physical corner momentum.
For example, return to the string with a massive endpoint $B$ in
Chapter~\ref{ch:04}, holding the other endpoint fixed. For two perturbations
$u,v$, the bulk two-form is $\int(\dot u v-\dot v u)\,dx$. Differentiate it
in time and use $\ddot u=\partial_x^2u-m^2u$ and the same equation for $v$.
The mass terms cancel, and spatial integration by parts leaves
$(\partial_n u)v-(\partial_n v)u$ at $B$. The endpoint equation converts
this into $-d[\kappa(\dot u v-\dot v u)]/dt$. Thus the bulk form alone
changes by exactly the amount stored at the endpoint.

The same endpoint argument tells us what changes when we open the cut.
There the first variation
has an additional one-form $\langle\Pi,\delta q\rangle_\Gamma$, so
\[
 (\omega_{\Sigma_f}-\omega_{\Sigma_i})(Y,Z)
 =-\langle Y\Pi,Zq\rangle_\Gamma+\langle Z\Pi,Yq\rangle_\Gamma.
\]
Here the integrated potentials are those of the opened action. Flux through
one region need not vanish. On matched solutions the two fluxes cancel,
since both variations preserve the common trace and the vanishing sum of
outward responses. Closedness, $\delta\omega_\Sigma=0$, follows from its
definition even before one asks about slice independence. A closed two-form
with a possible kernel is called \emph{presymplectic}. Only after identifying
its null directions and, when regular, taking the appropriate quotient does
one obtain a symplectic physical phase space.

\section{An action symmetry and the sign of its generator}
Let $X$ be an infinitesimal transformation tangent to the specified domain,
with any external symmetry parameter held fixed when fields vary. Test it as
an action symmetry \emph{off shell}:
\[
 \iota_X\delta S_R=\alpha_{X,\Sigma_f}-\alpha_{X,\Sigma_i}+\beta_X,
 \qquad \delta\beta_X=0.
\]
The endpoint functional $\alpha_X$ includes contributions from a changed
boundary Lagrangian or corner. It need not be zero merely because the field
equations are invariant. Comparing this identity with the on-shell variation
gives the Noether generator \cite{HW}
\begin{equation}
 H_{X,\Sigma}=\iota_X\theta_\Sigma-\alpha_{X,\Sigma}.
 \label{eq:05-noether}
\end{equation}
It has equal endpoint values if $\beta_X=0$; a nonzero field-independent
$\beta_X$ shifts their difference by a prescribed number. To obtain a
Hamiltonian vector field, the transformation must also have a differentiable
charge with the chosen boundary data. For the representative of the potential
being used, the endpoint check is
$\mathcal L_X\theta_\Sigma=\delta\alpha_{X,\Sigma}$ on solution variations
(including the corner). Cartan's identity then fixes the slot and the sign:
\[
 \delta H_{X,\Sigma}=
 \mathcal L_X\theta_\Sigma-\iota_X\delta\theta_\Sigma-\delta\alpha_{X,\Sigma}
 =-\iota_X\omega_\Sigma,\qquad
 \delta H_{X,\Sigma}(Y)=\omega_\Sigma(Y,X).
\]
If the endpoint check leaves a nonzero one-form, the stated charge does not
generate $X$ with this potential and these boundary data. The test is local
to each slice; an identity for the difference of two endpoints is not by
itself this test.

For example, a free particle with $L=m\dot q^2/2$ has the boost $Xq=vt$.
Off shell $XL=mv\dot q=d(mvq)/dt$, so $\alpha_X=mvq$. Since
$Xp=mv$ and $\mathcal L_X(p\,\delta q)=mv\,\delta q$ at fixed $t$,
\eqref{eq:05-noether} gives $H_X=v(tp-mq)$. Directly,
$\omega(Y,X)=v(t\,Yp-m\,Yq)$. The term $-mvq$ is what generates the change
of momentum; $vtp$ alone would only generate the change of position.

\section{Gauss' law and a charged cut}
Take source-free Maxwell theory on an ultrastatic slab,
$g=-dt^2+h_{ij}dx^i dx^j$, with the spatial metric fixed. Vary
$L=-F_{\mu\nu}F^{\mu\nu}\epsilon/4$, $F=dA$. Antisymmetry gives
\[
 -F^{\mu\nu}\nabla_\mu\delta A_\nu
 =(\nabla_\mu F^{\mu\nu})\delta A_\nu
     -\nabla_\mu(F^{\mu\nu}\delta A_\nu).
\]
Use walls with no corner kinetic term here, for instance fixed tangential
$A$ on a physical wall; take gauge parameters to vanish there. Define
$E^i=F^{i0}=h^{ij}F_{0j}$, so that $-F^{0i}\delta A_i=E^i\delta A_i$ is
the \emph{final}-slice coefficient. With $dV_h$ the spatial volume measure,
\[
 \theta_\Sigma=\int_\Sigma E^i\delta A_i\,dV_h,
 \qquad \omega_\Sigma=\int_\Sigma\delta E^i\wedge\delta A_i\,dV_h,
 \qquad D_iE^i=0.
\]
The last equation is the time component of Maxwell's equations; $D_i$ is
the connection of $h$. The gauge transformation $X_\lambda A=d\lambda$
leaves $F$, and hence $\theta_\Sigma$, invariant for fixed $\lambda$;
$\alpha_\lambda=0$. Integrating its generator by parts gives
\begin{align*}
 H_{\lambda,\Sigma}
 &=\int_\Sigma E^i D_i\lambda\,dV_h
 =-\int_\Sigma\lambda D_iE^i\,dV_h
    +\int_{\partial\Sigma}\lambda E_\perp\,dS,\\
 \omega_\Sigma(Y,X_\lambda)
 &=\int_{\partial\Sigma}\lambda\,Y E_\perp\,dS
 \quad\text{on the linearized Gauss constraint}.
\end{align*}
Here $E_\perp=r_iE^i$, with $r$ the \emph{outward} unit normal to
$\partial\Sigma$ within $\Sigma$. For a transformation preserving the boundary
problem, vanishing trace of $\lambda$ on every boundary component is
sufficient for a proper gauge variation: it lies in the
kernel and carries no boundary charge. At a cut $S=\Gamma\cap\Sigma$ an
allowed $\lambda$ with nonzero trace can instead generate the charge
$Q_\lambda=\int_S\lambda E_\perp\,dS$ if that flux can vary. Not every
nonzero parameter gives a nontrivial generator. For example, a spatial
constant has $d_\Sigma\lambda=0$ and sees only the total flux through the
whole closed boundary, which vanishes by Gauss' law. Also, a transformation
that changes prescribed data is not tangent to that fixed-$j$ problem: it
may act between fibers in the opened family. A slice charge there need not
be conserved in the presence of flux.

For a source-free join the outward electric fluxes on the two copies of $S$
are opposite. Equal gauge parameters on both sides have zero \emph{sum} of
cut charges; the regional charge was not an error to be discarded in advance.
Nor does a surface integral by itself put independent fields at the cut. To
see what the latter choice entails, introduce a boundary coordinate $u$ and
an independent momentum density $p$, adding $\int_S p\,\delta u$ to
$\theta_\Sigma$. If $X_\lambda u=\lambda|_S$ and $X_\lambda p=0$, the new
cut generator on Gauss data is $\int_S\lambda(E_\perp dS+p)$. Imposing
$p=-E_\perp dS$ makes this diagonal transformation null on the constraint
surface. We added variables, a symmetry action, and a constraint. None of
these follows just by integrating the original bulk Gauss law by parts.

\begin{exercise}
In one spatial dimension take $h=dx^2$ on $[0,L]$ and set $A_t=0$ at the
endpoints. Set $a=\int_0^L A_x\,dx$. What charge generates time-independent
endpoint gauge parameters, and what happens if the endpoint values agree?
\end{exercise}
\noindent\emph{Answer.} Gauss' law sets $E^x=e$ independent of $x$, giving
$\theta=e\,\delta a$. The transformation shifts $a$ by
$\lambda(L)-\lambda(0)$, so $H=e[\lambda(L)-\lambda(0)]$. Equal values give
zero even if neither value vanishes. The minus sign at the left endpoint is
its outward orientation, not a sign assigned to its gauge parameter.
\chapter{Causal response and the Peierls bracket}\label{ch:06}
How can a force applied for a finite time define a tangent vector to the
\emph{unforced} equations? Produce two responses to the same force: one with
no incoming disturbance, the other with no outgoing disturbance. Their
difference solves the homogeneous equations. The endpoint two-form of
Chapter~\ref{ch:05} tells us which difference is the Hamiltonian vector of
the source observable. This construction takes place near a specified
solution, and its domain matters as much as the differential operator.

\section{Fields, sources, and a causal boundary problem}
Fix a smooth solution $\bar\phi$. A perturbation $u$ is a section of the
linearized field bundle $V$. A smooth bulk source $f$ is a section of the
different bundle $V^\vee=V^*\otimes |\Lambda^{\mathrm{top}}T^*M|$; its value
on $u$ is a density, not another field. The pairing and the linearized Euler
operator are
\[
 \langle f,u\rangle=\int_M f(u),\qquad
 J=\left.D{\cal E}\right|_{\bar\phi}:
           \Gamma(V)\longrightarrow\Gamma(V^\vee).
\]
Here $\Gamma(V)$ means smooth sections of $V$, and the oriented Euler form
is read as a density using the spacetime orientation. $J$ is the Hessian
with the sign of the Euler term in
$\delta S$, not an arbitrarily signed wave operator. At this background the
derivative of an observable $F$ is a source when $\delta F(u)=\langle F',u\rangle$.
A choice of metric and measure may identify field components with source
components; it does not exchange their roles in this formula.

For a first source space take smooth dual densities compactly supported in
the spacetime interior. For a nongauge, normally hyperbolic equation on a
globally hyperbolic spacetime without a wall, the causal Green problem for
such tests is well posed \cite{BGP}; reversing the overall sign of the
operator reverses that of its inverses. With a timelike physical wall one
must supply a well-posed initial--boundary problem instead. Use the
homogeneous \emph{linearized} physical boundary condition for both responses
when that condition is its own adjoint. At a cut still held at fixed $j$ the
corresponding trace variation also vanishes. Vanishing past data for a
retarded solution do not mean imposing zero Cauchy data at both ends of a
slab. The advanced problem has the complementary future data.

There are two distinct adjoint assertions here. For interior test sections,
commuting two variations of the action gives formal symmetry $J^*=J$ by
integration by parts. In a slab containing the physical wall the same
calculation, before discarding endpoints, gives
\begin{equation}
 \langle Jv,u\rangle-\langle Ju,v\rangle
 =\omega_{\Sigma_f}(u,v)-\omega_{\Sigma_i}(u,v)
 \label{eq:06-green}
\end{equation}
for $u,v$ obeying the linearized wall condition, with the corner terms in
$\omega$. Indeed the antisymmetrized second bulk variation is
$\langle Ju,v\rangle-\langle Jv,u\rangle$; the complementary terms are the
two endpoint variations. This identity is for trial perturbations, not just
for solutions of $Ju=Jv=0$. Dropping the wall integral for arbitrary fields
would be false. For a scalar, real Robin data $\partial_n u+ku=0$ with
the same fixed $k$ for both variations give
$u\partial_n v-v\partial_n u=0$ on the wall. This is the extra domain check.
An absorbing wall may require a different adjoint boundary condition;
formal symmetry alone does not prove reciprocity for the chosen propagators.

\section{Inverting the endpoint form}
Suppose the adjoint-domain check, existence, uniqueness, and causal support
hold. Write $G^R f$ and $G^A f$ for the retarded and advanced solutions of
$Jv=f$, supported in the causal future and past of $\operatorname{supp}f$,
respectively. On interior source tests $JG^{R/A}=1$. To use an inverse on
the other side, note that $w-G^R Jw$ is the homogeneous solution with the
initial data of $w$. Thus on a finite slab $G^R Jw=w$ requires zero incoming
Cauchy data, and $G^A Jw=w$ requires zero outgoing Cauchy data. Admissible
test fields supported away from both caps meet both conditions; compactness
on a \emph{closed} slab alone does not. Using \eqref{eq:06-green} with a retarded
solution for $f$ and an advanced one for $g$ gives
$\langle g,G^R f\rangle=\langle f,G^A g\rangle$: before the sources the
retarded endpoint vanishes; after them the advanced endpoint vanishes.
Define
\[
 E=G^R-G^A.
\]
Thus $JEf=0$, $EJw=0$ on the common test-field domain just described, and
$E$ is skew with respect to the field--source pairing. For a compact source support
$K$, the support of $Ef$ is in $J^+(K)\cup J^-(K)$, which meets every Cauchy
surface in a compact set in the globally hyperbolic case without a wall.
One-sided inverses can have larger domains: a source turned on after a fixed
slice and left on forever can generate a retarded solution without admitting
a response with zero data after the source is over. One must specify a common
source space before subtracting the two inverses.

For $Ju=0$, take $v=G^R f$ in \eqref{eq:06-green} and put the final slice
after the source. There $G^A f$ vanishes, and at the initial slice $G^R f$
vanishes. Hence $\omega_{\Sigma_f}(u,Ef)=\langle f,u\rangle$. Since $Ef$
is homogeneous, the equality holds on every slice on which the pairing is
finite. Comparing with $\delta H(u)=\omega(u,X_H)$ gives
\begin{equation}
 X_F=E F',\qquad \{F,G\}_{\mathrm P}=\delta F(X_G)
                 =\langle F',E G'\rangle.
 \label{eq:06-peierls}
\end{equation}
With a source added as $S-\varepsilon G$, the linearized equation reads
$Jv=G'$, so its retarded response minus its advanced response is precisely
$X_G$. With $S+\varepsilon G$ each response would carry a minus sign.
Covectors differing by $Jh$ for $h$ in the common test-field domain give
the same tangent pairing and causal vector. Larger classes of sources need their own tests:
the Green operators must act on them, the dual pairings must exist, and
boundary terms must still satisfy \eqref{eq:06-green}. A force applied to an
opened trace pairs with that trace and drives an inhomogeneous boundary
problem; it is not automatically an interior delta source in the fixed-$j$
problem. Likewise singular functional derivatives may require wavefront
conditions before the pairing of distributions in \eqref{eq:06-peierls}
makes sense.

An oscillator determines all of the signs without appealing to a commutator.
For $S=\tfrac m2\int(\dot q^2-\nu^2q^2)dt$, $m,\nu>0$, one has
$J=-m(\partial_t^2+\nu^2)$. With $s=t-t'$ its causal inverses are
\[
 G^R(s)=-\vartheta(s)\frac{\sin(\nu s)}{m\nu},\qquad
 G^A(s)= \vartheta(-s)\frac{\sin(\nu s)}{m\nu}.
\]
Here $\vartheta$ is the step function. Both kernels are continuous and their
first derivatives jump by $-1/m$ at $s=0$, as required by $JG=\delta(s)$.
It follows that $\{q(t),q(t')\}=E(s)=-\sin(\nu s)/(m\nu)$. Differentiating
in the \emph{second} time gives $\{q(t),p(t')\}=m\partial_{t'}E(t-t')
=\cos(\nu(t-t'))$, hence $\{q,p\}=1$ at equal time. Evaluation at a time
can equivalently be obtained from smooth time smearings in this example.

\section{When the Hessian has gauge directions}
An infinitesimal proper gauge transformation $K\varepsilon$ is tangent to the
allowed fields and lies in the presymplectic kernel; a charged cut action
need not be proper. Because gauge transformations take solutions to solutions,
differentiating their action gives $JK=0$ at a solution. Hence $J$ cannot
have an inverse on all fields. The Batalin--Vilkovisky (BV) formalism
\cite{CMR} introduces ghosts for gauge parameters and dual variables
(antifields) for fields and ghosts. A differential $Q$ on these graded
variables encodes the equations, symmetries and their identities in
$Q^2=0$. At its zero $\bar\Phi$ (the classical solution $\bar\phi$ with
ghosts and antifields set to zero), $D=(DQ)_{\bar\Phi}$ squares to zero.
For an irreducible theory, with no further local gauge redundancies among
parameters in $U$, the relevant degrees of the deformation complex are
\[
 \underbrace{\Gamma(U)}_{-1}\xrightarrow{\ K\ }
 \underbrace{\Gamma(V)}_{0}\xrightarrow{\ J\ }
 \underbrace{\Gamma(V^\vee)}_{1}\xrightarrow{\ K^*\ }
 \underbrace{\Gamma(U^\vee)}_{2}.
\]
Here and below the spaces carry compatible support and boundary conditions.
The compositions vanish: the first says that a gauge variation of a solution
solves the linearized equations, the second is the linearized Noether
identity. Boundary tangency is a further requirement. If boundary conditions
are written $b(\Phi)=0$, then $Q$ acts within this problem only if
$b'_\Phi Q(\Phi)=0$ there. At $Q(\bar\Phi)=0$, differentiating in a tangent
direction $w$ gives $b'_{\bar\Phi}Dw=0$. This is why one may restrict the
complex to that boundary domain, rather than testing only its bulk equations.
For Maxwell, $K\varepsilon=d\varepsilon$; fixing the tangential potential on $B$ requires
$d_B(\varepsilon|_B)=0$, which the proper choice $\varepsilon|_B=0$ satisfies.
On a sewn interface matching the two parameter traces similarly preserves
matching of tangential potentials. For a full complex, each arrow must
preserve its specified boundary or transmission domain, not just the field
block in this example.

A \emph{retarded or advanced Green homotopy} $h^{R/A}$ lowers degree by one
and satisfies $Dh^{R/A}+h^{R/A}D=1$ on its corresponding causal support
complex. It is an inverse up to the neighboring gauge and identity arrows.
To construct one, let $Cu=0$ be a linear gauge condition, with $C$ mapping
field perturbations to parameter sections (for Maxwell one can take the
divergence). Choose a symmetric invertible bundle map $B:U\to U^\vee$;
this algebraic weight pairs gauge-condition violations. The gauge-fixed
field operator, the operator giving the change of the gauge condition along
a gauge orbit, and its formal adjoint are respectively
\[
 P=J+C^*BC,\qquad R=CK,\qquad R^*=K^*C^*.
\]
For either causal sign, suppress its superscript and assume $P,R,R^*$
have Green inverses $G_P,G_R,G_{R^*}$ on the corresponding one-sided
domains. The compositions below must preserve support, the field and proper
ghost boundary conditions, and any transmission and homogeneous cap data.
These are hypotheses on the gauge fixing, not consequences of adding a
quadratic term to the action \cite{Khavkine}.

Since $JK=0$, one has $PK=C^*BR$. Therefore $KG_R$ and $G_PC^*B$ solve
the same $P$ problem with the same causal data and are equal by uniqueness.
Likewise $K^*P=R^*BC$ gives $BCG_P=G_{R^*}K^*$. Use the following blocks
on inputs of degrees $0,1,2$, respectively:
\[
 h_0=G_R C,\qquad h_1=G_P,\qquad h_2=C^*G_{R^*}.
\]
For example, on a field perturbation the degree-zero homotopy equation is
$Kh_0+h_1J=KG_R C+G_PJ=G_P(C^*BC+J)=1$. On a source the second
intertwining relation gives $Jh_1+h_2K^*=(J+C^*BC)G_P=1$; the end blocks
are $h_0K=G_RR=1$ and $K^*h_2=R^*G_{R^*}=1$ on their respective domains.

For a concrete check, identify forms with their density duals using the metric on a boundaryless globally hyperbolic spacetime. Denote the formal adjoint of the spacetime exterior derivative by $d^*$, not by the field variation $\delta$. Maxwell gives $K=d$, $J=-d^*d$ on one-forms. Take $C=d^*$ and $B=-1$. With $W_k=dd^*+d^*d$ on $k$-forms, the field block is $P=-W_1$ and the scalar ghost block is $R=W_0$. If $S_k$ is the same-sign causal inverse of $W_k$, the three homotopy blocks are $S_0 d^*$, $-S_1$, and $d S_0$. For instance $d S_0d^*+S_1d^*d=S_1(dd^*+d^*d)=1$: the equality $d S_0=S_1 d$ follows from the differential-operator identity $W_1 d=d W_0$ and causal uniqueness. Adding a wall would oblige us to verify that $d$ really maps the allowed scalar boundary domain into the one-form domain. Commuting the formal differential expressions would not perform that test.

Restoring the retarded and advanced signs and subtracting their identities gives
$D(h^R-h^A)+(h^R-h^A)D=0$. For a closed source $f$ of degree one,
$K^*f=0$, the identity says $Jh^{R/A}f=f$. Thus
$(h^R-h^A)f$ is a homogeneous perturbation. Why should two such gauge choices
give the same bracket? For either causal sign set $w=(h-\widehat h)f$.
Then $Dw=0$. If both homotopies act on the same one-sided support complexes,
with proper parameter domains, their identity on $w$ gives
$w=D\widehat h w=K(\widehat h w)$. This is an allowed proper gauge change.
Pairing it with another closed source $g$ gives zero:
$\langle g,K\varepsilon\rangle=\langle K^*g,\varepsilon\rangle=0$, including
the boundary-domain check. For instance, in Maxwell theory this is the
familiar current conservation test: smearing $d\varepsilon$ against a
compact interior current $j^\mu$ gives
$-\int_M(\nabla_\mu j^\mu)\varepsilon\,|\epsilon|$ after integration by parts.
Therefore the field-block pairing of the causal difference is gauge
independent for the stated class. Skewness uses the adjoint relation of the
two causal field blocks on these sources, as in \eqref{eq:06-green}. A source
with an uncancelled surface term at the cut is not covered by this proof.
If the homotopy identity has right-hand side $1-P_{\mathrm{res}}$ for a
projection onto residual modes, those modes must be retained or treated
separately before making the same inference.

Finally, \eqref{eq:06-peierls} at one background gives a fixed bilinear
bracket for linear observables on the tangent theory. In a nonlinear theory
$J_\phi$, the Green operators, and possibly their domains change with the
solution. Antisymmetry at each background is not the nonlinear Jacobi
identity. One route to that identity requires a smooth (locally reduced)
phase space, Hamiltonian vectors for the admissible observables, and closure
of their algebra: the Jacobi identity then follows from $\delta\omega=0$.
The existence of one causal inverse at $\bar\phi$ establishes none of those
smooth-family assertions.

\begin{exercise}
On coordinates $(x,y,z)$ define a skew bracket by the Leibniz rule and
$\{x,y\}=1$, $\{y,z\}=y$, $\{z,x\}=0$. Does freezing the coefficients at
each point test Jacobi for the original bracket?
\end{exercise}
\noindent\emph{Answer.} Each frozen matrix gives a constant Poisson bracket.
For the varying bracket the cyclic sum on $(x,y,z)$ is
$\{x,\{y,z\}\}+\{y,\{z,x\}\}+\{z,\{x,y\}\}=1$. The missing operation in
the frozen test is the derivative of the bracket itself.

\part{Reconstructing the classical theory}
\chapter{Transmission and the reconstruction of solutions}\label{ch:07}
The variational principle exchanges only a few functions across a face. A smooth field has infinitely many derivatives there. Why should a few functions suffice? They do not suffice off shell. On shell the field equation may determine the missing derivatives, but even on a noncharacteristic face it matters which branch of a nonlinear theory the two sides occupy. We will prove a reconstruction result in which neither analyticity nor a previously known global solution is an input.

\section{The variable shared by two variations}
Work near the interior of an artificial face in a fixed geometric assembly. Write $s<0$ on $R_-$, $s>0$ on $R_+$, and use $y=(y^1,\ldots,y^{n-1})$ along the face. Use compatible frames; thus the face identification $D$ is temporarily invisible. Consider $m$ real components $u^A$ with a first-order action density $\mathcal L(s,y,u,v)\,|ds\,d^{n-1}y|$, where $v^A_\mu=\partial_\mu u^A$. Calculations are local, so a density can be represented by its coefficient in this coordinate system. Write
\begin{equation}
 p_A^\mu=\frac{\partial\mathcal L}{\partial v^A_\mu},\qquad E_A=\frac{\partial\mathcal L}{\partial u^A}-\partial_\mu p_A^\mu .\label{eq:normal-momentum}
\end{equation}
The primitive choice in this problem is to give the two fields a \emph{common value} $q(y)$. Vary a pair of smooth half-fields with this value but without yet relating their normal derivatives. A compactly supported variation of $q$ is realized by extending it into each collar with a cutoff equal to one at the face. Integration by parts on the two sides then gives
\begin{equation}
 \delta(S_-^o+S_+^o)\big|_\Gamma
   =\int_\Gamma (p_{-,A}^s-p_{+,A}^s)\,\delta q^A\,d^{n-1}y .\label{eq:summed-momentum}
\end{equation}
Thus common value is an identification, while common \emph{cooriented momentum} $p^s$ is an equation. Equivalently, \emph{outward} responses add to zero: $\Pi_-+D^\vee\Pi_+=0$, where the dual transport is fixed by the integral pairing with $\delta q$. These signs also apply to a spacelike face; one need not guess the sign from a drawing of two unit normals. Restricted variations, for instance those intrinsic to a constrained trace manifold, test only the corresponding tangent directions. In that case it is the pullback of the boundary one-form that must vanish. Equation~\eqref{eq:summed-momentum} with arbitrary $\delta q^A$ is not a universal set of boundary variables for every first-order constrained or higher-derivative theory.

\section{A prescribed load is not a new degree of freedom}
Let $[b]=b_+|_0-b_-|_0$. A continuous piecewise smooth field has weak first derivatives equal to its piecewise derivatives, with no delta function. Its Euler density does, however, have the singular part $-[p^s]\delta(s)$: this follows by differentiating the two step functions multiplying $p^s$. It is the distributional version of \eqref{eq:summed-momentum}. Adding a sheet action changes what it means to solve the global problem. For instance, with a prescribed surface density $j_A$ the term $-\int_\Gamma j_A q^A$ gives
\[
 p_{-,A}^s-p_{+,A}^s=j_A .
\]
The piecewise bulk field solves the equation $E_A=j_A\delta(s)$, not the smooth sourceless equation. For a scalar in flat spacetime with $s=x$ a spatial coordinate and $\mathcal L=-\tfrac12\partial_\mu u\,\partial^\mu u-V(u)$, $p^x=-\partial_x u$. The result is $[\partial_x u]=j$. This is why forcing all transverse jets to agree would discard the physical solution in the presence of a load.

For comparison, place a dynamical coordinate $\chi(t)$ on a timelike seam in two spacetime dimensions. In a fixed time coordinate let its action be
\[
 S_\Gamma=\int dt\left[\frac{\mu}{2}\dot\chi^2-\frac{\mu\Omega^2}{2}\chi^2-\frac{\kappa}{2}(\chi-q)^2-jq\right].
\]
Varying both $q$ and $\chi$ gives, per unit $t$,
\begin{align*}
 \Pi_-+\Pi_++\kappa(\chi-q)-j&=0,\\
 \mu(\ddot\chi+\Omega^2\chi)+\kappa(\chi-q)&=0 .
\end{align*}
There is now initial data for $\chi$, and the endpoint potential contains $\mu\dot\chi\,\delta\chi$. A clamped source $j$ alone has neither property. This is a physical sheet, not a choice of regulator for a transparent seam. On a higher-dimensional sheet a tangential integration by parts has an additional destination, its own boundary: the edge variation must be canceled, set to zero by allowed boundary data, or accounted for by an edge action. Appendix~\ref{app:A} fixes the incidence signs.

\paragraph{Where does a sheet current go at its edge?} Gauge invariance puts a restriction even on sources that are prescribed rather than varied. In $2+1$ dimensions take a timelike strip with coordinates $(t,y)$, $0\leq y\leq b$. Couple its current density $k^a$ to the pullback of a Maxwell potential, and an edge charge $Q_e(t)$ to its time component at $y=b$:
\[
 S_{\rm src}=\int dt\,dy\,(k^t A_t+k^y A_y)+\int dt\, Q_e A_t(t,b).
\]
Under $A_a\mapsto A_a+\partial_a\varepsilon$ choose a test parameter supported away from the time caps and $y=0$. Integration by parts gives
\[
 \delta_\varepsilon S_{\rm src}=-\int dt\,dy\,\varepsilon(\partial_t k^t+\partial_y k^y)
  +\int_{y=b} dt\,\varepsilon\,(k^y-\dot Q_e).
\]
The current must be conserved inside the strip, and current leaving the strip at its edge must increase the edge charge: $\dot Q_e=k^y(t,b)$. If the theory has no such charge and no neighboring face to receive current, $k^y(t,b)$ must vanish. These are consistency conditions for the source. They neither create an independent edge oscillator nor follow merely from counting faces. Allowing charge on an edge and choosing dynamics for that charge are separate decisions.

\section{How a finite trace recovers all normal derivatives}
Here is a useful positive answer for nonlinear variational equations. State it first in one collar; initial and physical outer boundary conditions, if present, are unchanged. The word smooth on a closed half-collar includes smooth one-sided derivatives up to its face. The result concerns such solutions, not the creation of smoothness from arbitrary finite-energy traces.

\begin{theorem}[Transparent sewing on one regular branch]\label{thm:branch}
Let $\mathcal L$ be the same smooth first-order density on both sides in the assembled coordinates and frames. Let $u_-,u_+$ be smooth up to their closed half-collars and solve $E_A=0$ on the interiors. At fixed $(s,y,u,v_i)$, $i$ tangential, suppose the normal Legendre map
\[
 w^B=v_s^B\longmapsto p_A^s(s,y,u,v_i,w)
\]
has a single smooth local inverse covering the two realized normal velocities at each face point. Equivalently for this purpose, equality of their momenta separates these two velocities and the normal Hessian $A^{ss}_{AB}=\partial p_A^s/\partial v_s^B$ is invertible at the matched jet. Then equality of $q$ and $p^s$ is necessary and sufficient for the pair to be restrictions of a unique smooth solution on the union. A given smooth, globally matching bulk source can be included on both sides.
\end{theorem}
\begin{proof}
The common value is a smooth function of $y$, so its tangential derivatives agree. Both sides therefore use the same normal Legendre map at the face. On the stated branch the common momentum gives the same $\partial_s u$. Tangentially differentiating this equality gives all tangential derivatives of the first normal jet, including $\partial_i\partial_s u$.

To see the next step rather than assume it, set $A^{\mu\nu}_{AB}=\partial p_A^\mu/\partial v_\nu^B$. Expand $\partial_\mu p_A^\mu$ by the chain rule. On a solution,
\begin{align}
 A^{ss}_{AB}\,\partial_s^2 u^B
   &=\mathcal L_{u^A}-(\partial_{x^\mu}p_A^\mu)_{u,v}
       -\frac{\partial p_A^\mu}{\partial u^B}\,\partial_\mu u^B
       -\sum_{(\mu,\nu)\ne(s,s)}A^{\mu\nu}_{AB}\,\partial_\mu\partial_\nu u^B .\label{eq:solve-normal}
\end{align}
Here $x^\mu=(s,y)$, and the derivative with subscript $u,v$ acts on explicit coordinate dependence only. The right side is known in terms of the jets already matched and their tangential derivatives. The inverse of $A^{ss}$ thus gives a common second normal jet. More generally, write \eqref{eq:solve-normal} with $\partial_s^2u$ isolated and differentiate $k$ times in $s$. The right side uses normal order at most $k+1$, with tangential derivatives, since no remaining second derivative has two normal indices. Induction yields normal order $k+2$. A common smooth source contributes its known Taylor coefficients to the same recursion.

The closed-half smoothness criterion of Chapter~\ref{ch:03} now gives one smooth field. It solves the equation at the face by continuity. Piecewise evaluation makes it unique. Conversely the restrictions of a smooth solution have the same value and momentum in a common frame. No convergent Taylor series was used: each equality is a trace of derivatives of actual half-fields.
\end{proof}

Here the principal symbol in the conormal direction is, up to the overall equation sign, $A^{ss}$. Its invertibility is noncharacteristicity for this second-order system. It is a pointwise algebraic statement, not a claim about stable evolution with $s$ as time. For example, a timelike surface can be noncharacteristic for a wave equation without being a Cauchy surface. Nor does the theorem assert that every prescribed trace occurs on a solution.

\paragraph{Why a branch matters.} Take $\mathcal L=\tfrac12 u_t^2-W(u_x)$, $W(v)=\tfrac14(v^2-1)^2$, in rescaled units. On the two sides of $x=0$ set $u=|x|$. Each side is a smooth static solution. At the cut $q=0$, $p^x=-W'(u_x)=0$ on both sides, and $-W''(\pm1)=-2$ is nonsingular at both realized velocities. Yet $u_x$ jumps from $-1$ to $+1$. It is a weak stationary configuration, not a smooth field. Knowing that the linearized equation is noncharacteristic \emph{on each side} did not tell us that the two velocities occupy the same branch. In particular a proof that inverts the normal symbol only after having silently equated $u_x$ would miss this example.

\paragraph{Which cone belongs to the field?} To apply the theorem to a branch with genuine derivative nonlinearity, take $\mathcal L=P(X)$ on a flat $(t,x)$ patch, $X=(u_t^2-u_x^2)/2$, $P(X)=X+\beta X^2$, with $\beta>0$ in the corresponding units. The constant-gradient solution $u=v t$ has $v\ne0$. Direct differentiation in the two velocities gives its normal-symbol matrix and characteristic speed:
\[
 (A^{\mu\nu})=\begin{pmatrix}1+3\beta v^2&0\\0&-(1+\beta v^2)\end{pmatrix},\qquad
 c_{\rm field}^2=\frac{1+\beta v^2}{1+3\beta v^2}<1 .
\]
Indeed a plane disturbance depending on $x-c t$ has principal coefficient $(1+3\beta v^2)c^2-(1+\beta v^2)$. For the metric-null face $s=t-x$, the contraction $A^{\mu\nu}\partial_\mu s\,\partial_\nu s=2\beta v^2$ is nonzero. In a sufficiently small neighborhood of that jet the scalar normal momentum has a local inverse; Theorem~\ref{thm:branch} applies to existing half-solutions on that branch even at this metric-null face. On $x-c_{\rm field}t=0$ it does not. As $\beta\to0$ the two cones coincide. The distinction between null and characteristic is therefore a statement about propagation, not a choice of name for a geometric cut.

\section{What changes when the field is a connection?}
The component $A_s$ of a gauge potential is a good way to spot an illegitimate use of Theorem~\ref{thm:branch}. Yang--Mills curvature $F_{si}=\partial_s A_i-\partial_i A_s+[A_s,A_i]$ contains no $\partial_s A_s$. The normal Legendre map for \emph{all} the components is singular. Demanding a gauge-invariant determination of $\partial_s A_s$ asks the equations to determine a convention.

A local result is nevertheless available. Take a flat product collar $g=\epsilon\,ds^2+h_{ij}dy^i dy^j$ with constant nondegenerate $h$ and $\epsilon=\pm1$. Use a boundary bundle identification to put tangential connections into the same frame; their common value and the equality of the cooriented $F^{si}$ are the transmission data for ordinary Yang--Mills theory with the same coupling. Assume smooth source-free half-solutions. On each side solve the ordinary differential equation $\partial_s g_0=-A_s g_0$, $g_0(0,y)=1$. The transformed connection $g_0^{-1}Ag_0+g_0^{-1}dg_0$ is in axial gauge $A_s=0$. This transformation fixes the frame on the face. Now $F^{si}=\epsilon h^{ij}\partial_s A_j$, so the transmission conditions match $A_i$ and $\partial_s A_i$. The tangential equations $D_\mu F^{\mu i}=0$ read
\[
 \partial_s^2 A_i=-\epsilon D^j F_{ji} .
\]
Tangential differentiation and the same normal induction give a smooth connection across the collar. The remaining equation $D^i\partial_s A_i=0$ is the constraint: it holds on both halves by hypothesis and on the joined face by continuity. It was not replaced by a new boundary condition. Any gauge transformation between two axial representatives satisfies $\partial_s g_0=0$, so two such arrows with the same face value also assemble. This proves a local result \emph{up to gauge}, including the arrows, not uniqueness of a preferred gauge potential on the entire manifold.

In more geometric language, restriction on solutions should be essentially surjective (a matching pair is isomorphic to a restriction) and fully faithful (its compatible changes of description are precisely restrictions of global ones). For ordinary fields of Theorem~\ref{thm:branch} this was a bijection. For the local product Yang--Mills problem the axial representatives exhibit the two stronger-looking conditions without discarding the arrows first. To assert an equivalence for a larger gauge system one must in addition justify its gauge slice or descent argument, the outer boundary conditions and any continuation across lower strata. At a corner, smoothness up to each sector and compatible jets in the incident sector chart are substantive assumptions; a singular corner solution is not covered by an interior-face proof.

\begin{exercise}[Identify the missing derivative]
In null coordinates $(u,v)$, take $\partial_u\partial_v w=0$ and cut at $u=0$. Does continuity of $w$ and its variational flux for the density $\partial_u w\,\partial_v w$ require $\partial_u w$ to agree? Construct a pair of smooth halves that answers the question.
\end{exercise}
\noindent\emph{Check.} The cooriented momentum is $\partial_v w$, already determined by the common trace. Take $w=0$ for $u<0$ and $w=u$ for $u>0$. Both the trace and the momentum vanish, and the piecewise solution even solves the equation weakly, but $\partial_u w$ jumps. The failure is now a characteristic one, not a choice between two locally invertible branches.
\chapter{Sewing a causal response}\label{ch:08}
A retarded solution answers a question with an arrow of time: where does an initially absent disturbance go after a source is applied? Knowing which derivatives of an existing solution match at a face does not give this answer. In particular, solving an algebraic mismatch equation at every frequency need not give a retarded inverse in time. We first construct an operator from independent experiments on the two regions; then we examine how the construction works differently at three kinds of face.

\section{The equation for the missing interface history}
Fix a smooth assembled background solution and linearize there. Let $P$ denote the resulting field operator, gauge fixed where necessary. On each region choose a well-posed auxiliary boundary problem at the cut, retaining the true outer boundary problem. For one choice of time direction, write
\begin{equation}
 v_a=G_a f_a+H_a\eta_a .\label{eq:particular-poisson}
\end{equation}
Here $G_a$ solves $P_a v_a=f_a$ with zero perturbation of the auxiliary trace; $H_a$ solves $P_a v_a=0$ with trace perturbation $\eta_a$. Both have zero incoming data for the retarded problem (zero outgoing data in the future for the advanced problem). $H_a$ here is a boundary-to-solution operator, not a Hamiltonian. The trace need not be Dirichlet, but \eqref{eq:particular-poisson} must parametrize all the allowed responses in the chosen class. This is an initial-boundary value hypothesis on a \emph{single} region, not information fetched from the global answer.

For a transparent first-order variational problem let $\mathcal B$ take a pair to the difference of the identified field traces and the sum of their outward linearized responses. Other problems use their actual transmission equations. Set $G=G_-\oplus G_+$ and $H=H_-\oplus H_+$. The two things we can measure independently are
\[
 Nf=\mathcal B Gf,\qquad M\eta=\mathcal B H\eta.
\]
The first is the error caused by freezing the cut; the second is how much an adjustment of the trace can change that error. The trace histories of a sourced sewn solution must therefore satisfy $M\eta=-Nf$. For a sheet that prescribes a linearized mismatch $b$, the right side is $b-Nf$ instead; a dynamical sheet generally adds unknowns and equations to $M$ itself.

\begin{proposition}[Causal response from regional data]\label{prop:causal}
Fix either the retarded or the advanced problem, with a time orientation and a propagation cone for the chosen hyperbolic system. The notation $J^\pm_{M_D}$ below refers to that cone on the assembled manifold; for a normally hyperbolic operator it is the metric cone. Take admissible global test sources restricted to the pieces, with their outer boundary and corner compatibility conditions. Suppose:
\begin{enumerate}
\item $G,H,\mathcal B$ act on specified solution and trace spaces as above, and zero mismatch assembles a solution of the declared global equation in its required regularity class, with no extra surface source.
\item $\operatorname{im}N\subseteq\operatorname{im}M$. On $\operatorname{im}N$ there is a linear right inverse $L$ of $M$ in those spaces, with trace assignments giving the required causal support and continuity. For retarded support this means no response outside $J^+_{M_D}(\operatorname{supp}f)$, including at the physical boundary; advanced support reverses the cone.
\item A homogeneous sewn solution in this causal class with zero past (respectively future) data is zero. For a physical quotient, replace zero by a proper gauge perturbation and the conclusion by equality of physical responses.
\end{enumerate}
Then, identifying matched pairs with assembled fields, the response is
\begin{equation}
 \mathcal G_\#=G-HLN .\label{eq:causal-sew}
\end{equation}
For the fixed-gauge problem this gives the causal Green operator on the declared domains; on compatible compact test fields both inverse identities hold.
\end{proposition}
\begin{proof}
Each component solves the correct interior equation because the correction is homogeneous there. Its mismatch is $Nf-MLNf=0$. The first two hypotheses give a genuine global solution with the right source and support. In particular this proves existence on the source domain; it does not assume a previously calculated global inverse. Uniqueness identifies it with any independently constructed causal response. For the second identity let $w$ obey the global wall and transmission conditions with $Pw$ in the source domain and with zero incoming Cauchy data at the past cap (zero outgoing data at the future cap for the advanced identity). Compact interior tests supported away from both caps have both properties. Then $\mathcal G_\#Pw-w$ is homogeneous with zero data at the required end, hence zero. Compactness merely up to a cap does not suffice: a nonzero homogeneous solution on a closed finite cylinder is a counterexample.
\end{proof}

A second choice of $L$ changes the answer by a field in $H\ker M$. A kernel can describe redundant trace coordinates ($H\eta=0$), a proper gauge transformation, or a genuine free excitation satisfying the matching equations. In the last case prescribing the causal data or the initial state of a sheet matters. Merely naming a right inverse says nothing about this distinction. Causal paths connecting two points of the face may travel through the bulk, so the support test is made in the assembled spacetime, not solely in the induced geometry on the face. In gauge theory one must apply the construction to the field, ghost and dual-ghost problems on domains that intertwine the gauge and gauge-fixing maps of Chapter~\ref{ch:06}; sewing the field block alone does not produce a Green homotopy.

\section{A spacelike cut carries data forward}
Suppose the cut is a Cauchy slice $t=t_*$ for a regular hyperbolic initial-value problem. For a retarded source before the cut, the earlier region supplies the field and its canonical momentum at $t_*$. Propagate this Cauchy pair forward with the homogeneous equation in the later region; add the response to subsequent sources with zero initial data. The future need not send a trace correction back to the past. For gauge fields the Cauchy pair must obey the constraint, and the evolution must propagate it. An arbitrary spacelike patch that is not Cauchy for the domain under study does not carry all incoming information.

One spatial Fourier mode makes this composition visible. For $D=\partial_t^2+\Omega^2$, its retarded impulse at time $t_0<t_*$ produces $(\sin(\Omega(t_*-t_0))/\Omega,\cos(\Omega(t_*-t_0)))$ at the cut. For $t>t_*$ its value is
\begin{align*}
 &\cos(\Omega(t-t_*))\frac{\sin(\Omega(t_*-t_0))}{\Omega}
 +\frac{\sin(\Omega(t-t_*))}{\Omega}\cos(\Omega(t_*-t_0))\\[-1mm]
 &\hspace{2.5em}=\frac{\sin(\Omega(t-t_0))}{\Omega}.
\end{align*}
The arbitrary time $t_*$ disappears. The zero-frequency formula follows by continuity, not by division by zero. On a curved spacetime the evolution need not be a trigonometric function, but uniqueness of Cauchy evolution gives the same composition law.

\section{A timelike cut exchanges histories}
Let $D=\partial_t^2-\partial_x^2$ on the line. Here we drive the equation as $Dw=f$; for the usual scalar action the Euler Hessian is $J=-D$, so $G_J^{R}=-G_D^{R}$ on the same source. The calculation below concerns $D$, not a change of sign in the Peierls convention. Cut the line at $x=0$ and set the trace to zero in the auxiliary half-line problems. Strike the left side at $(t,x)=(0,-a)$, $a>0$, with a unit impulse. Using the image source to enforce the Dirichlet trace, the independent retarded answers are
\[
 w_-^0=\tfrac12\{\Theta(t-|x+a|)-\Theta(t-|x-a|)\},\qquad w_+^0=0.
\]
The free kernel used for the half-line can itself be obtained by integrating the two characteristic transports of $D$. The two terms vanish against one another at $x=0$. Their outward \emph{spatial derivatives} on the two sides sum to $r(t)=-\delta(t-a)$: the first image term gives $-\tfrac12\delta(t-a)$ and the second gives another $-\tfrac12\delta(t-a)$. At a timelike face these outward derivatives have the opposite overall sign to the scalar action responses $\Pi$ in Chapter~\ref{ch:07}; their vanishing sum is the same transmission equation.

Now drive the common trace with a history $q(t)$. The homogeneous outgoing wave on either side is $q(t-|x|)$. The sum of its outward spatial derivatives is $2\dot q(t)$. In these reduced coordinates for the common trace, the mismatch equation and its \emph{retarded} solution are
\begin{equation}
 2\dot q=-r,\qquad q(t)=-\tfrac12\int_{-\infty}^{t}r(t')\,dt'
             =\tfrac12\Theta(t-a).\label{eq:retarded-trace}
\end{equation}
An arbitrary primitive would add a constant to $q$; the absence of a past disturbance removes it. At $x>0$ the detector reads $\tfrac12\Theta(t-a-x)$. On the left the correction cancels precisely the negative image term. The result on both halves is $\tfrac12\Theta(t-|x+a|)$: no premature arrival and no reflected pulse. The impulse is a distributional check; convolution with smooth compact pulses gives the corresponding solution on the smooth source domain. A source actually confined to $x=0$ belongs in the flux jump, not in one half's interior Dirichlet kernel.

\begin{figure}[htb]\centering
\begin{tikzpicture}[x=1.03cm,y=.77cm,>=Stealth]
\fill[pale] (0,0) rectangle (2.6,4.1);\draw[quiet,->] (-2.7,0)--(2.75,0) node[right] {$x$};\draw[marine,thick,->] (0,0)--(0,4.3) node[above] {$t$};
\draw[ochre,thick,-{Stealth}] (-1.65,.35)--(0,2);\draw[ochre,thick,-{Stealth}] (0,2)--(1.65,3.65);\draw[quiet,densely dashed] (0,2)--(-1.65,3.65);
\fill[ink] (-1.65,.35) circle (2pt) node[left] {source};\fill[ink] (0,2) circle (2pt);\fill[ochre] (1.65,3.65) circle (2pt) node[right] {detector};\node[below] at (0,0) {$0$};
\end{tikzpicture}
\caption{The dashed image wave of the auxiliary reflecting problem is canceled by the wave emitted from the corrected trace. Only the solid signal reaches the physical detector.}\end{figure}

The construction is not peculiar to a one-dimensional massless wave. For a planar cut in flat spacetime with a Klein--Gordon mass $m$, take a Laplace transform in positive time ($\Re z>0$) and a Fourier transform along the cut (momentum $k$). The outgoing/decaying half-space lift is $\widehat q\exp(-\kappa |x|)$, where $\kappa=\sqrt{z^2+|k|^2+m^2}$ is the branch with $\Re\kappa>0$. The two outward derivatives give $M=2\kappa$. Its inverse kernel in time at fixed $k$ is
\[
 \mathcal L^{-1}\!\left(\frac{1}{2\kappa}\right)
 =\tfrac12\Theta(t)J_0\!\left(\sqrt{|k|^2+m^2}\,t\right),
\]
where $J_0$ is the Bessel function of order zero. At $m=|k|=0$ this reduces to the primitive in \eqref{eq:retarded-trace}, with the overall minus furnished by the mismatch equation. On smooth causal data of suitable growth the transformed solution satisfies the equations, transmission and zero-past conditions. Time support follows from the Laplace prescription; \emph{spatial} finite speed also follows from the local energy estimate on backward cones, since the flux terms of the two internal faces cancel (Appendix~\ref{app:B}). A real-frequency reciprocal stripped of its branch prescription would not establish either causal inverse.

\section{A characteristic cut carries less transverse information}
For the metric wave operator set $u=t-x$, $v=t+x$. A null side $u=0$ is characteristic: the normal coefficient of the highest derivative vanishes. Indeed $4\partial_u\partial_v w=0$ has solutions $a(v)+b(u)$, and prescribing $w(0,v)$ alone does not determine $b(u)-b(0)$. Simply requesting the field and its conjugate flux on that side asks twice for related data. What is needed depends on the domain. In the future rectangle $0\le u\le U$, $0\le v\le V$, a second intersecting side supplies exactly the missing incoming information.

For a calculation that retains a mass term, prescribe $w(0,v)=a(v)$ and $w(u,0)=b(u)$ with $a(0)=b(0)=c$, and drive $(4\partial_u\partial_v+m^2)w=f$. Integrating once in each variable gives a causal integral equation:
\begin{equation}
 w(u,v)=a(v)+b(u)-c+\frac14\int_0^u d\xi\int_0^v d\eta\,[f(\xi,\eta)-m^2 w(\xi,\eta)].\label{eq:double-null}
\end{equation}
Its right side uses only the causal past rectangle of $(u,v)$. Iterating the integral part $k$ times bounds its operator norm on continuous data by $(m^2UV/4)^k/(k!)^2$: each ordered $u$-integral gives $U^k/k!$, and similarly for $v$. The successive integrals therefore converge uniformly on each finite rectangle, and iterating the same bound on the difference of two solutions proves uniqueness. With smooth compatible data the integral equation gives smoothness by differentiation. For $m=0$ it gives the explicit sum of the two traces minus their common corner value, plus the source integral. The shared corner was counted once, not twice.

This is a characteristic data problem on the stated domain, not a claim that every null hypersurface requires a second one. A complete cone with vertex conditions or an appropriate characteristic Cauchy surface calls for a different data theorem. More generally, a metric-null face is characteristic for a normally hyperbolic field; for a material field with a different principal symbol, the two notions can disagree. The spacetime atlas is common to all species, while the causal data problems need not be.

\begin{exercise}[A load at the cut]\label{ex:sheetpulse}
For $Dw=\delta(x)h(t)$ on the massless line with no past disturbance, compute the trace and the two outgoing fields using only the half-line lifts. What is measured at $x=b>0$ before a pulse supported at $t\geq0$ has traveled for time $b$?
\end{exercise}
\noindent\emph{Check.} Integration across $0$ gives $-[\partial_x w]=h$. The summed outward derivative is $2\dot q=h$, so $w(t,x)=\tfrac12\int_{-\infty}^{t-|x|}h(t')\,dt'$. The answer before the first light ray arrives is zero, independently of how the subsequent sheet history is chosen.
\chapter{Presymplectic sewing and the observables a cut can recover}\label{ch:09}

Draw a loop on a Cauchy surface so that half the loop lies in each of two regions. The integral of a local density over the surface can be split at the cut by ordinary additivity. The Wilson loop cannot: neither of the two pieces is a closed loop. What the cut must supply depends on what we want back. To add symplectic potentials we need compatible corner terms; to compare brackets we need the causal response; to \emph{build} the loop we need maps between endpoint fibres. These requirements can be tested separately.

\section{What adds on a matched solution}

Let $R_1,R_2$ meet along an artificial face $I$. Keep the physical wall actions and boundary conditions. Remove any action used only to close $I$, or cancel it explicitly against its counterpart. Work in a fixed gluing sector with a fixed comparison of the fields across the face. Write $\mathcal M$ for the pairs of opened regional solutions which satisfy the transmission equations and have the regularity required of the global field. There are three maps out of this space: the regional projections $r_1,r_2$ and the assembly map $\mathsf A:\mathcal M\to\widetilde{\mathcal P}_M$ to global solutions. We use a sector in which restriction and assembly are inverse at the level of fields in the identified bundles, as for the scalar in Theorem~\ref{thm:branch}. Gauge degeneracies have not been quotiented.

Let $\Sigma_a=\Sigma\cap R_a$ and $K_a=\Sigma_a\cap I$, oriented as part of $\partial\Sigma_a$. Write $\Theta_a$ for the bulk potential current, $C_a$ for the corner potential on the relevant face, and
\[
 \theta_a=\int_{\Sigma_a}\Theta_a-\int_{\partial\Sigma_a} C_a,
 \qquad \omega_a=\delta\theta_a .
\]
The minus sign is important: the final rim of a wall and the boundary of the final cap have opposite orientations (Appendix~\ref{app:A}; see also \cite{HW}). Suppose the bulk potentials are the restrictions of the same representative, the physical-wall corners agree with the global ones, and the two cut corners, transported to the same fibre and tangent variations, have the same $C$. The last condition holds, for example, if no integration by parts on the opened cut was needed and both $C$'s there are zero. With $d:K_1\to K_2$ the identifying map, the two internal contributions in the sum, evaluated on one matched variation, are
\[
 -\int_{K_1} C_1-\int_{K_2} C_2
        =-\int_{K_1}(C_1-d^*C_2)=0 .
\]
The bulk integrals add, and the physical corner integrals reassemble the actual boundary of $\Sigma$. We have proved, with the indicated choices,
\begin{equation}
 \mathsf A^*\theta_M=r_1^*\theta_1+r_2^*\theta_2,
 \qquad \mathsf A^*\omega_M=r_1^*\omega_1+r_2^*\omega_2 .
 \label{eq:09-sewn-form}
\end{equation}
The second equality uses $\delta$ commuting with pullback. If the cap actions differ by an endpoint functional $B$, the potentials can instead differ by $\delta B$ and give the same two-form. An arbitrary mismatch of corner potentials need not be exact. If the gluing map is itself varied, variations of that map enter the pullbacks; keeping its contributions while varying the fields is then part of the sewing problem.

There is a time-dependent check of this calculation. Write the opened boundary variation using a common positive face measure as $\langle\Pi_a,\delta q\rangle$, where $\Pi_a$ already includes its outward sign and the histories have been transported to one side. After the corner terms are included, the change of $\omega_a$ between two caps on regional solutions is minus the integral of $\langle\delta\Pi_a\wedge\delta q\rangle$ over the portion of cut between them. This follows by taking $\delta$ of the on-shell first variation; physical flux is assumed absent or accounted for by its wall terms. The sum vanishes since $\Pi_1+\Pi_2=0$ and a tangent variation satisfies $\delta(\Pi_1+\Pi_2)=0$. One open region can lose symplectic flux without there being any failure of conservation in the union.

\section{Which symmetries have charges?}

A symmetry is first checked off shell: its variation of the action is an endpoint difference $\alpha_X|_{f}-\alpha_X|_{i}$, possibly with a field-independent wall term. For a fixed generator $X$ preserving the data, in a representative with $\mathcal L_X\theta=\delta\alpha_X$ on the solution space, the Noether generator and its variation satisfy
\[
 H_X=\iota_X\theta-\alpha_X,\qquad
 \delta H_X=-\iota_X\omega .
\]
The latter identity is Cartan's formula; it fixes the contraction sign for $\omega=\delta\theta$. More generally the candidate charge one-form $-\iota_X\omega$ must be exact for a globally defined Hamiltonian. Closedness gives only local integrability on a contractible chart; periods can obstruct a single-valued charge. A flux through an unclosed wall may even obstruct closedness. Suppose a vector on $\mathcal M$ projects to compatible symmetry vectors on the regions and on the union. Contracting \eqref{eq:09-sewn-form} shows that the \emph{variations} of the charges add when regional Hamiltonians exist; their values add after fixing an additive constant on each connected sector. The sum of two nonintegrable regional variations can also be integrable on $\mathcal M$. Neither statement requires inverting a presymplectic form.

For example, apply the Maxwell charge of Chapter~\ref{ch:05} to a slice with a cut. Use a real potential $a_j$ and absorb the fixed spatial volume factor into the electric momentum density $\mathcal E^j$, so $\theta=\int\mathcal E^j\delta a_j\,\mathrm{d}^{n-1}x$. For a field-independent allowed transformation $X_\lambda a_j=\partial_j\lambda$, $X_\lambda\mathcal E^j=0$, one obtains on the Gauss constraint
\begin{equation}
 -\iota_{X_\lambda}\omega
 =\int_{\Sigma}\partial_j\lambda\,\delta\mathcal E^j\,\mathrm{d}^{n-1}x
 =\int_{\partial\Sigma}\lambda\,\delta E_{\rm out} .
 \label{eq:09-max-charge}
\end{equation}
Here $E_{\rm out}$ is the signed electric flux density; the interior term $-\int\lambda\,\partial_j\delta\mathcal E^j$ vanishes. A proper gauge direction is a redundancy of the chosen problem; infinitesimally it annihilates the presymplectic form on all allowed variations. Vanishing $\lambda$ on the boundaries is sufficient here. A cut transformation can instead have the charge $\int\lambda E_{\rm out}$ and is then active if that charge varies. Nonzero boundary value alone is not a test. Indeed, write $e_K=E_{{\rm out},1}=-E_{{\rm out},2}$ on the matched cut, using the same positive density convention. For parameters $\lambda_1,\lambda_2$ the two cut terms add to $\int_{K_1}(\lambda_1-\lambda_2)\,\delta e_K$. They cancel for a diagonal transformation; a relative transformation need not even preserve the gluing map. A charge on a genuine exterior boundary does not disappear by this argument.

\section{Comparing a bracket versus generating its arguments}

Fix a background solution and a class of differentiable, proper-gauge-invariant functionals whose first derivatives are admissible sources. Take the same physical boundary conditions in the global and regional response problems. Suppose a regional construction has produced $F_\#=F\circ\mathsf A$, $G_\#=G\circ\mathsf A$ on matched fields, with compatible source representatives and pairings, including any interface source. Solve the coupled retarded and advanced problems for the functional derivative $G'$. Under the transmission, regularity, support and uniqueness conditions of Proposition~\ref{prop:causal}, the assembled responses are $G_R G'$ and $G_A G'$ (up to proper gauge where appropriate). If $v_\#$ is the difference of the coupled responses, applying the chain rule gives
\begin{equation}
 \{F_\#,G_\#\}_{\#}
   =\delta F_\#(v_\#)
   =\langle F',(G_R-G_A)G'\rangle\circ\mathsf A
   =\{F,G\}_{\rm P}\circ\mathsf A .
 \label{eq:09-peierls}
\end{equation}
Proper-gauge ambiguity in $v_\#$ pairs to zero. This is not the sum of the brackets computed with two independently reflecting cuts. It is an equality for arguments already known to correspond. It gives no way, by itself, to choose regional generators for a proposed global argument.

\subsection{A source algebra for which generation can be proved}

Consider a real free scalar on a globally hyperbolic spacetime without physical boundary. Take $P=\Box_g-m^2$ (or add a smooth real potential), with the sign of the action Hessian. It is formally self-adjoint and normally hyperbolic. Choose as global target the unital real algebra of finite polynomials in
\[
 F_f(\phi)=\int_M f\phi\,\mathrm{d}V,\qquad f\in C_c^\infty(M),\qquad P\phi=0 .
\]
It is a Poisson algebra: the bracket on linear generators is the constant $\langle f,(G_R-G_A)g\rangle$. No completion and no product of fields at a point is intended.

Cut along a smooth separating timelike hypersurface. On each closed piece allow smooth tests compactly supported up to the face, not just tests supported in the interior. A regional linear observable pairs this test with the regional field, while the cut history is still allowed to vary. Form a real vector space $Q$ of pairs $(f_1,f_2)$ by requiring all jets to agree in a common signed collar coordinate, including the zeroth jet. Tests are scalar functions relative to the fixed smooth volume density. These conditions are imposed on the regional tests themselves. The piecewise function $f$ is smooth: matching the two one-sided values gives continuity and applying the same argument successively to derivatives gives smoothness to every order. Compact support follows from the finite union of the supports. Conversely restriction of any global test gives such a pair.

\begin{proposition}[Polynomial source reconstruction]
The homomorphism from the real symmetric algebra $\operatorname{Sym}(Q)$, sending a pair to the sum of its two regional smearings evaluated on matched solutions, has image the entire target just specified. Its kernel is the ideal generated by $P Q$, where $P$ acts piecewise with its smooth matched coefficients. In particular the linear kernel is precisely $P Q$.
\end{proposition}
\begin{proof}
By additivity of integration the image of a pair is $F_f$ for its piecewise assembly $f$. Every global generator is obtained. If $f=Ph$ with $h$ compactly supported and smooth, formal self-adjointness gives $F_f=0$. Conversely, suppose $F_f$ vanishes on every smooth solution. Put $E=G_R-G_A$ and test on the solutions $E k$ for arbitrary compact smooth $k$. Adjointness of retarded and advanced solutions yields $0=\langle f,E k\rangle=-\langle E f,k\rangle$. Hence $h=G_R f=G_A f$ satisfies $Ph=f$ and has support in $J^+(\operatorname{supp} f)\cap J^-(\operatorname{supp} f)$, compact by global hyperbolicity. Its restrictions give a pair in $Q$. This proves the linear assertion.

Any finite independent list of the resulting linear functionals on the vector space of solutions maps it onto $\mathbb R^m$, where $m$ is the length of the list. Otherwise a nonzero covector annihilating its image would give a linear relation. A polynomial vanishing on all solutions therefore vanishes as a polynomial in these independent coordinates. Choosing a basis in the finite span of the generators occurring in any proposed relation proves the ideal assertion. The source pairing descends through this ideal since $EP h=0$ for compact $h$.
\end{proof}

The source relation is \emph{joint}. On an open piece, integration by parts in $F_{P h_a}(\phi_a)$ leaves the outward flux of $\phi_a\nabla^\mu h_a-h_a\nabla^\mu\phi_a$ at the cut. The two such terms cancel for a matched $h$ and a smooth sewn field, not separately for arbitrary face jets. Declaring every $P h_a$ zero before sewing would impose relations not valid regionally. The proposition reconstructs an algebra from source pairs; computing its bracket by regional causal solves also needs the solvability and support in \eqref{eq:09-peierls}.

\section{Making a Wilson loop out of two open paths}

Let a based path cross the cut at $p,q$, travelling from $p$ to $q$ in $R_1$ along $\gamma_1$, then from $q$ back to $p$ in $R_2$ along $\gamma_2$. Take a finite-dimensional representation and let $V_{a,x}$ be its fibre on side $a$ at $x$. In a local frame write the connection as $\mathrm{d}+A_a$. The open transports are constructed \emph{in the two regions}: solve $\dot U=-A_a(\dot\gamma_a)U$, $U(0)=1$, along each path. At the endpoints this gives
\begin{align*}
 U_1&=\mathcal P\exp\!\left(-\int_{\gamma_1}A_1\right):V_{1,p}\to V_{1,q},\\
 U_2&=\mathcal P\exp\!\left(-\int_{\gamma_2}A_2\right):V_{2,q}\to V_{2,p} .
\end{align*}
Path ordering places later matrices on the left. The product $U_2U_1$ has no meaning until the endpoint spaces are compared. If $D_x:V_{2,x}\to V_{1,x}$ is the bundle identification, the endomorphism and the unnormalized Wilson trace are
\begin{equation}
 T=D_p U_2 D_q^{-1} U_1:V_{1,p}\longrightarrow V_{1,p},
 \qquad W=\operatorname{tr}_{V_{1,p}}T .
 \label{eq:09-cross-loop}
\end{equation}
Each intermediate index is contracted against its own fibre. To check the order, transform a section by $s'_a=g_a s_a$, so $A'_a=g_a A_a g_a^{-1}-(\mathrm{d}g_a)g_a^{-1}$. The initial-value equation gives $U'_1=g_{1,q}U_1g_{1,p}^{-1}$ and $U'_2=g_{2,p}U_2g_{2,q}^{-1}$. With $D'_x=g_{1,x}D_xg_{2,x}^{-1}$ the middle factors cancel in sequence:
\begin{align*}
 T'&=(g_{1,p}D_pg_{2,p}^{-1})(g_{2,p}U_2g_{2,q}^{-1})\\[-1mm]
   &\qquad\cdot(g_{2,q}D_q^{-1}g_{1,q}^{-1})(g_{1,q}U_1g_{1,p}^{-1})\\
   &=g_{1,p}Tg_{1,p}^{-1} .
\end{align*}
The trace, not either open matrix or generally the based matrix $T$, is invariant. The four vector spaces in Figure~\ref{fig:transport-cycle} make it possible to read the product without memorizing a convention for matrix order.

\begin{figure}[htb]\centering
\begin{tikzpicture}[>=Stealth,thick,font=\small]
\node (a) at (0,0) {$V_{1,p}$};\node (b) at (4,0) {$V_{1,q}$};
\node (c) at (4,-1.45) {$V_{2,q}$};\node (d) at (0,-1.45) {$V_{2,p}$};
\draw[->,ochre] (a)--node[above] {$U_1$} (b);
\draw[->,marine] (b)--node[right] {$D_q^{-1}$} (c);
\draw[->,ochre] (c)--node[below] {$U_2$} (d);
\draw[->,marine] (d)--node[left] {$D_p$} (a);
\end{tikzpicture}
\caption{The horizontal arrows are parallel transports measured within each region. The vertical arrows are identifications at one and the same endpoint, not a further segment of the physical path. Following the arrows once gives an endomorphism of $V_{1,p}$.}\label{fig:transport-cycle}\end{figure}

One may express $D_x=e_{1,x}e_{2,x}^{-1}$ using maps $e_{a,x}$ from a common reference fibre at $x$ into $V_{a,x}$. In these references the first arc is $e_{1,q}^{-1}U_1 e_{1,p}$, the second $e_{2,p}^{-1}U_2 e_{2,q}$; tracing their composition recovers \eqref{eq:09-cross-loop}. A simultaneous change of reference at an endpoint only changes coordinates on a contracted index. Such an edge reference records how to compare two experiments; supplying it has not given a new field an equation of motion. An extended regional phase space can assign a surface action or a flux pairing to edge variables \cite{DF}; what is dynamical then follows from that additional structure. Covariance under independent changes of local frame is not a declaration that all transformations at fixed boundary data are proper gauge. In fixed frames with fixed $D_x$, a symmetry preserving the identification must obey $g_{1,x}D_x=D_xg_{2,x}$.

\subsection{What is lost by keeping only closed regional transports?}

There is an exact answer in source-free Maxwell theory on a spatial circle. Normalize the charge of a probe to one, write the connection as $\mathrm{d}-i a$, and split the oriented circle into the two arcs above. In common endpoint frames put $z_i=\exp(i\int_{\gamma_i}a)$ and let $e_i$ be the constant electric momentum on arc $i$ (Gauss' law). Work over $\mathbb C$. Define the global target independently: finite Fourier sums in the circle holonomy $w$, with coefficients polynomial in the electric momentum $e$. Smooth initial data realize every $(w,e)\in U(1)\times\mathbb R$: take a constant connection with the desired integral modulo $2\pi$ and an arbitrary constant electric momentum.

The enriched regional input is the algebra of finite Laurent polynomials in the open $z_i$ and ordinary polynomials in $e_i$. At the two shared vertices the frame phases are $g_p,g_q$. Any such pair extends smoothly along either interval: choose a path in $U(1)$ constant near the ends. The two phases act with weights $g_q/g_p$ on $z_1$ and $g_p/g_q$ on $z_2$. Thus $z_1^m z_2^n$ has weight $(g_q/g_p)^{m-n}$ and the invariant subalgebra is precisely
\[
 \mathbb C[e_1,e_2,(z_1z_2)^{\pm1}] .
\]
The signed flux equation is $e_1-e_2=0$. The map from this invariant algebra sends $z_1z_2$ to $w$ and both $e_i$ to $e$. It is onto and its kernel is $(e_1-e_2)$: after substituting $e_2=e_1$, any remaining relation would be a finite Fourier polynomial vanishing for all $w$ and all real $e$. Taking its Fourier coefficients and then using polynomial uniqueness in $e$ rules that out. In a presentation with two symbols for $w$ and its inverse, the sole additional defining relation is $w w^{-1}=1$. The Poisson test can be made before and after matching. Locally $z_i=e^{i\alpha_i}$ on each arc, with potential $e_i\,\delta\alpha_i$ there. Hence $\{z_1z_2,e_i\}=i z_1z_2$ for either $i$; $e_1-e_2$ commutes with every invariant generator, so its ideal is Poisson. On the union $w=e^{i\alpha}$ and $\theta=\int e\,\delta a=e\,\delta\alpha$ give $\{w,e\}=i w$, exactly the induced bracket. This calculation concerns the stated Cauchy-data model. It does not assert admissibility of arbitrary unsmeared line sources in a continuum Peierls domain.

If, before sewing, one retains on each arc only transports invariant under \emph{independent} changes at its two ends, the only Laurent degree surviving there is zero. After flux matching the image is then only $\mathbb C[e]$. It cannot distinguish two smooth fields with the same $e$ and circle holonomies $+1$ and $-1$. The missing quantity was not a poor choice of bracket: its two open factors had already been discarded.

\begin{exercise}[A field-dependent gauge parameter]
On a two-dimensional chart of allowed boundary data $(b,e)$ suppose a proposed transformation gives the charge one-form $b\,\delta e$. Is $b e$ its Hamiltonian? \emph{Hint:} compute both $\delta(be)$ and $\delta(b\,\delta e)$ before restricting to the subspace $b=\text{constant}$.
\end{exercise}

\begin{exercise}[A third arc]
Insert a new marked point in the second arc. Using three independent vertex phases, find the invariant monomials in the three open $U(1)$ transports. \emph{Answer:} the zero-weight condition sets the exponents of incoming and outgoing edges equal at each vertex. The surviving monomials are the integer powers of the product around the full circle.
\end{exercise}
\chapter{Several cuts: junctions, elimination, and refinement}\label{ch:10}

Put a second artificial cut into a regional problem. An observer outside the region should not be able to detect which cut was erased first. There are two tests of that expectation. The first asks whether the same points and fibres were identified. The second asks whether a signal launched with the same source and outer conditions gives the same response. A correct point-set quotient passes only the first test. A calculation on the wave equation below passes both and exhibits the travel time that would be lost by a careless partial solve.

\section{The geometry at a junction}

Consider finitely many oriented pieces with faces, common edges of faces, and possibly deeper intersections. A \emph{stratum} here means a smooth piece of fixed incidence type: one can move on it without changing which faces meet. For example the four quadrants of the spatial plane times a time interval give four timelike half-faces meeting on a timelike line. The line is present in the bookkeeping even when the unsplit spacetime is flat and has no object there.

Fix the two-sided face identifications, their lifts to primitive bundles, the identifications on the remaining strata, and the local sector to be assembled. Require the identifications on a common stratum to be consistent: two routes to the same local germ give the same identification of that germ and its fibre. This tests transition maps, not parallel transport by a generally curved connection. Require compatible collar and junction charts which assemble to the intended smooth local models, with the required matched jets of smooth data. Assume the quotient is Hausdorff, has the desired time orientation, and carries the stipulated physical boundary and causal problem. Arbitrarily identifying a set of half-faces need not give a manifold: a cross section normal to an ordinary interior junction must be a disk with the indicated sectors, not two disks joined only at their centres. Metric jets must match for a smooth metric even when the topology is that of a disk.

Let $X$ be the disjoint union of the pieces. Colour one group of identifications blue and the rest red. Two points have the same image in the simultaneous quotient precisely if they are joined by a finite chain of steps of either colour. In the blue quotient the blue steps become equalities; the red steps descend as a relation \emph{generated by their images}. Every chain in the iterated quotient lifts to red steps interspersed with blue equivalences, so defines exactly the same final equivalence relation. By the defining property of the quotient topology, a set in the result is open just when its preimage in $X$ is open, in either order. The resulting homeomorphism is smooth with smooth inverse in the fixed assembled charts, and respects the bundle transitions by their consistency. This proves order independence of this geometry. It would not justify smoothing each intermediate corner in a newly chosen way or forgetting a still exposed transition map.

A \emph{defect} is additional physics supported on a stratum: for example a prescribed jump law, source, or independent dynamical field. It is not another name for the stratum itself. To see the distinction at the level of equations, vary the regional actions retaining the remaining faces as open. Use outward response densities and express identified traces in a common fibre. Up to the bulk and external terms the sum takes the form
\begin{equation}
 \sum_F\int_F\langle\Pi_{F,1}+\Pi_{F,2},\delta q_F\rangle
 +\sum_J\int_J\langle k_J,\delta z_J\rangle+\delta S_{\rm def} .
 \label{eq:10-junction-first-variation}
\end{equation}
Here $F$ ranges over paired face interiors, $z_J$ collects the compatible traces at a junction $J$, and $k_J$ is the \emph{signed sum} of terms arriving there after any tangential integrations by parts, expressed in the same trace variables. If an action uses more traces or jets they must also be included. Vary in face interiors first to obtain the flux equations. At $J$ vary only the traces that extend as a common admissible variation to the incident faces. If they are constrained, it is the restriction of the covector in \eqref{eq:10-junction-first-variation} to that variation space that must vanish. Treating each copy of $\delta z_J$ as independent can create fictitious equations. For the first-derivative bulk scalar with no face action the first integration by parts gives only face fluxes, no $k_J$. For a common corner potential on adjacent ordinary faces, equal pullbacks with opposite incidences cancel (Appendix~\ref{app:A}). Neither example demands an extra force at an artificial junction.

By contrast, on a timelike junction with proper time $t$, an independent coordinate $z$ with action $\int_J(\dot z^2-\mu^2z^2)\,\mathrm{d}t/2$ contributes $-(\ddot z+\mu^2z)\delta z$ to the integrand and $[\dot z\,\delta z]_i^f$ at its ends. If a coupling $\int_J c z q_J\,\mathrm{d}t$ is well defined on the chosen trace space, it also contributes $c q_J\,\delta z+c z\,\delta q_J$. The oscillator equation and trace law have changed; the endpoint term changes the potential. Sewing the same geometry but dropping these terms is a different theory.

\section{Keep the other cut open}

The actual equations on two cuts may be nonlinear. Denote them by $F_1(x,y)=0$ and $F_2(x,y)=0$ for independent coordinates $x,y$ on the joint trace domain; external data are held fixed. If the first equation has exactly one solution $x=\chi(y)$ for each admissible $y$ under consideration, the remaining equation is $F_2(\chi(y),y)=0$. Its solutions correspond to simultaneous solutions by $y\mapsto(\chi(y),y)$, not by $y\mapsto(\chi(0),y)$. If $x$ is multivalued, the relation of allowed $(x,y)$ and its branches must be retained. This elementary observation explains why the intermediate object is an \emph{open family} over the unsolved history. It must still respond when that history is changed. In particular a solution produced by closing all the other faces with zero trace is not that family.

There is an important qualification hidden in the word ``independent''. At a common edge the two face histories might have to agree. When a continuous lifting of the common trace exists, one can parametrize each face by a lifted common trace plus a remainder with zero edge trace, and retain the shared variable exactly once. More generally work on the joint compatible trace space without a splitting. The following block proof applies when honest product coordinates have been supplied for the block being eliminated and the variables kept. The block kept can include junction variables and their equations. It does not apply to independently varying incompatible corner values.

For a linear (or linearized) causal problem use the construction of Chapter~\ref{ch:08} with all traces prescribed. A particular response $G f$ to the collection of sources $f$ leaves a mismatch $Nf$ at the faces. A homogeneous lift $H\eta$ of the unknown trace history changes the mismatch by $M\eta$. Thus the equation to be solved is $M\eta=-Nf$, not just equality of boundary point values. In coordinates $\eta=(x,y)$ write
\begin{equation}
 A x+B y=b,\qquad C x+D y=d,
 \qquad \binom{b}{d}=-Nf .
 \label{eq:10-blocks}
\end{equation}
For example $B$ is the response measured on the first face when the history on the second changes with the first held fixed. It can be nonlocal in time. The blocks map histories $X,Y$ to mismatch spaces $U,V$: $A:X\to U$, $B:Y\to U$, $C:X\to V$, $D:Y\to V$. These names stand for specified trace and response domains, including regularity, compatibility at initial corners, gauge conditions if used, and support. There is no presumption that a Dirichlet parametrization valid for a timelike face works on a null or a Cauchy face.

\begin{proposition}[Two-stage causal solve]
Use the retarded trace spaces, with zero history before the source, and suppose all the displayed block operations act continuously on their stated domains. Suppose $A$ has a two-sided inverse there and $S=D-CA^{-1}B:Y\to V$ does also. Assume the blocks and inverses used propagate supports only into their allowed causal futures, with domains stable under the compositions. Then the simultaneous system \eqref{eq:10-blocks} has a unique retarded solution given by
\begin{equation}
 y=S^{-1}(d-CA^{-1}b),\qquad x=A^{-1}b-A^{-1}B y .
 \label{eq:10-staged}
\end{equation}
If the reverse pivots $D$ and $A-BD^{-1}C$ meet the corresponding assumptions, reversal gives the same response, including on all external faces left open. The advanced statement uses causal pasts instead.
\end{proposition}
\begin{proof}
At fixed $y$ the first equation gives $x=A^{-1}(b-B y)$. Substitution in the second gives $C A^{-1}b+(D-CA^{-1}B)y=d$. This proves necessity of \eqref{eq:10-staged}; substitution back into both rows proves sufficiency. Zero right-hand side gives $y=0$ and then $x=0$. Thus any reverse solve satisfying the two rows gives this same pair. A retained external measurement of the form $m_0(f)+m_1 x+m_2 y$ also agrees; so does the reconstructed bulk field $G f+H(x,y)$. Each operation is causal by hypothesis, and the composition of allowed causal paths is causal in the assembled spacetime. A matching solution with this support gives the global retarded response when assembly regularity and uniqueness of the global problem hold. In the notation of a single interface solve $L$ this is precisely $G-HLN$, now computed by two substitutions.
\end{proof}

For finitely many blocks apply the argument repeatedly, preserving the updated inhomogeneous term as well as the operator on every remaining history. On a restricted physical forcing range one can weaken surjectivity, but then $b-B y$ must actually be in the range of $A$ for the retained $y$, and $d-CA^{-1}b$ in the range of $S$. A right inverse alone leaves possible homogeneous solutions; either retain them or prove independence of the fields and observables being compared. An algebraic inverse need not be causal: the inverse of a time delay on unrestricted histories is a time advance. Even on its range it need not preserve retarded support relative to the applied force.

\begin{remark}[A bad pivot and a real zero mode]
Let both unknowns be histories on one time axis. The time-local equations $y=b$, $x+y=d$ have the unique causal solution $(d-b,b)$, but the coefficient of $x$ in the first row is zero. Eliminating $x$ by inverting that coefficient is impossible, although eliminating $y$ using the second row works. This is a block counterexample, not a failure of the geometric quotient. A different issue occurs for $x-y=b$, $-x+y=d$. The range requires $b+d=0$, and the common shift of $x,y$ is undetermined. It is not legitimate to set that shift to zero merely to write an inverse: it might be a gauge redundancy, an unprescribed incoming datum, or a physical junction mode. The identification and its effect on observables require separate information.
\end{remark}

\section{A refinement that remembers the travel time}

Return to the line of Chapter~\ref{ch:08}, with equation $u_{tt}-u_{\xi\xi}=f$, wave speed one, zero past disturbances and outgoing solutions at both ends. Put cuts at positions $\xi=0$ and $\xi=\ell>0$, with histories $q(t),r(t)$ respectively. Take both history and mismatch spaces to be $\mathcal H_0=\{h\in C^\infty(\mathbb R):h(t)=0\ \text{for}\ t\leq0\}$, with convergence of every derivative on compact time intervals. Start the Laplace calculation with smooth compact driving functions; the resulting time-domain maps below extend to $\mathcal H_0$. Measure the mismatch by the sum of the \emph{outward spatial derivatives}. For the action with kinetic density $(u_t^2-u_\xi^2)/2$ the variational coefficient is the negative of this sum. Multiplying both the homogeneous mismatch and that of the particular field by $-1$ leaves the interface equation unchanged.

For zero interior source and prescribed transforms $\widehat q,\widehat r$, the transformed middle field at position $\xi$ is the elementary solution of $u''=s^2u$,
\[
 \widehat u(\xi)=\frac{\sinh(s(\ell-\xi))}{\sinh(s\ell)}\widehat q
              +\frac{\sinh(s\xi)}{\sinh(s\ell)}\widehat r,
 \qquad 0\leq\xi\leq\ell .
\]
The two exterior solutions are $e^{s\xi}\widehat q$ for $\xi\leq0$ and $e^{-s(\xi-\ell)}\widehat r$ for $\xi\geq\ell$. At zero the interior outward derivative is $-\widehat u'(0)=s\coth(s\ell)\widehat q-s\operatorname{csch}(s\ell)\widehat r$, and the left exterior supplies $+s\widehat q$. At $\ell$ the interior derivative is $+\widehat u'(\ell)$ and the right exterior supplies $+s\widehat r$. Thus the matrix carrying the two histories to the two summed derivative mismatches is
\[
 M(s)=\begin{pmatrix}A&B\\ B&A\end{pmatrix},\qquad
 A=s(1+\coth(s\ell)),\qquad B=-s\operatorname{csch}(s\ell).
\]
Inhomogeneous bulk fields with fixed traces are accounted for in the two right-hand sides $\widehat b,\widehat d$. In particular they are not included a second time in the homogeneous map $M$.

Three identities do the entire elimination:
\[
 A^{-1}=\frac{1-e^{-2s\ell}}{2s},\qquad A^{-1}B=-e^{-s\ell},
 \qquad A-BA^{-1}B=2s .
\]
No pivot vanishes for $\operatorname{Re}s>0$. More is true than invertibility of numbers at each frequency. Applied to a causal function $b$, the first inverse gives $\tfrac12\int_{t-2\ell}^t b(\tau)\,\mathrm{d}\tau$ with zero extension to negative times. The factor $e^{-s\ell}$ gives a delay by $\ell$; $1/(2s)$ is half the integral from the past. The direct blocks also act causally: expanding $(1-e^{-2s\ell})^{-1}$ as a geometric series gives delayed derivatives with only finitely many arrivals on each bounded time interval. Thus all these maps are continuous on $\mathcal H_0$ and the inverse identities extend there: on any finite interval cut off the data smoothly beyond its end, where no change can affect the answer. The full solution, displayed in transform notation on data for which the transform exists, is
\begin{equation}
 \binom{\widehat q}{\widehat r}
 =\frac{1}{2s}
  \begin{pmatrix}1&e^{-s\ell}\\ e^{-s\ell}&1\end{pmatrix}
  \binom{\widehat b}{\widehat d} .
 \label{eq:10-two-cuts-line}
\end{equation}

Now view zero as a subdivision of the left half of a problem originally cut only at $\ell$. Eliminate $q$ while $r$ is still free. The surviving equation is $2s\widehat r=\widehat d+e^{-s\ell}\widehat b$. This is exactly the old one-cut operator: each of the two unbroken half-lines at $\ell$ contributes $s\widehat r$. A mismatch at the new face contributes to the remaining source only after time $\ell$.

Does the \emph{forcing} produced by a regional bulk experiment also reassemble correctly? Place an arbitrary smooth compact source $f(t,\xi)$ inside $0<\xi<\ell$ after time zero. Call the two coefficients of $\widehat q,\widehat r$ in the displayed middle field $h_0(\xi),h_\ell(\xi)$. The zero-trace particular solution $v$ in the middle satisfies $(s^2-\partial_\xi^2)v=\widehat f$. Integration by parts gives $\int_0^\ell h\widehat f\,\mathrm{d}\xi=[-h v'+h'v]_0^\ell$ for either coefficient. Thus $\widehat b=v'(0)=\int h_0\widehat f$ and $\widehat d=-v'(\ell)=\int h_\ell\widehat f$ are precisely the negatives of its outward derivatives. The combination $h_\ell+e^{-s\ell}h_0=e^{-s(\ell-\xi)}$ follows by substitution, or by solving the homogeneous equation with these two endpoint values. It gives the measured field at the retained cut:
\[
 \widehat r=\int_0^\ell \frac{e^{-s(\ell-\xi)}}{2s}\,\widehat f(s,\xi)\,\mathrm{d}\xi .
\]
This is the uncut answer. Indeed the full-line Green function for $s^2-\partial_\xi^2$ is $e^{-s|\xi-\xi'|}/(2s)$: it decays at both ends, is continuous, and its first derivative jumps by $-1$ at $\xi'$. In time, a signal from $\xi$ cannot arrive before $\ell-\xi$. Checking only that the homogeneous operator returned to $2s$ would not have tested this source dependence.

\begin{exercise}[Erase the other cut]
For the wave example retain $q$ instead of $r$. Write both the remaining source and the open lift of $r$ at arbitrary $q$. \emph{Answer:} $2s\widehat q=\widehat b+e^{-s\ell}\widehat d$ and $\widehat r=A^{-1}\widehat d+e^{-s\ell}\widehat q$.
\end{exercise}

\begin{exercise}[A branch lost by partial solving]
At a junction suppose the two equations are $x^2=y$ and $x+y=0$ for real traces. Compare the simultaneous solutions with a procedure retaining only $x=\sqrt y$. \emph{Hint:} solve $x^2+x=0$. The negative root is not a second description of the positive one; it gives a different matched trace.
\end{exercise}

\part{Sewing and quantization}
\chapter{What the linear commutator sews}\label{ch:11}
Two detectors, one on either side of a cut, can be used in two different
experiments. In the first, a source at one detector is varied and the
response at the other is recorded. In the second, the readings fluctuate in
repeated preparations of the same state and their correlation is measured.
For a linear bosonic field the first experiment fixes a commutator. It does
not in general predict the second experiment. Keeping the two questions
separate lets us quantize a sewn system even when neither a vacuum nor a
factorization of its Hilbert space has been specified.

\section{From physical linear observables to canonical commutators}
Fix the geometric identification, exterior boundary condition and classical
sector, and linearize about a solution in that sector. Tangents satisfying
the linearized equations, constraints and boundary conditions carry the
presymplectic form $\omega=\delta\theta$ of Chapter~\ref{ch:05}, including
its corner term. A proper gauge tangent pairs to zero with every allowed
tangent. A transformation with a variable surface charge generally does
not. Write $\widetilde{\mathcal S}$ for these constrained solution
tangents and $\mathcal G_{\rm proper}$ for their proper gauge directions.
The physical tangent space is $\mathcal S_{\rm phys}=
\widetilde{\mathcal S}/\mathcal G_{\rm proper}$. To call it symplectic
requires that on the specified sector no further radical of $\omega$
remains. Dividing by every transformation bearing a gauge parameter would
not establish that. A central charge fixed in the sector is a number; a
physical zero mode with a conjugate variable is still a canonical pair.
This distinction matters especially at a cut, where a charged transformation
on one piece can become a redundancy only after joining the pieces.

Here is a test-function realization that fixes the signs. For the real
scalar take $P=\Box-m^2-\xi R$, the Hessian with the sign of the action in
Chapter~\ref{ch:04}; $R$ here denotes scalar curvature. Use the metric
volume $dV$. Write $G_R,G_A$ for the retarded and advanced inverses
($G^R,G^A$ in Chapter~\ref{ch:06}) of this $P$, with the same physical
boundary condition and adjoint to one another. Put
\[
 E=G_R-G_A,\qquad F_f(\phi)=\int_M f\phi\,dV,\qquad
 E(f,h)=\int_M f(Eh)\,dV .
\]
Real sources $f,h$ are smooth and compactly supported in the interior. On
an ordinary globally hyperbolic spacetime without boundary, the label space
is $V=C_c^\infty(M;\mathbb R)/P C_c^\infty(M;\mathbb R)$: $F_{Pf}$
vanishes on solutions, and $E$ identifies this quotient with the real
solutions of spacelike-compact support \cite{WaldQFT}. At a physical wall
the test domain and Green identity must implement the wall condition. For
gauge fields one uses admissible linear observables that annihilate proper
gauge variations, modulo the equations and null observables. A propagator
of a gauge-fixed equation on all components is not by itself the physical
label space.

With $\theta=p\,\delta q$ on a canonical pair, the convention is
$\omega=\delta p\wedge\delta q$. Green's identity for the scalar gives
$\delta F_f(Y)=\omega(Y,Ef)$: $Ef$ is its Hamiltonian tangent. In
particular, the antisymmetric form on \emph{observable labels} is
\begin{equation}
 \sigma([f],[h])=\{F_f,F_h\}_{\mathrm P}
 =\delta F_f(Eh)=E(f,h)=-\omega(Ef,Eh).
 \label{eq:11-label-sign}
\end{equation}
The last minus sign trades the order of the arguments in the form on
solutions; it is not a change from retarded minus advanced. One check is a
spatial mode for which $P=-(\partial_t^2+\nu^2)$. Its causal kernel is
$-\sin[\nu(t-t')]/\nu$. Differentiating with respect to $t'$ at $t'=t$
gives $+1$, the bracket of displacement with its momentum.

\begin{definition}
For a real symplectic label space $(V,\sigma)$ the Weyl form of the
bosonic canonical commutation relations (CCR) defines the universal
algebra generated by unitaries $\mathcal W(v)$, $v\in V$, with
$\mathcal W(0)=\mathbf 1$ and
\begin{equation}
 \mathcal W(v)^*=\mathcal W(-v),\qquad
 \mathcal W(v)\mathcal W(w)=
 e^{-i\hbar\sigma(v,w)/2}\,\mathcal W(v+w).
 \label{eq:11-weyl}
\end{equation}
A regular representation is one in which $s\mapsto\mathcal W(sv)$ is
strongly continuous for every $v$. Its generator is written
$\mathcal W(sv)=e^{is\widehat X_v}$.
\end{definition}

The same relations make sense for a presymplectic space: labels in the
radical give central unitaries until their values are fixed or their labels
removed. The exponential relation avoids multiplying arbitrary unbounded
operators. In the field algebra, or by differentiation on suitable common domains,
$[\widehat X_f,\widehat X_h]=i\hbar E(f,h)\mathbf 1$; subscripts $f,h$
stand for their classes. No representation is needed to specify
\eqref{eq:11-weyl}. We set $\hbar=1$ from now on unless it is displayed.
If a positive-frequency basis is available and chosen, normalize physical
modes with $i\omega(u_I,\overline{u_J})=\delta_{IJ}$ and
$\omega(u_I,u_J)=0$. Then
\[
 \widehat\phi=\sum_I(a_Iu_I+a_I^\dagger\overline{u_I}),\qquad
 [a_I,a_J^\dagger]=\delta_{IJ}\mathbf 1 .
\]
For the oscillator this gives $u=e^{-i\nu t}/\sqrt{2\nu}$ and
$u(t)\bar u(t')-u(t')\bar u(t)=iE(t,t')$. Delta functions and spectral
integrals replace Kronecker symbols and sums for continuous labels. This is
one realization, not an instruction to choose positive frequency separately
on the two sides.

\section{Two routes to an algebra}
Suppose the classical assembly in the chosen sector gives the physical sewn
labels $V_{\rm sew}$ and a comparison $S:V_{\rm sew}\to V_M$ with the
independently specified whole. What licenses quantization of the comparison
is the bracket test
\[
 \sigma_M(Sv,Sw)=\sigma_{\rm sew}(v,w).
\]
For bulk sources this tests the causal propagator, including its blocks
between opposite sides of the cut. For instance, in the massless line
example of Chapter~\ref{ch:08}, a source at $x=-a$ and a detector at
$x=b$, $a,b>0$, are separated by a travel time $a+b$. For $t>t'$ the
transmitted kernel of the Hessian $P=-(\partial_t^2-\partial_x^2)$ gives
\[
 [\widehat\phi(t,b),\widehat\phi(t',-a)]
 =-\tfrac i2\,\Theta(t-t'-a-b)\mathbf1
\]
as a distribution on that time ordering. The minus converts the positive
wave operator used for the pulse to the action Hessian. Two independent
reflectors would give zero even after the first arrival. Matching just the
spaces of homogeneous solutions does not compute this block.

If $S$ is a symplectic isomorphism, $\mathcal W_{\rm sew}(v)\mapsto\mathcal W_M(Sv)$
is a star isomorphism: it preserves the phase in \eqref{eq:11-weyl}, and
$S^{-1}$ supplies the inverse. This is the \emph{sew, then quantize} route.
If the comparison covers only a symplectic subspace, the same argument only
identifies its generated subalgebra.

The converse construction starts with quantum generators of opened pieces.
What it means to impose a matching equation on them is not automatic. Even
for canonical pairs, setting $q_1-q_2=0$ in the full polynomial operator
algebra would set $[q_1-q_2,p_1]=i\mathbf 1$ to zero. The quotient by the
two-sided ideal so obtained would be the zero algebra. One must first find
observables that survive the matching. The following linear case gives a
precise meaning to that instruction.

\begin{proposition}
Let $U$ be a real symplectic space of extended linear labels and let
$I\subset U$ be isotropic: $\sigma(i,j)=0$ for $i,j\in I$. Suppose the
matched presentation is $I^\perp/I$, where $I^\perp$ consists of labels
pairing to zero with $I$. In the Weyl algebra of $I^\perp$, impose
$\mathcal W(i)=\mathbf1$ for $i\in I$. The resulting algebra has precisely
the Weyl relations of $I^\perp/I$ with its induced form. It is a
\emph{symplectic} presentation of the physical labels if the comparison
identifies their form and
$\operatorname{rad}(\sigma|_{I^\perp})=(I^\perp)^\perp=I$. Otherwise
the quotient can still be presymplectic.
\end{proposition}
\begin{proof}
In the indicated subalgebra the constraint unitaries are central. Since
$\mathcal W(v+i)=\mathcal W(v)\mathcal W(i)$ for $v\in I^\perp$, the
class of $\mathcal W(v)$ depends only on $v+I$. The phase of a product
also depends only on the two classes. This gives a map from the Weyl
algebra on the quotient labels; sending each old generator to the Weyl
generator of its class defines the inverse map. The radical of the induced
form is $(I^\perp)^\perp/I$: isotropy gives $I\subset I^\perp$ and hence
$(I^\perp)^\perp\subset I^\perp$. The double-orthogonal condition is
automatic in finite symplectic dimension, not for every infinite-dimensional
choice of subspace. Residual central Weyl generators must not silently be
set to one.
\end{proof}

For $I$ generated by $q_1-q_2$, the surviving total momentum is
$p_1+p_2$, not either $p_1$ or $p_2$ alone. With
$Q=(q_1+q_2)/2$ one finds $[Q,p_1+p_2]=i\mathbf1$: no factor of two is
lost when the displacements are identified. The proposition treats a
linear first-class constraint in its zero sector. A nonzero fixed charge
requires its specified central character; noncommuting (second-class)
constraints generally change the bracket of the survivors. A retarded
boundary-history problem need not admit this constraint presentation at
all. In particular, the direct sum of the label spaces of two independently
\emph{reflecting} theories carries their reflecting commutators. It does
not acquire a transmitting commutator by a choice of quantum state. A
quantize-first route must actually construct the physical label space and
prove the bracket comparison on its claimed domain.

\section{Extending correlations, not just commutators}
A state $\varpi$ is a normalized positive linear functional on the algebra;
positivity means $\varpi(A^*A)\ge0$. A centered quasifree (Gaussian)
state is specified by a real symmetric covariance $\mu$:
\begin{align}
 \varpi(\mathcal W(v))&=e^{-\mu(v,v)/2},\nonumber\\
 W(v,w):=\varpi(\widehat X_v\widehat X_w)
 &=\mu(v,w)+\tfrac i2\sigma(v,w) .\label{eq:11-covariance}
\end{align}
A mean would add its linear phase to the first formula and the product of
the means to the second. Higher centered moments are sums over pairings of
$W$, in the order of the fields. The necessary and sufficient Gaussian
positivity condition is that the Hermitian form $\mu+i\sigma/2$ be positive
semidefinite on the complexification of $V$ \cite{WaldQFT,Haag}. Indeed,
its quadratic form on $z=\sum c_j v_j$ is the expectation of
$(\sum c_j\widehat X_{v_j})^*(\sum c_j\widehat X_{v_j})$. In particular
\[
 \mu(v,v)\mu(w,w)\ \ge\ \mu(v,w)^2+\tfrac14\sigma(v,w)^2 .
\]
With $\hbar$ restored, replace $\sigma$ by $\hbar\sigma$ in these state
formulas. The commutator fixes the imaginary part; even a completely known
bracket therefore leaves a positivity problem for the real part.

For example, suppose two admissible regional subspaces have been embedded
in $V_M$. Seek a single covariance on $V_M$ whose restrictions are the
given $\mu_1,\mu_2$. They must agree on common data. In a presentation by
regional generators, the trial cross block $C(v_1,v_2)$ must make the
whole symmetric form descend through every relation: a label representing
zero must pair to zero with \emph{all} labels. For finite families of test
labels on the two sides the positivity test is the matrix inequality
\[
 \begin{pmatrix}
   \mu_1+i\sigma_{11}/2 & C+i\sigma_{12}/2\\
   (C+i\sigma_{12}/2)^* & \mu_2+i\sigma_{22}/2
 \end{pmatrix}\ \ge 0 .
\]
Here the $\sigma_{ab}$ are the blocks of the already sewn commutator, and
$*$ is conjugate transpose; the test must hold for every such family. If
the labels span only part of $V_M$, a positive extension to the rest is
still needed. For a field state the resulting two-point functional must
also be a distribution with the equations, boundary conditions and
short-distance behavior required of the whole theory. Equality of two
marginals alone performs none of these tests.

Take two commuting oscillator pairs ordered as $(Q_1,P_1,Q_2,P_2)$.
Write $J_2=\left(\begin{smallmatrix}0&1\\-1&0\end{smallmatrix}\right)$
and $1_2$ for the identity matrix. Consider
\[
 \sigma=J_2\oplus J_2,\qquad
 \mu_c=\begin{pmatrix}\nu 1_2&c1_2\\c1_2&\nu 1_2\end{pmatrix}.
\]
The sum and difference coordinates $(Q_1\pm Q_2)/\sqrt2$ and their
similarly normalized momenta give two canonical pairs with covariances
$(\nu+c)1_2$ and $(\nu-c)1_2$. Positivity holds precisely when
$\nu-|c|\ge1/2$. Thus $\nu=1$ admits both $c=0$ and $c=1/4$, with
identical individual Gaussian states and identical commutators, but
$\varpi(Q_1Q_2)=c$. A putative extension with $c=3/4$ fails the uncertainty
test in the difference pair. This finite calculation proves neither
nonuniqueness nor existence for every spacetime cover. It isolates exactly
which matrix entry the two regional commutator experiments did not measure.

\section{A sharp spatial cut and the vacuum}
There is a local version of this diagnostic that directly tests a proposed
product vacuum. In four-dimensional Minkowski space the massless scalar
vacuum has, at equal times and distinct spatial points,
\[
 W((t,\mathbf x),(t,\mathbf y))=
 \frac{1}{4\pi^2|\mathbf x-\mathbf y|^2},\qquad
 E((t,\mathbf x),(t,\mathbf y))=0 .
\]
The formulas may be smeared over two disjoint small balls on opposite
sides of a spatial plane. The first pairing is positive for nonzero
nonnegative test functions; the second is zero. The product of the two
centered marginal states would predict zero for both. As the balls approach
the same point from opposite sides, the missing correlation is the leading
short-distance singularity of the field, not an arbitrary convention at
one point. The Hadamard condition requires this local singularity on the
\emph{uncut} spacetime. It does not pick the vacuum: adding a smooth,
positive kernel $c\,u(x)u(y)$ with $c\ge0$ and $u$ a real smooth allowed
solution changes the Gaussian state while leaving the commutator and the
singularity unchanged. Appendix~\ref{app:C} gives the distributional test.

There is a stronger obstruction than the fact that one particular state is
entangled. In the vacuum representation of standard continuum free fields,
local von Neumann algebras for ordinary bounded open diamonds are type III
\cite{Haag}. In particular they are not full operator algebras
$\mathcal B(\mathcal H_1)$ of a Hilbert-space factor: the latter have
minimal projections and are type I. Thus a sharp spatial division does not
in general give a physical tensor decomposition in which each local algebra
is the full algebra of its factor, and a partial trace of a density matrix
is not the definition of restriction to a local algebra. Restriction is
simply evaluation of $\varpi$ on its subalgebra. An abstract Hilbert-space
isomorphism, or the tensor product for finitely many lattice sites, does
not resolve the operator-algebra obstruction. With a separating collar,
models satisfying the split property admit an intermediate type I factor;
shrinking the collar is not a proof of sharp factorization. Nor does a
spacelike division of initial data make successive \emph{time} intervals
independent quantum systems.

An algebra isomorphism always transports states by composition. To identify
two \emph{specified} cyclic state representations, one must also compare
the states: for centered Gaussian states transported by $S$, this means
$\mu_{\rm sew}(v,w)=\mu_M(Sv,Sw)$. If it holds, mapping the cyclic vectors
and the vectors generated from them preserves inner products and gives the
corresponding unitary after completion. A scalar commutator or a matched
classical flux alone does not give this unitary.

\begin{exercise}
Keep the two local covariances $\nu 1_2$ of the oscillator example, but
change the cross block to $\operatorname{diag}(c,-c)$. How large can
$|c|$ be at fixed $\nu\ge1/2$? Can a local pair with vacuum covariance
$\nu=1/2$ have any such correlation?
\end{exercise}
\noindent\emph{Hint.} In the sum pair the displacement variance is
$\nu+c$ while the momentum variance is $\nu-c$; in the difference pair
they are interchanged. Apply Gaussian positivity to each entire pair.
\chapter{Perturbations about a sewn solution}\label{ch:12}
Send in two weak pulses. A linear field gives the sum of the two signals,
but an interacting field can respond to their overlap. The first coefficient
sensitive to that overlap is a mixed derivative with respect to the source
strengths. Its calculation makes a good test of sewing: the pulses, the
vertex where they meet, and the detector need not be on the same side. We
will compute the coefficient twice, first from a classical initial-value
problem and then as a tree term of finite-time quantum evolution.

\section{Choosing the background and the free problem}
Let $\bar\phi$ be a solution on the whole slab, or a matched regional pair
proved to assemble to such a solution. Fix the outer data and geometric
sector. Write $\phi=\bar\phi+\eta$ in a field chart. A source expansion
will vary the bulk driving force while holding initial data fixed. Thus the
response $\eta$ has zero incoming Cauchy data, although the background
need not have zero data. With the exterior conditions held fixed, and for
variations supported away from the caps, the first variation of the joined
action vanishes on $\bar\phi$. This includes the cancellation of its two
outward responses. Otherwise the proposed fluctuation theory would start
with a source on the seam even before the pulses are sent.

The quadratic action is one half of the second variation evaluated twice
on $\eta$; its operator is the Jacobi operator at $\bar\phi$. Its domain includes the linearized transmission and
outer conditions. Higher variations give interaction vertices, including
any allowed physical boundary vertices. Linearizing independently around
two convenient regional vacua need not give these operators: their fields
or normal momenta might not be the traces of one solution. For a gauge
theory a Green operator may require a gauge choice; the external
observables must still descend to the physical theory of
Chapter~\ref{ch:11}. We now choose a scalar self-interaction for which the
mixed response can be found without any gauge assumption.

Use a fixed background metric, $P=\Box-m^2-\xi R$, and add to the free
action the terms
\begin{equation}
 -\int_M\left(\frac{\lambda}{3!}\chi\phi^3+j\phi\right)dV .
 \label{eq:12-source-action}
\end{equation}
Here $\chi$ is smooth, $j$ is an applied bulk source, and $\lambda$ counts
cubic vertices. Use compact switching where end effects or spatial support
matter. The action sign fixes the equation
$P\phi-\lambda\chi\phi^2/2=j$. Suppose $\bar\phi$ solves it at $j=j_0$,
with the transparent scalar matching of Chapter~\ref{ch:04}. For
$j=j_0+s f+r h$ the \emph{exact} equation for the fluctuation is
\begin{equation}
 P_{\bar\phi}\eta=s f+r h+\tfrac{\lambda}{2}\chi\eta^2,
 \qquad P_{\bar\phi}=P-\lambda\chi\bar\phi .
 \label{eq:12-jacobi}
\end{equation}
The term linear in $\eta$ from expanding the interaction belongs to the
free operator about this solution. The quadratic coefficient is
$\tfrac12\int\eta P_{\bar\phi}\eta\,dV$ when the integrations by parts
have no uncancelled surface terms; the remaining interaction is cubic in
$\eta$. Since the interaction in \eqref{eq:12-source-action} contains no
derivatives, it does not change the scalar trace or the outward normal
momentum. This is one reason for the choice of model.

\section{The mixed response and its regional construction}
Let $G_R^{\bar\phi}$ be the retarded inverse of $P_{\bar\phi}$ on this
problem. Assume the finite-time initial-boundary problem is well posed for
the smooth sources used here and has the requisite derivatives at zero;
alternatively read the following as a formal recursion in $s,r$ with a
well-defined linear inverse. Write $u_f=\partial_s\eta|_0$,
$u_h=\partial_r\eta|_0$, $w_{fh}=\partial_s\partial_r\eta|_0$.
Differentiating \eqref{eq:12-jacobi} and imposing zero initial data at each
order gives
\begin{align}
 u_f&=G_R^{\bar\phi} f,\qquad u_h=G_R^{\bar\phi}h,\nonumber\\
 P_{\bar\phi}w_{fh}&=\lambda\chi u_f u_h,\qquad
 w_{fh}=\lambda G_R^{\bar\phi}\!\left(\chi u_f u_h\right).
 \label{eq:12-response}
\end{align}
There are two ways to assign the incoming derivatives to the two factors
of $\eta^2$. They cancel the $1/2$ in its coefficient. Thus the coefficient
of $sr$ in $\eta$ is $w_{fh}$, whereas the coefficient of $s^2$ alone is
$\lambda G_R^{\bar\phi}(\chi u_f^2)/2$. The expression retains the
background dependence of the inverse; it is not an assertion that the
background is of order zero in an independent expansion in $\lambda$.

For a numerical sign and normalization check, take a massless flat patch,
$\bar\phi=0$, and $\chi=f=h=1$ in the causal past of the observation
between the initial cap and elapsed time $T$. Cutoffs outside that domain
do not change this retarded calculation. With zero initial perturbation
the solution there has no spatial dependence. Its first derivative is
$u_f(T)=u_h(T)=-T^2/2$, because $P=-\partial_T^2$ on this field. The
second derivative is
\[
 w_{fh}(T)=-\frac{\lambda}{4}\int_0^T(T-t)t^4\,dt
          =-\frac{\lambda T^6}{120} .
\]
In particular the nonlinear signal is not obtained by simply multiplying
the two first-order readings at the detector. Their product acts as a new
source; it must propagate from its point of formation to the detector.

To reproduce \eqref{eq:12-response} from regional data, first compute
$u_f,u_h$ by solving the linearized response problem on the pieces. On a
timelike scalar cut this means varying their common boundary history until
the two outward normal responses cancel. Do the same for the source
$\lambda\chi u_f u_h$ at the next step. Why should the result equal a
whole-slab response? For this scalar, continuity kills a delta derivative
in the wave operator of a piecewise field; equality of the two normal
derivatives in a common signed coordinate kills the delta itself. The
piecewise field then solves the whole inhomogeneous equation weakly. With
the prescribed past and outer data, uniqueness in the solution class gives
the claimed response. This uses the reconstruction and Green-operator
hypotheses of Chapters~\ref{ch:07} and~\ref{ch:08}; it is not a rule for
manipulating Heaviside functions times quantum fields.

There is a useful way to compare the two computations directly without
placing a singular Green kernel on the cut. Smear the output with a smooth
detector profile $k$. By adjointness of retarded and advanced inverses,
\begin{equation}
 \int_M k\,w_{fh}\,dV
 =\lambda\sum_{a=1,2}\int_{M_a}
 (G_A^{\bar\phi}k)\,\chi\,u_f u_h\,dV .
 \label{eq:12-sewn-tree}
\end{equation}
The restriction to each piece is now an integral of smooth fields. The
advanced solution carries the detector profile back to the vertex; the
retarded solutions carry the sources forward. If $f$ and $h$ lie in
different pieces, replacing the transmitted fields by independently
reflecting fields can set the overlap to zero. With the correctly joined
Green operators the two integrals instead add to the whole answer, without
an independent interaction assigned to the artificial wall. Were the
interface equation itself nonlinear, differentiating it would produce an
additional inhomogeneous interface condition for $w_{fh}$; it is absent
here by the particular action and domain.

\section{Finite-time evolution and the same tree}
Choose the same free fluctuation problem, with an admissible state if an
expectation is wanted. Write $\widehat\eta_I$ for the interaction-picture
free field. For this nondifferentiated potential the interaction Hamiltonian
has the opposite sign to the added Lagrangian density. For a foliation with
$dV=N\,dt\,d\Sigma_t$ its insertions are
\[
 V_I(t)=\int_{\Sigma_t} N\,d\Sigma_t\left[
  \frac{\lambda}{3!}\chi\, :\widehat\eta_I^3:
     +(s f+r h)\widehat\eta_I\right].
\]
The colons denote Wick products in a specified subtraction. In the field
theory $V_I(t)$ denotes an operator-valued distribution in time obtained
from spacetime insertions, integrated with smooth switching; the notation
does not assert that a sharp-time Wick cube is an operator on a Hilbert
space. No stationary Hamiltonian or choice of particles at infinity is
needed: the free linear evolution identifies the slice algebras. With
$\hbar=1$, the Dyson series on the finite slab is
\begin{align}
 U(t_f,t_i)&=\mathcal T\exp\!\left(-i\int_{t_i}^{t_f} V_I(t)\,dt\right)
 \nonumber\\
 &=\mathbf 1-i\int_{t_i}^{t_f}V_I(t)\,dt
 -\frac{1}{2!}\int_{t_i}^{t_f}dt\int_{t_i}^{t_f}dt'\,
   \mathcal T\{V_I(t)V_I(t')\}+\cdots .
 \label{eq:12-dyson}
\end{align}
Later times stand on the left. In a field theory this is an order-by-order
formula for smeared, renormalized insertions, not an operator-norm
convergence claim. One may use switching supported in the interior of the
slab. If an actual time-dependent Hamiltonian evolution is defined,
$U(t_f,t_i)=U(t_f,t_*)U(t_*,t_i)$ follows by splitting the ordered
integration domains. For local perturbations the corresponding statement
is causal factorization for suitably ordered supports. Neither operation
makes two spatial regions independent, and neither takes an infinite-time
scattering limit \cite{FR}.

The quantity appropriate to our detector is a response in the specified
initial state, namely the expectation of
$U(t_x,t_i)^*\widehat\eta_I(x)U(t_x,t_i)$, understood after smearing.
At first order in $V_I$, its change is
$i\int_{t_i}^{t_x}[V_I(t),\widehat\eta_I(x)]\,dt$. For the source term,
$i[\widehat\eta_I(y),\widehat\eta_I(x)]=E_{\bar\phi}(x,y)\mathbf1$.
With retarded ordering the result is $u=G_R^{\bar\phi}(s f+r h)$.
All higher source commutators of this c-number vanish: in the free theory
a source simply displaces the field by $u$. To keep that displacement
exact while counting cubic vertices, factor out the evolution $U_j$
generated by the source alone, writing $U=U_j U_{\rm rel}$. The cubic
insertion in $U_{\rm rel}$ is then conjugated by $U_j$: its field at $y$
is $\widehat\eta_I(y)+u(y)$. Use one Wick prescription for the powers,
so that the ordinary binomial rule holds for a shift by a smooth function.
The commutator needed at first cubic order is
\[
 i[:\!(\widehat\eta_I(y)+u(y))^3\!:,\widehat\eta_I(x)]
 =3 E_{\bar\phi}(x,y):\!(\widehat\eta_I(y)+u(y))^2\!: .
\]
The three identical terms cancel a factor of three in $3!$. With retarded
ordering the one-vertex correction to the mean field is therefore
\begin{equation}
 \frac{\lambda}{2}\int_M G_R^{\bar\phi}(x,y)\chi(y)
  \left\langle :\!\big(\widehat\eta_I(y)+u(y)\big)^2\!:\right\rangle dV_y .
 \label{eq:12-displaced-tree}
\end{equation}
Equivalently, insert the source-only solution into the quadratic source on
the right side of the renormalized field equation at first vertex order.
For a centered free state, the term quadratic in source amplitudes is just
$u(y)^2$. Its mixed derivative is $2u_fu_h$, so
\eqref{eq:12-displaced-tree} reproduces \eqref{eq:12-response} with both
its sign and its factor of $\lambda$. A local expectation of the Wick
square, if nonzero in the chosen subtraction, is source independent at
this order. It is not part of the two-pulse response.

\section{What survives the comparison at the next order?}
A classical response graph is drawn with its output at the root and a
retarded line for each inversion of $P_{\bar\phi}$. Equation
\eqref{eq:12-sewn-tree} has one three-legged vertex and no closed chain.
Iterating the classical equation gives more vertices but still trees. In a
quantum calculation two fluctuations at vertices can instead be contracted
with each other. Two lines joining the same pair of vertices already make
one loop. The distinction is not the power of $\lambda$: two cubic
vertices can occur both in a tree and in a loop. For example, with $n$
vertices and $I$ internal lines a connected graph has $L=I-n+1$ independent
loops. One way to keep the counting straight is to expand a connected
functional with external classical fields held fixed: internal contractions
carry $\hbar$, action insertions carry $1/\hbar$, and the overall
$\hbar$ multiplying the connected logarithm gives $\hbar^L$. The tree
response is of order $\hbar^0$ in this counting.

In a centered quasifree reference state Wick's rule reduces free-field
expectations in the Dyson expansion to pair contractions
$G_F(x,y)=\langle\mathcal T\widehat\eta_I(x)\widehat\eta_I(y)\rangle$,
without a further factor of $i$, once the coincident products are defined.
For a general initial state higher connected correlations are additional
data. The time-ordered contraction is not the retarded Green operator used
in a classical response. Even away from coincident vertices, an ordinary
quantum correlator computed in independently chosen regional states need
not have the correct cross contraction. At coincidence the product of two
lines in a loop can also require a local extension. Agreement of the
transmitted Hessian and all the classical vertices gives the tree response
derived here; agreement of a quantum coefficient further requires the
chosen state's contractions, the same operator and subtraction
normalizations, and compatible time-ordered products including contact
terms. Chapter~\ref{ch:13} works out the first example where a genuine
extension choice enters. The one-vertex calculation supplies a comparison
at that order, not an all-orders commutation of quantization with sewing.

\begin{exercise}
Couple a third source $b\,l$ to the same cubic problem. Derive the fully
mixed third derivative of the retarded classical field at zero source in
terms of the first responses and $w_{fh}$ of \eqref{eq:12-response}. How
many distinct terms remain when the three source labels are different?
\end{exercise}
\noindent\emph{Hint.} Differentiate $\eta^2$ three times at $\eta=0$.
One factor must receive one derivative, the other two. Pair the two possible
orders of these factors before applying $G_R^{\bar\phi}$.
\chapter{Time-ordered products at a cut}\label{ch:13}
In the mixed response of Chapter~\ref{ch:12} the point where two pulses
meet is integrated against smooth solutions. A pair of quantum vertices is
different. Even if each vertex has already been defined as a Wick
polynomial, contracting two lines between the vertices asks for the square
of a singular propagator. Putting the vertices on opposite sides of a drawn
line does not remove the issue: both points can approach the same point of
the line. What does the regional calculation have to reproduce there?

\section{Causal ordering fixes the separated-vertex answer}
A time-ordered product $T_n(F_1,\ldots,F_n)$ is a symmetric multilinear
operator-valued distribution in its bosonic local insertions. More
precisely, its kernels are distributions; the displayed product denotes
their evaluation on smooth compactly supported couplings. An important
part of its definition is \emph{causal factorization}. If the supports of
a first group do not meet the causal past of the supports of the remaining
group, then
\begin{equation}
 T_n(F_1,\ldots,F_n)=
 T_k(F_1,\ldots,F_k)\,T_{n-k}(F_{k+1},\ldots,F_n).
 \label{eq:13-causal}
\end{equation}
The product on the right is the ordinary quantum product, with the later
group to the left. For spacelike-separated groups the opposite ordering is
also possible and the products commute. For two fields in a specified
state, the expectation of this ordered operator product is $W(x,y)$ when
$x$ is later, and $W(y,x)$ when $y$ is later;
at spacelike separation the two agree. It does not say to put all the
vertices on the left of a \emph{spatial} division later than all those on
the right. With lower time-ordered products known, causal factorization
constructs the next one off the total coincidence diagonal in an ordinary
globally hyperbolic bulk. Extending to the diagonal is a separate step
\cite{BF,HWT,FR}.

For instance, introduce smooth parts $\chi_1+\chi_2=\chi$ of the
switching function, supported in two overlapping neighborhoods of a cut.
Let $F_a=\lambda\int\chi_a\,:\!\phi^3\!:\,dV/3!$. The quadratic Dyson
term for $F_1+F_2$ is
\[
 -\tfrac12 T_2(F_1,F_1)-T_2(F_1,F_2)-\tfrac12 T_2(F_2,F_2).
\]
The cross term has coefficient one, not one half. Its contractions require
the two-point function \emph{between} neighborhoods. Matching the
commutator and the two separately renormalized regional Wick cubes need
not supply it. A smooth partition is a way of labeling the vertices. It
introduces neither a new boundary condition nor a discontinuous coupling;
the answer for a transparent cut must be independent of the partition when
the pieces are added back together.

\section{One loop with its normalization}
Use a four-dimensional flat interior patch to see the extension explicitly.
Take the massless scalar vacuum and use its normal ordering in this
calculation. For $z=x-y$ set $z^2=-(z^0)^2+|\mathbf z|^2$. In units
$\hbar=1$, our time-ordered contraction is
\[
 G_F(z)=\frac{1}{4\pi^2}\frac{1}{z^2+i0},\qquad
 \Box G_F=i\delta^{(4)} .
\]
The $i0$ is a boundary value, not a small real cutoff. To check the source
sign, spatial Fourier transformation gives
$e^{-i|\mathbf p|\,|t|}/(2|\mathbf p|)$, whose first time derivative has
jump $-i$ at $t=0$; $\Box$ contains $-\partial_t^2$.
Appendix~\ref{app:C} gives the slightly different boundary value for the
unordered two-point function.

Take two cubic vertices at $x,y$. Choosing two of the three fields at each
point and pairing them across the vertices gives
$\binom32^2 2!=18$ double contractions. Away from coincidence this part of
the second-order operator in \eqref{eq:12-dyson} is
\begin{equation}
 -\frac{\lambda^2}{4}
 \int dV_x\,dV_y\,\chi(x)\chi(y)\,
  G_F(x-y)^2\,:\!\phi(x)\phi(y)\!: .
 \label{eq:13-fish}
\end{equation}
Indeed the coefficient before counting pairings is
$(-i\lambda)^2/[2!(3!)^2]$. The remaining fields in
\eqref{eq:13-fish} have \emph{not} been contracted with external
detectors. This is the two-line loop, or fish. On the punctured relative
space $z\ne0$ its scalar distribution is
$t^\circ(z)=(16\pi^4)^{-1}(z^2+i0)^{-2}$. This notation includes the
noncoincident light cone, where the boundary value is singular but has a
consistent frequency direction. Under $z\mapsto a z$, $a>0$, the degree
is $-4$. The normal space to $x=y$ has dimension four, so the divergence
is logarithmic. An ordinary integral using the pointwise formula at the
origin would not define the coefficient.

Here one can write an extension rather than just count the free parameters.
Take $\ell>0$ of dimension length and set
\begin{equation}
 t_\ell(z)=-\frac{1}{64\pi^4}\Box
 \left[\frac{\log((z^2+i0)/\ell^2)}{z^2+i0}\right].
 \label{eq:13-fish-extension}
\end{equation}
Off the origin the bracket is defined by the same complex boundary value;
its scaling degree is two, so it extends uniquely there without increasing
its scaling degree. The derivative is then taken on distributions. To see
why the coefficient works, for $s=z^2$ one has
$\Box h(s)=4s h''(s)+8h'(s)$. For $h(s)=\log(s/\ell^2)/s$ this is
$-4/s^2$; the identity continues to the prescribed boundary value away
from $z=0$. Hence \eqref{eq:13-fish-extension} agrees with $t^\circ$ and
has scaling degree four. Two such power-counting extensions can differ by
$c\,\delta^{(4)}(z)$, but not by derivatives of that delta. In particular
the displayed family gives
\[
 t_{\ell'}-t_\ell=\frac{i}{8\pi^2}
       \log(\ell'/\ell)\,\delta^{(4)}(z),
\]
by $\Box(1/(z^2+i0))=4\pi^2 i\delta^{(4)}(z)$. Inserting a contact
difference $c\,\delta^{(4)}$ into \eqref{eq:13-fish} changes the full
operator by $-c\lambda^2\int\chi^2:\!\phi^2\!:\,dV/4$. Its mixed part
for the smooth partition $\chi=\chi_1+\chi_2$ is
\[
 -\frac{c\lambda^2}{2}\int_M\chi_1\chi_2:\!\phi^2\!:\,dV .
\]
This contact is exactly zero for disjoint supports, even though the
nonlocal fish contraction can still connect the two supports. The bulk
mass normalization matters for overlapping supports and for the full
coefficient, not as a new sheet of mass on the cut. Reality and unitarity
further restrict admissible constants in the chosen propagator convention.

\section{Power counting in the normal directions}
The fish displays the general geometry of power counting. Near the diagonal
of $n$ bulk vertices in spacetime dimension $D$, hold the common point
fixed and scale the $k=(n-1)D$ relative coordinates to zero. If the
relative distribution has scaling degree $\rho$, an extension preserving
that degree is unique for $\rho<k$. Otherwise the local ambiguity consists
of derivatives of a delta in the relative coordinates of orders
$|\alpha|\le\lfloor\rho-k\rfloor$; tangential coefficients are restricted
by the normalization axioms. In a scalar graph in $D>2$, $I$ singular
internal lines of leading degree $D-2$ and $b$ derivatives on such lines
give the superficial bound
\[
 \rho-k\ \le\ (D-2)I+b-D(n-1)=D L-2I+b
\]
for a connected cluster with $L=I-n+1$. This counts allowed local terms;
tensor cancellations can lower the actual degree. Proper subclusters must
first have their products defined, and their extension choices propagated
to larger clusters. On a curved smooth bulk the short-distance counting
is local, although allowed coefficients can now involve the metric, mass
and curvature \cite{HWick,HWT}.

In a geometry with distinguished faces the collision set can itself have
strata. For the stratum with all $n$ bulk points coincident \emph{on} a
face of codimension $q$, the normal dimension in $M^n$ is
$(n-1)D+q$: one must bring the common point to the face as well as bring
the points together. Vertices constrained from the start to the face have
a different configuration space, and hence a different normal dimension.
On a smooth stratum with a normal scaling-degree bound, possible additional
ambiguities are normal derivatives of deltas supported there, subject to
that bound. If the coefficients themselves are singular at an edge or
corner, that smaller stratum needs its own analysis; one cannot count a
bulk delta once and call it a boundary counterterm. This reasoning
classifies local freedoms \emph{if} the respective distributions and
extensions exist. Bulk theorems on smooth globally hyperbolic manifolds do
not supply a renormalization theorem for every proposed boundary problem.

For the flat fish just computed, write $X=(x+y)/2$. The uncut distribution
is independent of the coordinate $X_\perp$ normal to a planar cut.
Under joint scaling of $(z,X_\perp)$ its degree is still four, whereas
$\delta^{(4)}(z)\delta(X_\perp)$ has degree five. Thus the latter is not
an additional ambiguity with this scaling bound. More generally, an
\emph{artificial} wall is absent from the whole theory's local background
data. Regional subtractions, including any auxiliary sheet terms needed by
a regulator, must combine into the same whole-spacetime prescription at a
transparent seam. A surviving wall coupling that changes correlators when
the plane is moved fails this comparison. A \emph{physical} defect, by
contrast, may have a prescribed action, a jump equation or its own fields.
The face and its geometry then are local data of the theory. Face, corner
or junction operators allowed by power counting and symmetries must be
fixed by the physical normalization conditions when needed. Their
coefficients cannot be inferred by demanding that this defect disappear.

\section{Ward identities are more than matching conditions}
Time-ordered products obey differential identities with contact terms.
For the free scalar, $P_xW(x,y)=0$ but $P_xG_F(x,y)=i\delta(x,y)$.
Differentiating the time ordering has detected the equal-time commutator.
Treating time-ordered fields as classical solutions at the coincident
point would miss it. At zero coupling the flat massless scalar has a shift
symmetry $\phi\mapsto\phi+\epsilon$, with current
$j^\mu=-\partial^\mu\phi$ and $j^0=\partial_t\phi$. In the interior the
Ward identity for one field insertion is
\[
 \partial_\mu^x\langle\mathcal T j^\mu(x)\phi(y)\rangle
 =-\Box_x G_F(x-y)=-i\delta^{(4)}(x-y).
\]
The sign is that of $[j^0(t,\mathbf x),\phi(t,\mathbf y)]
=-i\delta^{(3)}(\mathbf x-\mathbf y)$. Integrating the identity over a
small spacetime neighborhood of $y$ detects that the inserted field shifts
by one. A cubic potential breaks this shift explicitly; that breaking is
not a quantum anomaly. For other symmetries the contact terms involve the
variations of the insertions, together with any physical boundary fluxes.
Matching the smooth classical fluxes alone does not test the quantum
contact terms.

Two failures of a local comparison have different diagnoses. The two routes may use different allowed local terms, such as
different mass normalizations in \eqref{eq:13-fish}. Aligning the
prescriptions can remove that discrepancy. Or a symmetry condition may
itself be obstructed. In the usual local perturbative framework, when the
Ward identities hold at lower orders, their first possible local violation
$\mathcal A$ obeys a consistency condition. For a gauge theory written
with a nilpotent symmetry differential $s$ this reads $s\mathcal A=0$
modulo the relevant total derivatives. Adding an admissible local
counterterm $B$ shifts the violation by $sB$ in that quotient. At an actual
boundary a total derivative has a boundary term; it is the combined bulk
and boundary identity that must be tested. A nontrivial class, not
removable by any allowed $B$, is a quantum anomaly of that Ward identity,
not an arbitrary failure to match two choices of subtraction. Conversely,
absence of such a class does not furnish a state, a transmission Green
operator or equality of every pair of renormalized products. These are
separate inputs to an order-by-order comparison \cite{FR}.

\section{Fixed sectors and an optional weighted sum}
Suppose the classical gluing problem allows several inequivalent global
sectors: for instance, transport around a loop through the seam may depend
on transition data which cannot be removed by a gauge transformation.
Choose one such sector $s_0$, its exterior conditions and its fluctuation
problem. For insertions in the common observable domain let $\iota_{s_0}$
be the comparison from the matched regional algebra to the whole algebra,
including transport of any crossing observable. The \emph{strong,
fixed-sector} assertion at order $n$ is an equality there, for every
admissible collection of insertions,
\begin{equation}
 \iota_{s_0}\!\left(T^{\rm sewn}_{n,s_0}(F_1,\ldots,F_n)\right)
 =T^M_{n,s_0}(\iota_{s_0}F_1,\ldots,\iota_{s_0}F_n).
 \label{eq:13-fixed}
\end{equation}
Here $F_i$ are matched labels for local insertions on which both products
are defined. For numerical correlators one uses corresponding states as
well. The preceding fish calculation identifies one local ambiguity that
must be fixed for this assertion even at $n=2$; it is not a proof of
\eqref{eq:13-fixed} at all orders. In a fixed flux or fixed bundle sector
there is no instruction in causal factorization to sum over others.

A different observable may deliberately leave sector data unobserved.
First map the sector answers to a common target, for example complex
amplitudes $A_s$ with the same fixed external data. If a \emph{chosen}
measure is a finite groupoid counting measure with weights $w_s$, the
combined amplitude takes the form
\[
 A_{\rm sum}=\sum_{[s]}\frac{w_s}{|\operatorname{Aut}(s)|}\,A_s .
\]
The meaning of the denominator can be checked without a path integral. If
a finite gauge group $G$ acts on a set of presentations $X$, then dividing
$\sum_{x\in X}w_x A_x$ by $|G|$ counts an orbit $Gx$ with coefficient
$|Gx|/|G|=1/|G_x|$ for invariant $w_x A_x$. Here $G_x$ is the stabilizer,
namely the automorphism group of that presentation. If the weights are
already a measure on equivalence classes, one does not divide a second
time. For continuous labels a measure and treatment of residual symmetry
volumes must be supplied; there is no default convergent sum. A statistical
mixture of states would instead require nonnegative normalized weights.
Summing complex amplitudes and mixing states are different experiments.

For a numerical example, measure the holonomy $H=ba$ in the two-strip
circle of Chapter~\ref{ch:02}. Each of the two classes $H=+1,-1$ has a
diagonal $\mathbb Z_2$ stabilizer, so unit groupoid weights give each
class counting mass $1/2$. The normalized mean of the measured sign is
$[(+1)/2+(-1)/2]/[1/2+1/2]=0$. In a fixed periodic sector the same
measurement gives $+1$, and in a fixed antiperiodic sector it gives $-1$.
Zero is the mean for the chosen ensemble, not a third possible holonomy.

Equality in each fixed sector implies equality of common weighted sums
when those sums are defined with the same measure. Equality only after a
sum is weaker, since sector discrepancies can cancel. In particular,
releasing a common classical trace to obtain transmission, contracting
quantum insertions across a cut, and summing over sectors held fixed in
\eqref{eq:13-fixed} are three different operations.

\begin{exercise}
Suppose an order-two distribution in flat four-dimensional spacetime is a
single relative derivative $\partial_{z^\mu}t_\ell(z)$ of the renormalized
fish, constant in $X$. At the coincident-on-plane stratum, what sheet
contact term is no longer ruled out by the scaling-degree count alone?
Why does that count not allow you to introduce a new freely adjustable
sheet coupling in the transparent theory?
\end{exercise}
\noindent\emph{Hint.} Scale $z$ and $X_\perp$ together. A derivative can
raise the degree by one. Then ask what additional data a coefficient with
index $\mu$ would need, and whether differentiation of the already fixed
whole distribution leaves such a choice.

\appendix
\cleardoublepage\phantomsection\addcontentsline{toc}{part}{Appendices}
\chapter{Orientations and calculus at a cut}\label{app:A}

An outward flux has a sign even when the face has no unit normal. A corner has a sign too, but it depends on which face one took first. Both facts can be read off from a volume form. Throughout this appendix $n$ is the spacetime dimension, the metric has signature $(-,+,\ldots,+)$, and $\epsilon$ is the volume form for a chosen spacetime orientation. The initial and final slices, when named as slices, are both oriented as the boundary of their past. Thus it is $\Sigma_f-\Sigma_i$, not $\Sigma_f+\Sigma_i$, that occurs in the oriented boundary of a time slab.

\section{Which way does the corner point?}

On a smooth face $F$ of a region $R$, choose a defining function $s$ with $s\leq 0$ inside, $s=0$ on $F$, and $\mathrm{d}s\ne0$ there. If $N$ is any transverse vector with $\mathrm{d}s(N)>0$, a tangent frame $(e_1,\ldots,e_{n-1})$ of $F$ is positive precisely when $(N,e_1,\ldots,e_{n-1})$ is positive in $R$. This outward-first rule is the one in Stokes' formula. There is no metric in it. In particular at the top of a slab one can use $N=\partial_t$; raising the outward covector $\mathrm{d}t$ in Minkowski space would instead give $-\partial_t$.

At the transverse intersection of two faces choose outward vectors $N_1,N_2$, each tangent to the other face, in a corner chart. For example if the coordinate inequalities are $s_1,s_2\leq0$, use $N_i=\partial_{s_i}$. The two successive boundary orientations are represented by
\[
 \iota_{N_2}\iota_{N_1}\epsilon
       =-\iota_{N_1}\iota_{N_2}\epsilon
\]
on their common tangent space. Thus every codimension-two face of an ordinary manifold with corners occurs twice with opposite incidence signs in $\partial(\partial R)$. It does not vanish as a set; it vanishes in the oriented sum. If one sums over only some faces, there is no reason for their boundary to vanish. For example, if $\partial R=\Sigma_f-\Sigma_i+\Gamma$ with $\Gamma$ the entire lateral wall, then
\begin{equation}
 \partial\Gamma=-\partial\Sigma_f+\partial\Sigma_i .
 \label{eq:app-a-wall-ends}
\end{equation}
When $\Gamma$ itself is divided into faces, the right side does not include their common edges: those cancel within $\partial\Gamma$. For face-dependent integrands there may nevertheless be an edge term, since the two values being integrated need not agree.

Here is a sign check that works in any $n\geq2$. Take $t_i\leq t\leq t_f$, $x\leq0$, with volume $\mathrm{d}t\wedge\mathrm{d}x\wedge\epsilon_Y$ in a product chart. At the corner $t=t_f$, $x=0$, the final cap carries $\mathrm{d}x\wedge\epsilon_Y$ and its boundary carries $+\epsilon_Y$. The wall $x=0$ carries $-\mathrm{d}t\wedge\epsilon_Y$ and its final rim carries $-\epsilon_Y$. Performing the same computation at $t_i$ reverses the cap's sign as a face of $R$. In two spacetime dimensions $\epsilon_Y$ here is simply $1$; the calculation determines the sign of an endpoint in the elementary integration-by-parts formula.

\section{Use a flux density, not a unit null normal}

Write $\nu=\mathrm{d}s|_F$ for the outward conormal as a covector in the ambient cotangent space, and normalize a transverse vector by $\nu(N)=1$. If $j:F\hookrightarrow R$ is inclusion, put $\epsilon_\nu=j^*\iota_N\epsilon$. Replacing $N$ by $N+T$ with $T$ tangent does not change this form: evaluation would involve $n$ tangent vectors to an $(n-1)$-dimensional space. Decomposing any flux vector $V$ on the face as $\nu(V)N+T$ gives
\begin{equation}
 j^*\iota_V\epsilon=\nu(V)\epsilon_\nu,
 \qquad \mathrm{d}\Sigma_\mu=\nu_\mu\,|\epsilon_\nu| .
 \label{eq:app-a-flux-density}
\end{equation}
The first formula is for integration of forms with the induced orientation; the second uses a positive density and places the outward sign in the covector. If $s$ is replaced by $a s$, $a>0$, its differential on the face is $a\nu$ and $\epsilon_{a\nu}=a^{-1}\epsilon_\nu$. Consequently $\mathrm{d}\Sigma_\mu$ does not depend on this rescaling, even when $a$ varies along $F$. On a non-null face one can normalize $\nu$ to have squared norm $\varepsilon=\pm1$. Then $N=\varepsilon\nu^\sharp$ has the required pairing and $|\epsilon_\nu|$ is the induced metric volume density. On a null face $\nu^\sharp$ is tangent, not transverse. Nothing in \eqref{eq:app-a-flux-density} requires raising $\nu$.

On opposite copies of one cut there are two equivalent ways to add fluxes. Pull a single flux \emph{form} back to both copies and integrate over their opposite boundary orientations. Alternatively, integrate two outward flux \emph{densities} with the same positive measure: it is then the outward conormals which are opposite. Inserting an extra orientation reversal into the density sum counts the sign twice. For the scalar Lagrangian $-\tfrac12(\nabla_\mu\phi\nabla^\mu\phi+m^2\phi^2)\epsilon$ the potential flux vector is $V^\mu=-\nabla^\mu\phi\,\delta\phi$. In a flat chart with orthonormal $t,x$, the potential on a future cap is $+\partial_t\phi\,\delta\phi$ per unit spatial volume. At the wall $x=0$, with $x\leq0$ to the left and $x\geq0$ to the right, the sum of the outward densities per unit wall volume is
\[
 (-\partial_x\phi_{\rm left}+\partial_x\phi_{\rm right})\,\delta q,
 \qquad q=\phi_{\rm left}|_{x=0}=\phi_{\rm right}|_{x=0} .
\]
Thus releasing one trace enforces equality of the two derivatives taken along the \emph{same} $x$ direction, not equality of the two outward derivatives.

For comparison take a null chart with metric $-2\,\mathrm{d}u\,\mathrm{d}v+h_{AB}\,\mathrm{d}y^A\mathrm{d}y^B$, oriented by $\mathrm{d}u\wedge\mathrm{d}v\wedge\epsilon_h$, and keep $u\leq0$. Here $\nu=\mathrm{d}u$, $N=\partial_u$, $\nu^\sharp=-\partial_v$, and $\epsilon_\nu=\mathrm{d}v\wedge\epsilon_h$. The same scalar potential has outward coefficient $V^u=+\partial_v\phi\,\delta\phi$. This derivative runs \emph{along} the generator. One can see exactly what it misses. In flat $1+1$ dimensions the massive equation $(-2\partial_u\partial_v-m^2)\phi=0$ restricts to a transport equation for the transverse derivative $b(v)=\partial_u\phi(0,v)$. If $q(v)=\phi(0,v)$, it gives
\[
 b(v)=b(v_0)-\tfrac12 m^2\int_{v_0}^{v}q(v')\,\mathrm{d}v' .
\]
Thus equal traces fix the change in $b$ \emph{along} the cut but not its starting value. A datum on a transversely intersecting null side can fix $b(v_0)$ (Chapter~\ref{ch:08}); without it even equal potential flux need not give a smooth join. The massless kink in Chapter~\ref{ch:07} is the special case in which the undetected difference is constant along the entire generator.

\section{From a wall derivative to a cap potential}

Let $\mathrm{d}$ denote the spacetime exterior derivative and $\delta$ the exterior derivative on field space, with $\mathrm{d}\delta=\delta\mathrm{d}$. Write $\delta L=E_a\delta\phi^a+\mathrm{d}\Theta$, so $\Theta$ is a spacetime $(n-1)$-form linear in variations. Keep the wall Lagrangian $\ell$ in $S=\int_R L+\int_\Gamma\ell$. Suppose the allowed wall variations obey the Harlow--Wu condition \cite[eqs.~(2.20)--(2.27)]{HW}
\[
 (\Theta+\delta\ell)|_\Gamma=\mathrm{d}_\Gamma C .
\]
The corner potential $C$ is a spacetime $(n-2)$-form and a field-space one-form. Applying \eqref{eq:app-a-wall-ends} after integrating the wall derivative, one finds
\begin{equation}
 \begin{aligned}
  \delta S&=\int_R E_a\delta\phi^a+\theta_{\Sigma_f}-\theta_{\Sigma_i},\\
  \theta_\Sigma&=\int_\Sigma\Theta-\int_{\partial\Sigma}C,
  \qquad \omega_\Sigma=\delta\theta_\Sigma .
 \end{aligned}
 \label{eq:app-a-cap-potential}
\end{equation}
This assumes no separate cap action. Adding $B_f-B_i$ to the action adds $\delta B$ to the cap potential, but not to its two-form. There is a second useful freedom: replacing $\Theta$ by $\Theta+\mathrm{d}Y$ and $C$ by $C+Y|_\Gamma$ leaves \emph{the potential itself} unchanged, by Stokes on $\Sigma$. Changing just $C$ at an artificial cut is not this operation.

One can check conservation without extending $C$ into the interior. On a solution and a pair of linearized solutions, $\mathrm{d}\delta\Theta=-\delta E_a\wedge\delta\phi^a$ vanishes. The pullback of $\delta\Theta$ to the wall need not vanish: it equals $\mathrm{d}_\Gamma\delta C$ on the allowed variations. Stokes on the slab and then on its wall consequently gives $\omega_{\Sigma_f}=\omega_{\Sigma_i}$. If the boundary history $q$ is opened instead, the face one-form may contain a response term $\Pi\,\delta q$ in addition to $\mathrm{d}_\Gamma C$. Its field variation is $\delta\Pi\wedge\delta q$, and the same computation leaves precisely this flux to be matched. Ignoring a surviving flux would make the answer depend on where the cap was drawn.

\begin{exercise}[A tilted defining function]
For the wall $x=0$ in the product chart, replace $x$ by $e^{b(t,y)}x$. Compute the wall density of a vector $V$ anew. \emph{Hint:} use $N=e^{-b}\partial_x$ on the wall. The factor multiplying $V^x$ in the conormal cancels the one in the surface element.
\end{exercise}

\begin{exercise}[An unmatched corner]
Two adjacent cap pieces induce opposite orientations on $K$. Suppose that in the same corner orientation their potentials obey $C_1-C_2=p\,\delta q$ there, with $p,q$ independent data. What remains in the sum of the two cap two-forms? \emph{Answer:} $-\int_K\delta p\wedge\delta q$. It vanishes on a fixed-$q$ slice of field space, which is not a justification for dropping it after $q$ is released.
\end{exercise}
\chapter{Traces, support, and well-posed wave problems}\label{app:B}
An inverse of a differential operator is not specified by a formula alone. The functions it acts on, the boundary conditions of the functions it returns, and its support in time are part of the inverse. This appendix gives practical tests for those three pieces of information. The elementary estimates are given for a scalar; they are not substitutes for the energy and constraint estimates of a gauge system.

\section{What can actually be evaluated at a hypersurface?}
For a smooth field, restriction to $s=0$ is evaluation. For a square-integrable field it is not: changing its values on a set of measure zero changes no $L^2$ equivalence class. The local Sobolev space $H^r$ measures $r$ derivatives in a mean-square sense; in a coordinate patch its norm is given by the Fourier weight $(1+|\xi|^2)^{r/2}$. This is an auxiliary positive norm on the coordinates, not the indefinite norm of a Lorentz metric.

On a flat hyperplane $s=0$, write $k$ for the Fourier variable tangential to the plane and $\lambda$ for the one dual to $s$. The tangential Fourier transform of $u|_0$ is, up to the Fourier normalization, $\int \widehat u(k,\lambda)\,d\lambda$. Cauchy--Schwarz bounds its square by
\[
 \left(\int_{\mathbb R}(1+|k|^2+\lambda^2)^{-r}\,d\lambda\right)
 \left(\int_{\mathbb R}(1+|k|^2+\lambda^2)^r|\widehat u(k,\lambda)|^2\,d\lambda\right).
\]
The first integral is $C_r(1+|k|^2)^{1/2-r}$ for $r>1/2$. Multiplying by $(1+|k|^2)^{r-1/2}$ and integrating over $k$ proves the loss of half a derivative:
\[
 \|u|_0\|_{H^{r-1/2}}\leq C\|u\|_{H^r},\qquad r>1/2.
\]
Localization and a smooth extension of half-space Sobolev fields give the corresponding local result in a collar. For the classical trace of $\partial_s u$ the same argument needs $r>3/2$. It is a sufficient condition, not the only way a normal flux can exist. For example on a bounded Lipschitz coordinate patch a flux vector density $p$ with $p$ and $\partial_\mu p^\mu$ in $L^2$ has a weak normal trace in $H^{-1/2}$: pair a test trace $q$ with
\[
 \int_R \big((\partial_\mu p^\mu)\,\widetilde q+p^\mu\partial_\mu\widetilde q\big),
\]
where $\widetilde q\in H^1$ extends $q$. Integration by parts on zero-trace extensions proves independence of the extension; the extension bound controls the pairing. These facts justify a weak mismatch equation in appropriate norms. They do not upgrade it to matching infinite jets. At a corner one cannot assign mutually incompatible traces on intersecting faces. Starting with traces of actual smooth fields, or specifying the appropriate compatibility subspace in a completion, avoids treating those assignments as independent.

\section{Sources, trace histories, and causal inverses}
Let $P$ be a wave operator with specified physical boundary conditions on a time-oriented spacetime. $J^+(K)$ consists of points reachable from $K$ by future-directed causal curves in the specified spacetime, allowing the physical boundary as appropriate. On smooth compact sources a retarded Green operator must satisfy
\[
 PG^R f=f,\qquad G^R P v=v,\qquad \operatorname{supp} G^R f\subseteq J^+(\operatorname{supp} f).
\]
The test field $v$ in the middle identity must satisfy the boundary domain of $P$, have $Pv$ in the source domain, and itself have zero past Cauchy data. On a finite slab a smooth field supported away from both caps is a common test for both inverses; mere compactness up to the caps is not enough. Reverse $+$ to $-$ and past to future for $G^A$. On finite slabs, take sources vanishing near the past cap for a smooth zero-past solution, or impose the corner compatibility conditions of forced initial data; the advanced problem uses the future cap instead. Retarded and advanced do \emph{not} mean that data are prescribed independently at both caps. Existence on a globally hyperbolic spacetime without timelike boundary for a normally hyperbolic operator is a standard input \cite{BGP}; a timelike physical boundary brings a separate initial-boundary value problem.

Here is one way to specify the spaces used in Chapter~\ref{ch:08}. In a static half-space with the scalar wave operator, take $f$ smooth of compact support consistent with zero past; for the elementary half-line check one may start with supports away from the artificial corners. Take smooth trace histories $\eta$ that vanish in the distant past, locally bounded with all derivatives on finite time intervals. Impose compatible support in the directions parallel to the cut, or work with compact spatial wave packets and causal completion. On the massless half-line the lift $H\eta(t,x)=\eta(t-|x|)$ is explicitly continuous on these smooth causal histories. On the corresponding smooth past-zero mismatch space $M=2\partial_t$ is onto: $Ly(t)=\tfrac12\int_{-\infty}^t y(s)ds$ is continuous on each fixed finite interval with a common past cutoff. It need \emph{not} return a history compact in time. Insisting on compact support for the response as well as for the source would incorrectly impose $\int y=0$. For more general equations these spaces and bounds have to be supplied; an everywhere-defined right inverse of abstract vector spaces does not give them. Energy norms often give the first existence result and smooth compatible data a second regularity result.

\section{Flux cancellation and finite speed}
For $D=\partial_t^2-\nabla^2+m^2$, $m^2\geq0$, multiply $Du=f$ by $u_t$. For a real field the elementary conservation law is
\[
 e=\tfrac12\big(u_t^2+|\nabla u|^2+m^2u^2\big),\qquad
 \partial_t e-\nabla\!\cdot(u_t\nabla u)=f\,u_t .
\]
Integrating on a spatial region with a static boundary gives an outward-derivative term $\int_{\partial}u_t\partial_nu$ on the right of the energy balance. At an artificial cut the two $u_t$ equal $\dot q$, and the outward derivatives sum to zero. The internal energy contributions cancel. At a fixed homogeneous Dirichlet or Neumann outer wall the same power vanishes by the corresponding boundary condition. A passive or dynamical boundary may instead require an estimate or its own stored energy. On a closed system with zero initial data the integrated energy estimate gives uniqueness. It gives continuous dependence in the energy norm by estimating $\int f u_t$; existence still needs a construction, such as mode expansion on the line, or the appropriate hyperbolic initial-boundary value theorem.

Cancellation of the total flux alone proves neither transparency nor finite speed: two separately reflecting Dirichlet walls already have zero power at the cut because $u_t=0$ there. To test support for the matched problem, use smooth data and repeat the calculation on a ball shrinking with radius $T-t$ toward an observation point. The moving sphere contributes $-e$ per unit area as well as $u_t\partial_nu$. Since
\[
 u_t\partial_nu-e=-\tfrac12(u_t-\partial_nu)^2
                    -\tfrac12|\nabla_{\rm tan}u|^2-\tfrac12m^2u^2\leq0,
\]
energy cannot enter the backward cone through its lateral characteristic side. Any paired internal faces cancel \emph{within that cone}, not just after an integral over the entire cut. If the cone contains no source and its initial ball has zero data, its tip sees no disturbance. For $m=0$ the energy controls derivatives and the initial value fixes the constant. The argument gives the spacetime support statement used for the planar example in Chapter~\ref{ch:08}. Sources or interactions assigning data instantaneously between separated parts of a face would not pass this local argument.

\section{Which analytic input belongs to which cut?}
At an actual Cauchy surface the wave equation propagates a pair of initial data, typically $H^1\times L^2$ at energy regularity. On a timelike wall those two traces generally cannot be prescribed independently. One admissible boundary history produces the other as a response. In a spatial Fourier mode growing into a half-space one can see why an arbitrary sideways prescription does not even give a stable estimate. For the massless wave equation a time-independent tangential mode $e^{iky}$ obeys $\partial_s^2 u=k^2u$ in a flat collar. The continuation $e^{|k|s}$ has exponentially large amplification at fixed $s>0$ as $|k|$ increases. Noncharacteristicity of $ds$ did not turn $s$ into a good time variable.

On a characteristic surface even counting the data as a Cauchy pair is usually wrong. In a double-null rectangle, Equation~\eqref{eq:double-null} provides existence as well as a support estimate using one trace on each incoming side. The trace on a single characteristic Cauchy surface with suitable causal/global conditions is a distinct problem; see B\"ar and Tagne Wafo~\cite{BW}. For gauge fields add propagation of constraints, compatible gauge-parameter boundary conditions, and the definition of proper gauge at exterior faces to the list of analytic inputs. Existence of the geometric quotient supplies none of these estimates.

\begin{exercise}[A retarded inverse on the wrong space]
Let $y$ be a smooth compactly supported function of unit integral on the time line. Can there be a compactly supported $q$ with $2\dot q=y$? Does this rule out a retarded solution?
\end{exercise}
\noindent\emph{Answer.} The time integral of a derivative of a compactly supported function is zero. The retarded solution nevertheless exists: it is zero sufficiently early and tends to $1/2$ sufficiently late. Causal support is not compact support.
\chapter{Hadamard singularities and distributional tests}\label{app:C}
At separated spacelike points the free vacuum two-point function can look
like an ordinary function of distance. At null separation its prescription
matters, and at coincidence it cannot be evaluated at a point. These are
different issues. The Hadamard form specifies a local singularity together
with a frequency orientation. The wavefront set tests whether one can
restrict or multiply the resulting distribution. Scaling degree answers a
third question: if an expression is already defined off coincidence, what
freedom remains in extending it? We give the tests for a real scalar in
the smooth bulk, where all three questions already arise.

\section{World function, frequency and normalization}
Work first in a smooth geodesically convex bulk neighborhood in four
spacetime dimensions with signature $(-,+,+,+)$. For the unique geodesic
$\gamma$ from $x'$ to $x$, affinely parametrized on $[0,1]$, Synge's world
function is
\[
 \sigma_g(x,x')=\tfrac12\int_0^1 g(\dot\gamma,\dot\gamma)\,ds .
\]
Thus it is \emph{half} the signed squared interval, negative for a timelike
separation. It is distinct from the antisymmetric form on CCR labels. Let
$T$ be a smooth time function increasing to the future, in units of length.
Use the boundary value, for $\epsilon\downarrow0$,
\[
 \sigma_\epsilon=\sigma_g+i\epsilon(T(x)-T(x'))+\epsilon^2/2 .
\]
For $T=t$ in flat space this is exactly
$\big(|\mathbf x-\mathbf x'|^2-(t-t'-i\epsilon)^2\big)/2$. Use the
branch of the logarithm real on the positive real axis. With $\hbar=1$ and
$\ell>0$ a length, the local singular part of the \emph{unordered}
(Wightman) two-point function is written
\begin{equation}
 H_\ell(x,x')=\frac{1}{8\pi^2}\lim_{\epsilon\downarrow0}
  \left[\frac{U(x,x')}{\sigma_\epsilon(x,x')}
    +V(x,x')\log\!\frac{\sigma_\epsilon(x,x')}{\ell^2}\right].
 \label{eq:C-hadamard}
\end{equation}
$U,V$ are geometric coefficients for the chosen Klein--Gordon operator,
not functions to be fitted to a state. For example $U$ obeys
\[
 2\nabla^a\sigma_g\nabla_a U+(\Box\sigma_g-4)U=0,\qquad U(x,x)=1 .
\]
To see the first equation, act with $\Box$ on $U/\sigma_g$ away from the
cone and use $\nabla_a\sigma_g\nabla^a\sigma_g=2\sigma_g$. The
coefficient of $\sigma_g^{-2}$ must vanish. The remaining inverse power
determines the leading logarithmic coefficient by transport; higher
coefficients can be constructed successively. On a smooth, not necessarily
analytic, metric this notation denotes a local parametrix modulo smooth
terms. It can be used to any specified differentiability order by
truncating its asymptotic expansion sufficiently late; convergence of an
infinite power series is not assumed \cite{WaldQFT,HWick}.

The flat massless benchmark is $U=1,V=0$. It can also be calculated from
positive-frequency modes: with $\tau=t-t'$, $r=|\mathbf x-\mathbf x'|$,
\begin{align*}
 W_0(x,x')&=\lim_{\epsilon\downarrow0}
 \int\frac{d^3p}{(2\pi)^3 2|\mathbf p|}
 e^{-i|\mathbf p|\tau+i\mathbf p\cdot(\mathbf x-\mathbf x')-\epsilon|\mathbf p|}\\
 &=\lim_{\epsilon\downarrow0}\frac{1}{4\pi^2 r}
       \int_0^\infty e^{-(\epsilon+i\tau)p}\sin(pr)\,dp
 =\lim_{\epsilon\downarrow0}\frac{1}{4\pi^2[r^2-(\tau-i\epsilon)^2]} .
\end{align*}
This checks both the $1/(8\pi^2)$ in \eqref{eq:C-hadamard} and the sign of
the prescription. Reversing the arguments reverses the frequency
orientation. Time ordering instead selects $W_0(x,x')$ for $\tau>0$ and
$W_0(x',x)$ for $\tau<0$; its flat contraction is
$G_F=(4\pi^2)^{-1}(r^2-\tau^2+i0)^{-1}$, as used in
Chapter~\ref{ch:13}. The Wightman function is a bisolution, whereas this
$G_F$ has an $i\delta$ source. Other spacetime dimensions have different
Hadamard powers or logarithms, not the four-dimensional prefactor printed
here.

A scalar Hadamard state has $W-H_\ell$ smooth locally (or the corresponding
finite regularity after truncation). Changing $\ell$ shifts the smooth
representative and finite local subtractions, not the leading singularity.
The parametrix is not in general an exact positive bisolution.
To specify an actual centered quasifree state the full $W$ must satisfy
the equation and the true boundary condition in both entries,
$W(x,x')-W(x',x)=iE(x,x')$, Hermiticity and positivity as in
Chapter~\ref{ch:11}. Two Hadamard two-point functions for the same ordinary
globally hyperbolic free problem differ by a smooth symmetric bisolution.
Their fluctuation predictions can therefore differ even if every
short-distance subtraction is the same.

\section{Which directions are singular?}
A distribution $u$ is smooth near a point if its localized Fourier
transform decays faster than any inverse power in every momentum
direction. The \emph{wavefront set} $\operatorname{WF}(u)$ records pairs
$(x,k)$ with $k\ne0$ where this test fails in direction $k$: no cutoff
nonzero at $x$ gives rapid decay in an open conical neighborhood of $k$.
We take Fourier phases $e^{-ik\cdot x}$ for this test. A positive-frequency
mode has oscillatory covector $k=(-|\mathbf p|,\mathbf p)$, so its
metric dual $k^\sharp$ is future directed. For a state of the ordinary
free scalar on a smooth globally hyperbolic spacetime the Hadamard
condition is equivalently the microlocal spectrum condition \cite{Rad}:
\begin{equation}
 \operatorname{WF}(W)=\{(x,k;x',-k'):(x,k)\sim(x',k'),\quad
                              k^\sharp\ \text{future null}\} .
 \label{eq:C-wavefront}
\end{equation}
Here $\sim$ means that $k,k'$ are cotangents to the same null geodesic
(with metric dual tangent to it) related by parallel transport. The
coincident pairs $x=x',\ k=k'$ on a null direction are included. The minus
on the second covector comes from the two arguments, not a past-directed
first frequency. A mass does not move these singular directions off the
light cone: the principal part of the operator is unchanged. At a genuine
timelike boundary reflected singularities may need to be included in a
global criterion; \eqref{eq:C-wavefront} is the standard smooth-bulk
result, not an assertion about arbitrary characteristic or corner problems.

In the ordinary bulk problem this gives a short proof of the smoothness of
state differences. Let $D=W_1-W_2$ for two Hadamard states with the same
commutator. The wavefront set of the difference is contained in the union
of their wavefront sets, so its first covector must be future null. But
$D(x,x')=D(x',x)$ by the commutator. Exchanging the arguments in
\eqref{eq:C-wavefront} would make that same first covector past null.
There are no such nonzero covectors; hence $D$ is smooth. The proof shows
that smoothness fixes neither the size nor the sign of a difference in
fluctuations.

\section{Test before restricting or multiplying}
Let $\iota:S\to M$ be a smooth embedding. Its conormal consists of the
nonzero covectors at $\iota(s)$ which annihilate every tangent to the
image. The standard pullback $\iota^*u$ is defined by the wavefront
criterion if
\[
 \operatorname{WF}(u)\ \cap\ N^*\iota=\varnothing .
\]
For a kernel of two variables apply this criterion to the actual map, for
example $\iota\times\iota:S^2\to M^2$. Both cotangents must annihilate
their respective tangent spaces to give an obstruction. On a spacelike
Cauchy surface the conormals are timelike, whereas both covectors in
\eqref{eq:C-wavefront} are null. The restriction of $W$ to two Cauchy
arguments is therefore a distribution on the product of the surface with
itself. Spacetime derivatives of $W$ do not enlarge its wavefront set, so
normal derivatives can be taken before this pullback. This is how
field--momentum as well as field--field covariances can be given by Cauchy
data. It does not define the value $W(x,x)$.

Now take the null plane $u=t-x^1=0$ and two points on the same generator,
with the same transverse coordinates. A future null covector at the first
point is $-a\,du$, $a>0$; parallel transport gives the same covector at
the second point. The pair $(-a\,du,+a\,du)$ lies both in the wavefront
set of $W$ and in the conormal to the product of null planes. The usual
pullback theorem no longer allows us to restrict both fields in this
manner. It does not prove that no characteristic observable whatsoever
exists. Smeared fluxes, specially differentiated data or other trace
constructions require their own domains and checks. The classical
existence of a smooth trace for a particular solution is not the quantum
distributional check.

Multiplication is pullback to the diagonal of a tensor product. A
sufficient condition for $uv$ is that there be no $(x,k)\in
\operatorname{WF}(u)$ with $(x,-k)\in\operatorname{WF}(v)$. Smooth
switching functions always pass this test. For a sharp step across a smooth
face $s=0$, $ds\ne0$, the wavefront set of $\Theta(s)$ contains both
normal directions $a\,ds$, $a\ne0$. Hence, for example,
$\Theta(s)\delta(s)$ fails this test: assigning a fraction of a surface
contact to one side by writing $\Theta(0)=1/2$ is not a distributional
proof. The problem can be seen by smoothing the step and shifting its
transition by a distance of the order of its width. Its product with the
delta depends on the value selected at zero even when all these smooth
steps have the same limit. Conversely, a step in only the first argument
of a bidistribution has zero cotangent in the other slot in this product
test. One must check the actual distributions, not declare every time
ordering or every step inadmissible.

This criterion also explains the puncture in the fish calculation. Two
factors of $W$ can multiply as distributions in the \emph{pair} of
variables: their nonzero first covectors are both future directed. Away
from the diagonal $G_F$ selects $W$ or its transpose locally near each
singular pair, so its square there likewise passes the test. But on the
diagonal $\Box_x G_F=i\delta(x,y)$ in the flat example. A differential
operator cannot add directions to a wavefront set, and the delta has all
diagonal conormals $(k,-k)$ there, with both signs of $k$. The square of
$G_F$ at coincidence therefore fails the product test. Extending an
already defined off-diagonal distribution, not assigning a value to a
step at zero, is the operation that remains.

\section{How much freedom does extension leave?}
Suppose a scalar distribution $t^\circ$ is given for $z\ne0$ in
$\mathbb R^k$. For $\psi$ a test function set
$\psi_a(z)=a^{-k}\psi(z/a)$, $t_a(\psi)=t^\circ(\psi_a)$ on the
punctured space. The scaling degree is
\[
 \operatorname{sd}(t^\circ)=
 \inf\{d:a^d t_a\longrightarrow0\ \text{as }a\downarrow0
                     \text{ in distributions}\} .
\]
A homogeneous singularity $|z|^{-b}$ has scaling degree $b$ off zero;
logarithms do not change the infimum. The bulk extension theorem for finite
scaling degree $\rho$ says that an extension of degree $\rho$ exists and
is unique at that degree if $\rho<k$ \cite{BF}. If $\rho\ge k$, two such
extensions may differ by
\[
 \sum_{|\alpha|\le\lfloor\rho-k\rfloor} c_\alpha\,
                                \partial_z^\alpha\delta^{(k)}(z) .
\]
Indeed their difference is supported at zero; distributions so supported
are finite sums of delta derivatives, and each derivative has degree
$k+|\alpha|$. To see why the bound also suggests a construction, subtract
from the test function its Taylor polynomial through degree
$\lfloor\rho-k\rfloor$, multiplied by a smooth cutoff equal to one near
zero. The remainder vanishes fast enough that removing a small neighborhood
of zero has a limit. Assigning the finitely many missing Taylor moments
amounts to choosing the constants $c_\alpha$. An extension is a new
prescription; it need not preserve exact homogeneity, even when it
preserves the scaling degree.

At a smooth collision stratum use local coordinates $(X,z)$ with $z=0$
on the stratum. Scale the $k$ normal coordinates $z$ only, keeping the
tangential variables $X$ as test parameters. The analogous ambiguities
are normal delta derivatives with coefficients on the stratum; in general
those coefficients are distributions subject to further requirements.
Local covariance, allowed field dimensions, unitarity, causal
factorization and Ward identities reduce the freedom in an actual field
theory. This counting works only once the distributions away from the
stratum and the requisite normal estimates are known. A null pullback that
failed the wavefront test is not automatically defined by quoting the
codimension of a coincident-point stratum. Chapter~\ref{ch:13} instead
starts with a defined product away from coincidence and extends it, which
is precisely the situation addressed by the theorem.

\begin{exercise}
For the flat massless vacuum restrict the two arguments to $t=t'=0$.
Starting from the momentum integral above, find the singularity as the
spatial separation $r$ tends to zero. Does the distribution on the pair of
Cauchy surfaces need an arbitrary spatial contact term at $r=0$ if its
scaling degree is to be preserved? Does this give a finite value for the
field squared at a point?
\end{exercise}
\noindent\emph{Hint.} The equal-time pullback passed the conormal test.
Compare $1/r^2$ with the radial volume $r^2 dr\,d\Omega$ in three relative
spatial dimensions; distinguish a locally integrable distribution from its
pointwise value at the origin.

% !TeX root = ../main.tex
\chapter{Geometry from the ground up: building a seam}\label{app:D}\markboth{Appendix D}{}
Two strips of paper come with the instruction ``identify these edges.'' What does this instruction give us? At first, only a rule specifying which points are to count as one. It has not told us how to differentiate across the join, how to compare two vectors attached to the joined point, or how to transport the value of a field around the completed strip. We will construct each of these operations. At the end two flat strips will form a spatial circle. Different, locally indistinguishable ways of identifying their field values will give different allowed wave frequencies.

This appendix can be read before Chapter~\ref{ch:01}. It assumes sets and functions, elementary calculus of several real variables, and linear algebra, but not a course in topology or bundles. A useful question to keep asking is: \emph{what is the domain and what is the codomain of the map I am using?} A map of point sets, its derivative, a map of fibers over it, and the pullback of a differential form go between different spaces.

\section{Give the set of points a topology}\label{D-topology}

A \emph{topology} on a set $X$ is a collection of subsets called \emph{open}: it contains $\varnothing$ and $X$, is closed under arbitrary unions, and is closed under finite intersections. A set is \emph{closed} if its complement is open. A neighborhood of $p$ contains an open set containing $p$. In $\mathbb R^d$, $U$ is open when each $p\in U$ is the center of some ball contained in $U$. What we retain in passing to a topological space is the possibility of testing a statement near each point, not the notion of distance that supplied these particular balls.

A \emph{basis} is a collection of open sets such that every open set is a union of members. Euclidean balls with rational centers and positive rational radii suffice: inside any ball about $p$, one can find a smaller ball of this sort that still contains $p$. A space with a countable basis is called \emph{second-countable}. A different condition, \emph{Hausdorff separation}, says that any two distinct points admit disjoint open neighborhoods. Euclidean space has this property: take balls of radii less than half the distance between the points. These two conditions are part of the definition of a manifold in this book; looking Euclidean near each point does not by itself guarantee them globally.

A map $f:X\to Y$ is \emph{continuous} if the inverse image $f^{-1}(V)=\{x:f(x)\in V\}$ of every open $V\subset Y$ is open in $X$. Here $f^{-1}(V)$ makes sense even if $f$ has no inverse function. It suffices to test a basis of $Y$, since inverse images preserve unions. A \emph{homeomorphism} is a bijection continuous in both directions. A \emph{cover} is a family of sets whose union is the space. A space is \emph{connected} when two nonempty disjoint open sets cannot cover it. In the discrete topology every subset is open. A continuous function from a connected strip to the discrete set $\{+1,-1\}$ is therefore constant: the two inverse images of the singletons would otherwise separate the strip.

We need three elementary ways of making spaces. A subset $A\subset X$ has the \emph{subspace topology}: its open sets are $A\cap U$ with $U$ open in $X$. Thus $[0,a)$, for $0<a<L$, is open \emph{in} $[0,L]$. It is not open in $\mathbb R$. The \emph{product topology} on $X\times Y$ has basis $U\times V$ for $U,V$ open. The \emph{disjoint union} $X=R_-\sqcup R_+$ keeps a tag on each point, even if both pieces use the same numerical coordinates; a set in it is open precisely when its part on each piece is open. We use the product to thicken a face, and the disjoint union to start with two genuinely independent regions.

\subsection{How can we tell what is open after identification?}
An \emph{equivalence relation} is reflexive ($p\sim p$), symmetric ($p\sim r$ implies $r\sim p$), and transitive ($p\sim r$, $r\sim z$ imply $p\sim z$). For a list of face pairings, take the relation they \emph{generate}: a point is equivalent to every point reachable by finitely many pairings and their inverses. Its classes form the set $M=X/{\sim}$. Write $q:X\to M$ for the map $p\mapsto[p]$. To specify the \emph{quotient topology}, declare $U\subset M$ open if and only if $q^{-1}(U)$ is open in $X$. This is a topology: take inverse images in the axioms above. The \emph{saturation} of $A\subset X$ is $\operatorname{Sat}(A)=q^{-1}(q(A))$, the union of all classes it meets. The set $A$ is \emph{saturated} when $\operatorname{Sat}(A)=A$. Saturated open sets upstairs are exactly inverse images of open sets downstairs. The image of an arbitrary open set need not be open.

\begin{proposition}[The test for descending a continuous map]\label{D-quotient-test}
Suppose a map $f:X\to Y$ has the same value at equivalent points. There is a unique map $\bar f:M\to Y$ with $f=\bar f\circ q$. It is continuous exactly when $f$ is continuous.
\end{proposition}
\begin{proof}
Define $\bar f([p])=f(p)$. Equality on a class is exactly what makes the formula single valued, and every class has a representative. For an open set $V\subset Y$,
\[
 q^{-1}\big(\bar f^{-1}(V)\big)=f^{-1}(V).
\]
By the quotient topology the set inside $q^{-1}$ is open exactly when the right side is. Applying this to all $V$ proves the test.
\end{proof}

For example, identify $0$ with $L$ in $[0,L]$. A small open arc around the new point has inverse image $[0,a)\cup(L-a,L]$ for $0<a<L/2$. Keeping only the first interval fails: its saturation must contain $L$ but then contains no relative neighborhood of $L$. Figure~\ref{D-arc-figure} shows the missing half of the neighborhood.

\begin{figure}[htbp]\centering
\begin{tikzpicture}[x=.94cm,y=.94cm,>=Stealth,font=\small]
\draw[quiet] (0,0)--(3.9,0);\draw[line width=2pt,marine] (0,0)--(.80,0) (3.10,0)--(3.90,0);
\fill[marine] (0,0) circle (2pt) (3.90,0) circle (2pt);
\draw[marine,fill=white] (.80,0) circle (2pt) (3.10,0) circle (2pt);
\node[below] at (0,-.07) {$0$};\node[below] at (.80,-.07) {$a$};\node[below] at (3.10,-.07) {$L-a$};\node[below] at (3.90,-.07) {$L$};
\draw[->,ink] (4.5,0)--node[above] {$q$} (5.65,0);
\begin{scope}[shift={(7.1,0)}]
\draw[quiet] (0,0) circle (.83);\draw[marine,line width=2pt] (142:.83) arc (142:218:.83);
\draw[marine,fill=white] (142:.83) circle (2pt) (218:.83) circle (2pt);
\fill[marine] (180:.83) circle (2pt);\node[above left,inner sep=1pt] at (-.91,.60) {$[0]=[L]$};
\draw[thin,quiet] (-1.03,.52)--(-.85,.10);
\end{scope}
\end{tikzpicture}
\caption{One open arc downstairs has two relatively open end intervals upstairs. Filled points are included; the two filled points on the interval represent the single marked point on the circle.}\label{D-arc-figure}
\end{figure}

To identify the quotient with the unit circle, use the map
\[
 c:[0,L]\longrightarrow S^1\subset\mathbb R^2,\qquad
 x\longmapsto\big(\cos(2\pi x/L),\sin(2\pi x/L)\big).
\]
Its values agree at the two ends, and no other distinct pair has the same image. It induces a continuous bijection onto the unit circle in $\mathbb R^2$. Why is its inverse continuous? A space is \emph{compact} if every cover by open sets has a finite subcover. The closed bounded interval is compact (the Heine--Borel theorem); its continuous image is compact since one can pull a cover back. A compact subset of a Hausdorff space is closed: separate an outside point from each point of the subset by disjoint neighborhoods, keep finitely many neighborhoods covering the compact subset, and intersect their partners around the outside point. A closed subset of a compact space is compact (add its open complement to a cover). Our continuous bijection thus sends closed sets to closed sets, which is the continuity test for its inverse. This proof uses the topology on the quotient, not how round a circle can be drawn.

\begin{exercise}[Where separation can fail]
Take two copies of the quadrant $(r,y)\in[0,\infty)^2$. On the faces $r=0$ pair equal $y$ only when $y>0$, leaving the two origins distinct. Can a neighborhood of one origin in the quotient be disjoint from a neighborhood of the other?
\end{exercise}
\noindent\emph{Solution.} Pull both neighborhoods back. Each contains points $(0,y)$ for all sufficiently small positive $y$ on its own copy. Such a pair represents the same point downstairs, so the two neighborhoods must meet. Closing the equivalence relation by also pairing the origins changes the space; it is not a choice of coordinates on the original quotient.

\section{Turn two regional charts into one smooth chart}\label{D-collars}
A \emph{chart} is a homeomorphism from an open neighborhood of a point onto an open subset of a coordinate model. Use $\mathbb R^d$ for an interior and $[0,\infty)^k\times\mathbb R^{d-k}$ at a corner with $k$ meeting walls. Here ``open'' means relatively open in the model. In Euclidean open sets, \emph{smooth} means having continuous partial derivatives of all orders. A function on the model is \emph{smooth up to its walls} if locally it is the restriction of such a function on an open set of $\mathbb R^d$. Smooth maps are defined componentwise this way. A change of charts is the map $\psi\circ\varphi^{-1}$ on the overlapping coordinate images. A collection of charts covering a Hausdorff, second-countable space, whose changes are smooth with smooth inverse, gives a \emph{smooth atlas}; adjoining all compatible charts does not change what is smooth. A bijection of such spaces smooth with smooth inverse is a \emph{diffeomorphism}. A homeomorphism only knows open sets: $x\mapsto x^3$ on the line, for instance, has continuous inverse, but its inverse is not differentiable at zero.

At a point of \emph{depth} $k$ exactly $k$ independent nonnegative wall coordinates vanish. The connected pieces of a fixed depth and incidence type are its \emph{strata}; a face in our regions is a specified embedded closed boundary hypersurface. \emph{Embedded} means it has the subspace topology and locally is a coordinate slice, not a curve self-crossing in the picture. A wall ending on another wall does not thereby create a third hypersurface: in $[0,\infty)^2\times\mathbb R^{d-2}$ the intersection of the two independent walls has dimension $d-2$. The number of independent restrictions is its \emph{codimension}. These labels distinguish artificial walls that disappear from physical walls that remain.

\subsection{Construct the map and its inverse}
Let $F_\pm\subset R_\pm$ be paired faces, with a specified diffeomorphism $f:F_-\to F_+$. For the moment stay away from corners. A \emph{collar} is a diffeomorphism from $[0,\epsilon)\times F$ onto a neighborhood of $F$ taking $(0,p)$ to $p$. Local widths suffice here. The coordinate $r$ measures inward parameter, not proper distance. We supply collars as data; proving a collar-existence theorem for every choice of further structures is not needed. Make the quotient of $R_-\sqcup R_+$ identifying $p$ with $f(p)$ and \emph{no} points at positive collar parameter.

An example with different coordinates on the two walls is a good test of the inverse map. Suppose the face coordinate on the minus side is $y$ and on the plus side is $z$, with $z(f(p))=2y(p)+1$ (dimensionless chart coordinates). Use collar coordinates $(r,y)$ and $(r,z)$, both with $r\geq0$. On the saturated union of corresponding collar boxes define
\begin{equation}
 \Psi([r,y]_-)=(-r,y),\qquad
 \Psi([r,z]_+)=\big(r,(z-1)/2\big).
 \label{D-signed-chart}
\end{equation}
The brackets mean the class of the indicated regional point. Both formulas give $(0,y)$ at the seam. They yield all $(s,y)$ in a small open rectangle about it, exactly once after identification. For the inverse take the class of $(-s,y)_-$ when $s\leq0$, or of $(s,2y+1)_+$ when $s\geq0$. The forward map is continuous by Proposition~\ref{D-quotient-test}. The inverse is continuous by \emph{finite closed pasting}: if continuous maps on finitely many closed subsets covering a space agree on overlaps, the inverse image of any closed set under their union is a finite union of closed sets. Here the two closed subsets are the half-rectangles \emph{relative to the open rectangle}. Thus we have a chart across the seam. Using only the coordinate $r$ on both sides would have lost which half a point came from.

To use many such charts, explicitly compose their maps on overlaps. If, for example, the new tangential coordinate on the minus side is $Y=3y$ and on the plus side is $Y=3(z-1)/2$, then $Y=3y$ on the whole signed rectangle. If a second proposed chart instead uses $S=s$ for $s\leq0$ and $S=2s$ for $s\geq0$, the transition is a homeomorphism but fails the first derivative test. Each collar separately is perfectly smooth; the two proposed signed charts do not belong to a single smooth atlas. A stated map on the face alone would not have detected the failure.

\subsection{Intersections and global checks}
At $k$ transverse artificial walls meeting $l$ retained exterior walls, make one chart on each of the $2^k$ local sectors in the product model. If $r_i$ measures inward distance in coordinates (not metric length), put $s_i=\sigma_i r_i$ with $\sigma_i=\pm1$; keep the same nonnegative exterior coordinates $\rho_a$. The proposed charts must cover \emph{all} representatives of this neighborhood; identifications must join exactly the common subfaces and preserve all the indicated other coordinates there. Their images fill
\[
 (-\epsilon,\epsilon)^k\times[0,\epsilon)^l\times V,
 \qquad V\subset\mathbb R^{d-k-l}\ \text{open}.
\]
The same inverse-by-sectors proof works. In a plane four quadrants can erase both axes; if the horizontal axis is a physical boundary, two quadrants erase the vertical cut and leave a half-plane. Three half-lines glued at their origins instead form a letter Y: deleting the central point of a small connected neighborhood leaves three components, not the one or two allowed in a one-dimensional manifold with boundary. Its product with tangential coordinates is the branched local model one would create by identifying three ordinary half-spaces along one open face. Counting incident sectors, not just assigning signs to boundary integrals, is essential.

There are two more checks in calling the result a manifold. Here is a way of carrying them out for the finite product constructions of Chapter~\ref{ch:01}. Suppose each of finitely many pairings is a homeomorphism between \emph{closed} faces, with inverse included; suppose the sector test bounds the number of representatives of a class by $b$. To find all points equivalent to one point, follow chains of pairings, removing repeated points from each chain. No more than $b-1$ steps are needed. There are only finitely many such words. Each word is a homeomorphism on a closed domain to a closed range: when composing, restrict the first map to the inverse image of the second domain. Hence the saturation of a closed set $C$ upstairs is closed, being the finite union of the images of $C$ intersected with the word domains. Distinct finite classes have disjoint open neighborhoods upstairs by Hausdorff separation (separate each pair of points and take finite intersections). Within any such neighborhood $U$ of a class,
\[
 X\setminus\operatorname{Sat}(X\setminus U)
\]
is an open \emph{saturated} neighborhood of that class. These neighborhoods descend, proving Hausdorff separation downstairs. For countability, pull a cover by our assembled charts back to the disjoint union upstairs, which has a countable basis. For each basis member contained in one of these pullbacks, choose one such chart. This countable selection still covers: at each representative a basis member fits inside a pullback. Each selected chart itself has a countable Euclidean relative basis. Their union gives a basis on the quotient. Neither check assumed that there was already a whole smooth spacetime.

\begin{exercise}[An external corner survives]
Take two copies of $r\geq0$, $\rho\geq0$, with a common tangential coordinate $y\in\mathbb R$. Pair the entire $r=0$ faces with equal $(\rho,y)$; do not pair any open part of $\rho=0$. Write the chart of the class of $(0,0,0)$ and identify its depth before and after sewing.
\end{exercise}
\noindent\emph{Solution.} The signed variable is $s=-r$ or $s=r$ according to the copy. With $(s,\rho,y)$ the local image is open in $\mathbb R\times[0,\infty)\times\mathbb R$; at $(0,0,0)$ only $\rho$ is a boundary coordinate. A depth-two point on each piece has become a depth-one point. Extending both $r$ and $\rho$ to signed variables would introduce points that no piece supplied.

\section{Tangent vectors and pullbacks}\label{D-tangent-maps}

In a signed collar, $+\partial_s$ points out of the piece $s\leq0$. No path starting on its face and proceeding into that piece has this initial velocity. Yet a smooth field on the closed piece has a derivative in the $+\partial_s$ direction. To see why, consider what a tangent vector does to a function.

A \emph{germ} at $p$ is a smooth function defined near $p$, with two functions regarded as the same if they agree on some neighborhood of $p$. A tangent vector $X$ at $p$ is a real linear operation on such germs, with values in $\mathbb R$, satisfying the product rule
\[
 X(hk)=h(p)X(k)+k(p)X(h).
\]
An operation with this rule is called a \emph{derivation at $p$}. It kills constants, since $X(1)=2X(1)$. In a chart $x=(x^1,\ldots,x^d)$, expand a germ as
\[
 h(x)-h(p)=\sum_i (x^i-x^i(p))h_i(x),\qquad h_i(p)=\partial_i h(p).
\]
To obtain $h_i$, integrate $\partial_i h$ along the segment from $x(p)$ to $x$ with parameter running from zero to one, using a smooth extension in a small coordinate box. At a corner the derivatives at $p$ do not depend on the extension: two extensions agree on the nearby interior, and their derivatives agree at $p$ by continuity. Applying the product rule now gives $X(h)=\sum_i X(x^i)\partial_i h(p)$. Conversely, every choice of the $d$ numbers $X(x^i)$ gives a derivation by this formula. The $d$ coordinate derivations $\partial_i|_p$ therefore form a basis of the \emph{tangent space} $T_p R$. In new coordinates $z^a$ the components of this same vector are $X(z^a)=\sum_i(\partial_i z^a)(p)X(x^i)$.

The space $T_p R$ has dimension $d$ even at a corner. By contrast, let $\gamma(\tau)$ be a smooth curve in $R$ for $0\leq\tau<\eta$, with $\gamma(0)=p$. Each active wall coordinate $r_i$ starts at zero and stays nonnegative, so $(r_i\circ\gamma)'(0)\geq0$. These inequalities describe exactly the cone of possible right velocities: for every tuple satisfying them, the straight ray stays in the chart for a short time. A two-sided curve staying in $R$ has derivative zero in each active wall coordinate at $p$, since each has a local minimum. Neither collection of velocities should replace the full tangent space when differentiating across a cut.

For a smooth map $H:R\to M$, the \emph{differential} $\mathrm{d}H_p:T_p R\to T_{H(p)}M$ is defined by $(\mathrm{d}H_p X)(h)=X(h\circ H)$. Applying this to the target coordinate functions gives the matrix of partial derivatives of $H$, its Jacobian. The \emph{cotangent space} $T_p^*R$ consists of linear maps from $T_p R$ to $\mathbb R$. For instance $\mathrm{d}h_p(X)=X(h)$; the coordinate covectors $\mathrm{d}x^i$ are dual to the coordinate vectors because $\mathrm{d}x^i(\partial_j)=\delta^i_j$. A smooth choice of covector at each point is a one-form. To bring a one-form $\beta$ on $M$ back to $R$, set
\[
 {(H^*\beta)}_p(X)=\beta_{H(p)}(\mathrm{d}H_p X).
\]
Thus vectors go forward and covectors pull back. No inverse of $H$ is required; in particular $H$ can be the inclusion of a face.

For example, write a minus collar in a common two-dimensional chart by the map $\Phi(r,y)=(s,z)=(-r,y+\lambda r)$, where $\lambda$ is constant. It gives
\[
 \mathrm{d}\Phi(\partial_r)=-\partial_s+\lambda\partial_z,
 \qquad \Phi^*\mathrm{d}s=-\mathrm{d}r,\qquad
 \Phi^*\mathrm{d}z=\mathrm{d}y+\lambda\mathrm{d}r .
\]
Consequently a one-form $\alpha=A\,\mathrm{d}s+B\,\mathrm{d}z$ has the regional expression
\[
 \Phi^*\alpha=\bigl(-(A\circ\Phi)+\lambda(B\circ\Phi)\bigr)\,\mathrm{d}r
                  +(B\circ\Phi)\,\mathrm{d}y .
\]
On the face all these maps are $y\mapsto z=y$. Pulling back to the face itself gives only $B(0,y)\,\mathrm{d}y$: a vector tangent to $r=0$ pairs to zero with $\mathrm{d}r$. Thus the face map knows how to compare tangential traces, but it cannot tell us the coefficient of $\mathrm{d}r$ in the collar.

\begin{exercise}[Clean is not the same as transverse]
Two embedded submanifolds $A,B$ meet \emph{cleanly} if $A\cap B$ is a submanifold and $T_p(A\cap B)=T_pA\cap T_pB$ at each point of intersection. They meet \emph{transversely} if $T_pA+T_pB=T_pM$ there. Compare two distinct coordinate planes in $\mathbb R^3$, a plane taken twice, and the curves $y=0$, $y=x^2$ in $\mathbb R^2$ at zero.
\end{exercise}
\noindent\emph{Solution.} Two coordinate planes meet in a coordinate line; their tangent spaces intersect in that line and span $\mathbb R^3$. A plane taken twice meets cleanly but the sum is the plane, not $\mathbb R^3$. The two curves meet in a point, whose tangent space is zero, but both tangent lines are the $x$-axis. Their meeting is neither clean nor transverse. Independent defining covectors at a product corner provide more than a drawing of intersecting faces.

\section{Forms and the oriented flux through a face}\label{D-forms-flux}

A differential $k$-form assigns to each point an alternating, multilinear function of $k$ tangent vectors, with coefficients smooth in a chart. A function is a zero-form. The coordinate $k$-forms $\mathrm{d}x^{i_1}\wedge\cdots\wedge\mathrm{d}x^{i_k}$, $i_1<\cdots<i_k$, form a basis at each point: their values on $X_1,\ldots,X_k$ are the determinants of the matrices with entries $\mathrm{d}x^{i_a}(X_b)$. The \emph{wedge product} concatenates lists of covectors and extends by bilinearity; interchanging two covectors changes the sign, so a repeated covector gives zero. For example,
\[
 (\mathrm{d}s\wedge\mathrm{d}z)(V,W)=V^s W^z-V^z W^s,
 \qquad \iota_V(\mathrm{d}s\wedge\mathrm{d}z)=V^s\mathrm{d}z-V^z\mathrm{d}s .
\]
The second operation, \emph{contraction}, puts $V$ in the first slot and leaves the other slots free. Pullback of a form puts $\mathrm{d}H_p X_i$ in \emph{each} slot, so pullback preserves wedge products.

The \emph{exterior derivative} differentiates coefficients and puts the new covector first. Writing $\mathrm{d}x^I$ for an increasing wedge as above, the rule is $\mathrm{d}(a_I\,\mathrm{d}x^I)=\mathrm{d}a_I\wedge\mathrm{d}x^I$, extended by addition. Thus $\mathrm{d}\alpha=(\partial_s B-\partial_z A)\,\mathrm{d}s\wedge\mathrm{d}z$ for the one-form just computed. Commuting second partial derivatives occur with opposite wedge signs, so $\mathrm{d}^2=0$. The ordinary product rule gives $\mathrm{d}(\omega\wedge\eta)=\mathrm{d}\omega\wedge\eta + {(-1)}^k\omega\wedge\mathrm{d}\eta$ for $\omega$ of degree $k$. On functions the chain rule reads $H^*\mathrm{d}h=\mathrm{d}(h\circ H)$; applying the wedge and product rules to the coordinate expansion then gives $H^*\mathrm{d}\omega=\mathrm{d}(H^*\omega)$ in every degree.

An orientation chooses positive ordered bases of $T_p R$ consistently from point to point. Let $\epsilon$ be a volume form, a top-degree form positive on these bases. Suppose $F$ is given locally by $\rho=0$, with $\rho\leq0$ inside and $\mathrm{d}\rho\ne0$ on $F$. The tangent space $T_p F$ is the kernel of $\mathrm{d}\rho_p$. Covectors vanishing on this kernel form the \emph{conormal line}; a \emph{coorientation} selects one of its two rays continuously. Here $\nu_p=\mathrm{d}\rho_p$ selects the outward ray. It is a covector on the \emph{ambient} space at $p\in F$, not the pullback $j^*\mathrm{d}\rho$ by inclusion $j:F\hookrightarrow R$; the latter vanishes because $\rho\circ j=0$.

There is also a \emph{normal quotient} $T_pR/T_pF$. Two ambient vectors give the same element of this one-dimensional space if their difference is tangent. The conormal line is its dual: a covector that vanishes on all tangents assigns one value to the entire equivalence class. No metric has so far selected a preferred vector representing that class.

Choose an outward transverse vector $N$, meaning $\nu(N)>0$. The \emph{outward-first} boundary orientation makes $e_1,\ldots,e_{d-1}$ positive on $F$ if $\epsilon(N,e_1,\ldots,e_{d-1})>0$. Any other such vector has the form $cN+T$, with $c>0$ and $T$ tangent to $F$. Its tangential part contributes zero by alternation, so the sign does not change. This convention fixes Stokes' formula for a smooth compactly supported $(d-1)$-form $\omega$ up to the faces:
\[
   \int_R\mathrm{d}\omega=\sum_{F\text{ a boundary face}}\int_F j_F^*\omega .
\]
To check the sign, orient the strip $a\leq s\leq b$, $z\in\mathbb R$, by $\mathrm{d}s\wedge\mathrm{d}z$ and use $\omega=\alpha$ with compact support in $z$. Integrating $\partial_s B-\partial_z A$ over the strip gives $\int_{\mathbb R}(B(b,z)-B(a,z))\,\mathrm{d}z$ by the fundamental theorem of calculus. At $b$ the positive boundary direction is $+\partial_z$; at $a$ it is $-\partial_z$. Splitting the integral at an intermediate $s$ gives the upper-end sign on the minus piece and the lower-end sign on the plus piece.

For a face pairing $f:F_-\to F_+$, write $j_\pm$ for the regional inclusions. An oriented sewing requires $f$ to reverse the induced boundary orientations: the same signed transverse direction is outward on one piece and inward on the other. Changing variables in the second face integral gives
\[
 \int_{F_-}j_-^*\omega_-+\int_{F_+}j_+^*\omega_+
   =\int_{F_-}\bigl(j_-^*\omega_- - f^*(j_+^*\omega_+)\bigr).
\]
In particular, restrictions of a single form on the sewn collar cancel there. It is the tangential pullbacks that enter this calculation, not all the coefficients of the regional forms.

Normalize $N$ by $\nu(N)=1$ and put $\epsilon_\nu=j^*\iota_N\epsilon$. For any vector field $V$, the difference $V-\nu(V)N$ is tangent on $F$, so $j^*\iota_V\epsilon=\nu(V)\epsilon_\nu$. If instead one integrates against a positive measure, the \emph{outward flux density} is $\nu(V)|\epsilon_\nu|$. The absolute value means that the measure transforms by the absolute Jacobian, without an orientation sign. Rescaling $\nu$ by a positive function rescales the normalized transverse vector by the inverse factor, up to a tangent vector. Hence the density is unchanged. When adding outward densities one must not insert a second boundary sign (Appendix~\ref{app:A}).

\begin{exercise}[Flux through a moving face]
In the $(t,x)$ plane orient by $\mathrm{d}t\wedge\mathrm{d}x$ and keep $x\leq h(t)$. Find the outward flux of $V$ per unit $|\mathrm{d}t|$ on the curve $j(t)=(t,h(t))$.

\emph{Answer.} Here $\nu=\mathrm{d}x-h'(t)\mathrm{d}t$ and $N=\partial_x$ has $\nu(N)=1$, giving $\epsilon_\nu=-\mathrm{d}t$. Direct pullback gives $j^*(V^t\mathrm{d}x-V^x\mathrm{d}t)=(h'V^t-V^x)\mathrm{d}t$, with coefficients evaluated on the curve. Relative to the boundary orientation $-\mathrm{d}t$, the outward coefficient is therefore $V^x-h'V^t$. In particular a vector moving with the face has zero flux even if $V^x\ne0$.
\end{exercise}

\section{Lorentzian cones and Cauchy surfaces}\label{D-causal-geometry}

A Lorentzian metric $g$ assigns to each tangent space a symmetric bilinear form varying smoothly with the point. In a suitable basis its matrix is $\operatorname{diag}(-1,1,\ldots,1)$. It is therefore nondegenerate: the only vector pairing to zero with every vector is zero. A nonzero $V$ is timelike, null, or spacelike according as $g(V,V)$ is negative, zero, or positive. Timelike and null vectors are called causal. The timelike vectors form two cones; a \emph{time orientation} selects one continuously as future. Given a future timelike vector $T$, a causal $V$ belongs to the future cone or its null boundary exactly when $g(T,V)<0$. Choosing future is independent of orienting the spacetime for integration.

To compare two metrics one must give them the same vectors as arguments: ${(H^*g)}_p(X,Y)=g_{H(p)}(\mathrm{d}H_p X,\mathrm{d}H_p Y)$. For the collar map $\Phi$ above, $\mathrm{d}s^2-\mathrm{d}z^2$ pulls back to $\mathrm{d}r^2-{(\mathrm{d}y+\lambda\mathrm{d}r)}^2$; the squares here are tensor squares, not wedges. For every $\lambda$ the induced metric at $r=0$ is $-\mathrm{d}y^2$, but its value on tangential vectors does not determine the mixed coefficient in the collar. For a smooth metric on the chosen sewing, compare all coefficients and their one-sided derivatives of every order in the \emph{same signed chart}. The two choices of future must also agree there.

The metric identifies vectors with covectors by $V^\flat(W)=g(V,W)$. Its inverse sends $\nu$ to the vector $\nu^\sharp$ determined by $g(\nu^\sharp,W)=\nu(W)$; in coordinates ${(\nu^\sharp)}^\mu=\sum_\lambda g^{\mu\lambda}\nu_\lambda$, where $(g^{\mu\lambda})$ is the inverse metric matrix. For a hypersurface the orthogonal complement of $T_p F$ is thus the line spanned by its raised conormal. Write $g^{-1}(\nu,\nu)=\nu(\nu^\sharp)$. If this is negative, an orthonormal basis starting with the normalized timelike normal leaves only positive directions in $T_p F$: the hypersurface is \emph{spacelike}. If it is positive, the normal removes a positive direction, leaving a timelike direction in $T_p F$: the hypersurface is \emph{timelike}. If it is zero, $\nu^\sharp$ belongs to $\ker\nu=T_p F$ and is orthogonal to all of $T_p F$. The induced metric is then degenerate and the hypersurface is \emph{null}. Its metric normal cannot serve as the transverse vector used in Stokes' formula.

For example take the Minkowski plane with $g=-\mathrm{d}t^2+\mathrm{d}x^2$, future $\partial_t$, and orientation $\mathrm{d}t\wedge\mathrm{d}x$. Set $u=(t-x)/\sqrt{2}$ and $v=(t+x)/\sqrt{2}$. Substituting $\mathrm{d}t=(\mathrm{d}u+\mathrm{d}v)/\sqrt{2}$ and $\mathrm{d}x=(\mathrm{d}v-\mathrm{d}u)/\sqrt{2}$ gives
\[
 g=-\mathrm{d}u\otimes\mathrm{d}v-\mathrm{d}v\otimes\mathrm{d}u
     =-2\,\mathrm{d}u\,\mathrm{d}v,\qquad
 \mathrm{d}t\wedge\mathrm{d}x=\mathrm{d}u\wedge\mathrm{d}v .
\]
Keep $u\leq0$. The inverse metric has $g^{uv}=g^{vu}=-1$ and zero diagonal, so the outward conormal $\mathrm{d}u$ raises to $-\partial_v$, tangent to the cut. Instead $N=\partial_u$ gives the positive boundary form $\mathrm{d}v$. A unit time translation $V=\partial_t$ carries flux $\mathrm{d}u(V)=1/\sqrt{2}$ per unit $|\mathrm{d}v|$. Indeed $\iota_{\partial_t}(\mathrm{d}t\wedge\mathrm{d}x)=\mathrm{d}x$, which pulls back to $\mathrm{d}v/\sqrt{2}$ there.

A causal curve has a consistent time direction: orient it toward the future and its tangent must be future causal. Piecewise smooth curves suffice for the following examples; more generally one can take locally Lipschitz curves with future causal tangent almost everywhere. The \emph{causal future} $J^+(p)$ consists of $p$ and the points reachable from it by such curves. In the null chart $g(V,V)=-2V^u V^v\leq0$ and $g(\partial_t,V)=-(V^u+V^v)/\sqrt{2}<0$ together require $V^u,V^v\geq0$, not both zero. Integrating along the curve, both coordinates are nondecreasing. Conversely a straight segment reaches any point $u,v\geq0$ from the origin, so this entire quadrant is $J^+(0)$ in the full plane.

A spacelike hypersurface is \emph{Cauchy} if every inextendible causal curve meets it exactly once. Inextendible means that the curve cannot be continued within the spacetime at either end, even after reparametrization. In the full plane a future causal curve has increasing $t$ and $|\mathrm{d}x/\mathrm{d}t|\leq1$. If its $t$ range had a finite upper endpoint, the bound would give a limit for $x$ there; the curve could continue along a vertical timelike line. The same argument holds at the lower end, so every line $t=\text{constant}$ is Cauchy. The cut $u=0$ misses all the inextendible null lines $u=u_0\ne0$.

\begin{exercise}[A spacelike curve that misses a light ray]
Consider the entire graph $t=\sqrt{1+x^2}$ in the Minkowski plane. Is it spacelike? Is it Cauchy?

\emph{Answer.} Its tangent for increasing $x$ is $Y=(x/\sqrt{1+x^2})\partial_t+\partial_x$, so $g(Y,Y)=1/(1+x^2)>0$ everywhere. Yet the inextendible null line $t=x$ never meets the graph: $\sqrt{1+x^2}>x$ for all real $x$. Spacelikeness was a test at each point of the graph; the Cauchy condition also tests curves elsewhere in spacetime.
\end{exercise}

For a time-oriented spacetime without boundary, existence of a spacelike Cauchy hypersurface characterizes \emph{global hyperbolicity}. It is not a matching condition on a collar. For example, identifying the time caps $(0,x)\sim (L,x)$ of a flat slab with $L>0$ sews both the metric and the future direction smoothly but closes every vertical timelike line. No surface meets such a repeatedly traversed line exactly once. Even on the globally hyperbolic plane, a single null trace need not determine a wave: both $\phi=0$ and $\phi=u$ solve $-2\partial_u\partial_v\phi=0$ and vanish on $u=0$. A well-posed evolution requires existence, uniqueness, and continuous dependence in the chosen data and solution norms. These tests concern a specified equation and admissible data, including conditions at any surviving physical wall; smoothness of the sewn metric alone does not imply them (Appendix~\ref{app:B}).

\section{Building a bundle from local products}\label{sec:D-bundles}

A tangent vector at $x$ comes with a base point before it comes with any components. A charged field value also belongs to a space over $x$; reporting it as a complex number requires a choice of phase convention there. To compare two reports at a cut, we need to know how each convention names the elements of the fiber.

A smooth \emph{fiber bundle} consists of a total space $E$, a smooth surjection $\pi:E\to M$ (the projection), and local product descriptions with a common model fiber $Y$. More precisely, there is an open cover $\{U_i\}$ of $M$ and there are diffeomorphisms
\[
 \tau_i:\pi^{-1}(U_i)\longrightarrow U_i\times Y,
 \qquad \operatorname{pr}_1(\tau_i(v))=\pi(v) .
\]
Such a map is a \emph{local trivialization}: it identifies each actual fiber $E_x=\pi^{-1}(x)$ over $U_i$ with $Y$. It need not extend to a trivialization over all of $M$. A \emph{section} $s:M\to E$ is a smooth choice of one element of each fiber, so $\pi\circ s=\mathrm{id}_M$. In the $i$th trivialization it is described by a function $z_i:U_i\to Y$ through $\tau_i(s(x))=(x,z_i(x))$. The function $z_i$ is not the field $s$: a second report of the very same value generally uses another function.

On $U_i\cap U_j$ write $\tau_i\tau_j^{-1}(x,w)=(x,t_{ij}(x)w)$, where $t_{ij}(x)$ is a diffeomorphism of $Y$ and juxtaposition denotes its action on $w$. The index on the left is the \emph{destination}: $z_i=t_{ij}z_j$. At a point covered by three patches, passing from $k$ to $j$ and then to $i$ is a change of description of just one element. Hence
\begin{equation}
 t_{ii}=1,\qquad t_{ji}=t_{ij}^{-1},\qquad t_{ij}t_{jk}=t_{ik} .
 \label{eq:D-bundle-cocycle}
\end{equation}
The last equality on triple overlaps is the \emph{cocycle condition}. The rightmost map acts first.

Conversely, start with the products $U_i\times Y$ and smooth transition maps satisfying \eqref{eq:D-bundle-cocycle}. Form their disjoint union, identify $(x,w)_j$ with $(x,t_{ij}(x)w)_i$ on each overlap, and use the quotient topology. Project the class $[x,w]_j$ to $x$. Over $U_i$ there is an explicit chart and inverse:
\[
 [x,w]_j\longmapsto (x,t_{ij}(x)w),
 \qquad (x,v)\longmapsto [x,v]_i .
\]
The first map gives the same answer on a representative from $k$ because $t_{ik}t_{kj}w=t_{ij}w$. It is continuous by the quotient topology, and the inverse is continuous by restriction of the quotient map to the $i$th product. Their chart changes are the prescribed smooth transitions. This constructs a local product rather than merely a collection of sets of the same size. To construct a section of it, give smooth functions $z_i$ with $z_i=t_{ij}z_j$ and set $s(x)=[x,z_i(x)]_i$; the transition rule is precisely what removes the choice of $i$.

In a rank-$n$ real \emph{vector bundle}, $Y=\mathbb R^n$ and the transitions are invertible linear maps. One can add vectors in $E_x$ by adding their columns in any trivialization, since a linear transition gives the same sum in the next one. Complex bundles use $\mathbb C^n$ instead. The \emph{tangent bundle} $TM$ is the collection of all $T_xM$, with projection $(x,V)\mapsto x$. Each coordinate basis gives a local product, and its transition matrix is the Jacobian $J=\partial x_i/\partial x_j$: a column of vector components changes by $V_i=J V_j$. The \emph{cotangent bundle} $T^*M$ is built out of the dual spaces. A column of covector components changes by $\eta_i=(J^{-1})^{\mathsf T}\eta_j$, since the pairing $\eta_i^{\mathsf T}V_i$ must give $\eta_j^{\mathsf T}V_j$ for every $V_j$. A \emph{frame} $p:\mathbb R^n\to E_x$ is the linear isomorphism sending the standard basis to a chosen ordered basis of the fiber. Once $p$ is chosen, every other frame at $x$ is $p g=p\circ g$ for a unique $g\in\mathrm{GL}(n,\mathbb R)$; none is preferred.

Collecting all frames gives a \emph{principal bundle} with group $\mathrm{GL}(n,\mathbb R)$. More generally take a Lie group $G$, a manifold whose group multiplication and inversion are smooth. A principal $G$-bundle $P\to M$ carries a smooth \emph{right} action, $(p g)k=p(gk)$, free and transitive on each fiber: for any $p,q$ in that fiber there is exactly one $g$ with $q=p g$. Its local trivializations identify this action with right multiplication on the $G$ factor of $U_i\times G$. The local section $e_i$ given by $e_i(x)=\tau_i^{-1}(x,1)$ lets us write every point over $U_i$ uniquely as $e_i(x)a_i$. If $e_j=e_i g_{ij}$, then $e_i a_i=e_j a_j$ implies $a_i=g_{ij}a_j$, not $a_j g_{ij}$.

Now let $G$ act smoothly on $W$ from the \emph{left}, so $\rho(gk)w=\rho(g)(\rho(k)w)$. Make the \emph{associated bundle} $P\times_G W$ out of pairs $(p,w)$ by declaring
\[
 [p g,w]=[p,\rho(g)w],\qquad [p,w]\longmapsto\pi(p) .
\]
Equivalently, changing $p$ to $p g$ while holding an element fixed changes its component to $\rho(g^{-1})w$. Thus $[e_j,z_j]=[e_i,\rho(g_{ij})z_j]$: the same principal frames give the component changes for every associated field. With the frame bundle and $W=\mathbb R^n$ in the defining representation, the map $[p,w]\mapsto p(w)$ recovers the original vector bundle. Other left actions need not be linear.

\begin{example}[A line with no global frame]\label{ex:D-sign-line}
Regard the unit circle as the complex numbers of modulus one. Cover it by $U=S^1\setminus\{-1\}$ and $V=S^1\setminus\{1\}$. Their intersection has two components. Glue real lines by $t_{VU}=+1$ on the upper semicircle and $t_{VU}=-1$ on the lower semicircle, using the inverse maps on the return. These smooth transitions give a line bundle by the construction above. A nowhere-zero section would have a constant sign in each of the connected patches. Its two signs would have to be equal on the upper overlap and opposite on the lower one. Thus there is no such section and the bundle cannot be a product line, whose constant unit section is nowhere zero. It still has a zero section.

The transitions preserve $|w|$, so it makes sense to keep only the unit vectors. They give a principal $\{+1,-1\}$-bundle: either unit vector in a fiber determines the other by multiplication by $-1$. Lift a circuit starting at $1\in S^1$ with coordinate $+1$ in $U$. In a continuous lift to these two sheets the coordinate cannot change within a patch. Cross to $V$ on the upper overlap, obtaining $+1$, and back to $U$ on the lower one, obtaining $-1$. The returned sign is the holonomy of this discrete bundle. Switching charts and back \emph{at one point} gives $t_{UV}(x)t_{VU}(x)=1$; the circuit switches at two different points. Its minus sign does not violate the cocycle condition.
\end{example}

\section{Lifting a seam map and comparing connections}\label{sec:D-bundle-lift}

Let a diffeomorphism $f:F_-\to F_+$ identify faces of two regions carrying principal bundles $P_\pm\to R_\pm$. The plus-side fiber to be used at $y\in F_-$ lies over $f(y)$. For any smooth $b:N\to M$ the \emph{pullback bundle} implements this change of base:
\[
 b^*E=\{(y,v)\in N\times E:\ b(y)=\pi(v)\},\qquad (y,v)\longmapsto y .
\]
Over $b^{-1}(U_i)$ use $y$ and the second coordinate of $\tau_i(v)$ as a trivialization. Its transitions are $t_{ij}(b(y))$. A section $s$ pulls back to $y\mapsto(y,s(b(y)))$; inclusion of a face gives ordinary restriction. For instance, pull the twisted line of Example~\ref{ex:D-sign-line} back by the double covering $b(e^{i\phi})=e^{2i\phi}$. Start with a unit vector at $\phi=0$ and continue it by following its sheet. When $\phi$ reaches $2\pi$ the image has circled twice: the returned sign is $(-1)^2=1$. The choices close to a global unit section, smooth because its local components are constant. Multiplying this section by an arbitrary real coordinate trivializes the pullback line.

At a cut, pulling back $P_+|_{F_+}$ by $f$ puts two principal bundles over $F_-$, but does not identify them. Even the identity of the circle cannot lift to a linear isomorphism from the product line to the twisted line: the image of its constant unit section would be nowhere zero.

A seam lift is an equivariant bundle isomorphism $D:P_-|_{F_-}\to P_+|_{F_+}$ with $\pi_+D=f\pi_-$ and $D(p a)=D(p)a$. The first equality is read along the two routes in Figure~\ref{D-lift-diagram}.

\begin{figure}[htbp]\centering
\begin{tikzpicture}[>=Stealth,thick,font=\small]
\node (p) at (0,1.4) {$P_-|_{F_-}$};\node (q) at (4.5,1.4) {$P_+|_{F_+}$};
\node (x) at (0,0) {$F_-$};\node (y) at (4.5,0) {$F_+$};
\draw[->,marine] (p)--node[above] {$D$} (q);\draw[->,ink] (x)--node[below] {$f$} (y);
\draw[->,quiet] (p)--node[left] {$\pi_-$} (x);\draw[->,quiet] (q)--node[right] {$\pi_+$} (y);
\end{tikzpicture}
\caption{A lift maps actual fiber elements; the face map below it maps base points. Starting with a frame over $y$ on the minus side gives base point $f(y)$ along either route. Commutativity specifies where the frame must land, not which frame there to choose.}\label{D-lift-diagram}\end{figure}

Equivalently, $p\mapsto(\pi_-(p),D(p))$ is an isomorphism over the identity of $F_-$ onto $f^*(P_+|_{F_+})$. To construct it, choose frames on corresponding patches of the faces and specify smooth $h_i$ on the minus patches by
\begin{equation}
 D(e_i^-(y)a)=e_i^+(f(y))h_i(y)a,\qquad
 h_i g_{ij}^-=(g_{ij}^+\circ f)h_j .
 \label{eq:D-lift-frames}
\end{equation}
Indeed, $D(e_j^-)$ computed as $D(e_i^-g_{ij}^-)$ has coefficient $h_i g_{ij}^-$ in the plus $i$ frame. Starting with $e_j^+$ instead gives the coefficient $(g_{ij}^+\circ f)h_j$. Equality makes the displayed prescription a single map, whose inverse uses $f^{-1}$ and $h_i(y)^{-1}$. Glue $p$ to $D(p)$, not just $\pi_-(p)$ to $f(\pi_-(p))$. Choose product descriptions of the collar bundles as pullbacks of their face bundles, as in Chapter~\ref{ch:02}, and extend the frames and $h_i(y)$ constantly in the collar parameter. A common group coordinate is $h_i(y)a_-$ on the minus side and $a_+$ on the plus side. Equation~\eqref{eq:D-lift-frames} now makes its transition between patches the same on both halves. The two descriptions agree on the identified face; they do not identify base points at positive collar distance. For associated fields the induced map is $[p,w]\mapsto[D(p),w]$. In the common plus frame the two values to compare on the face are therefore $\rho(h_i(y))z_i^-(y)$ and $z_i^+(f(y))$.

Changing local frames to $e_i^\pm k_i^\pm$ describes the \emph{same} $D$ by $(k_i^+\circ f)^{-1}h_i k_i^-$. In contrast, an allowed re-marking acts on the pieces. With the base fixed, let $u_\pm$ be automorphisms of the whole regional bundles respecting the exterior markings, and use hats for their face restrictions. The pair carries $p\sim D(p)$ to $u_-p\sim u_+D(p)$, so the new identification is $D'=\widehat u_+D\widehat u_-^{-1}$. The \emph{stabilizer} consists of the pairs with $D'=D$. This tests the face identification; an isomorphism of smooth sewn configurations must also be smooth in the sewn charts and carry the specified fields to their counterparts.

Cut the circle example at one point of each overlap and allow either sign at each cut. Denote the first crossing sign from $U$ to $V$ by $a$ and the return sign by $b$. A smooth automorphism of the discrete bundle over either complete arc is a single sign, say $u$ or $v$. The crossings become $vau^{-1}$ and $ubv^{-1}$, so their ordered product on return is $u(ba)u^{-1}=ba$. Taking $u=1$, $v=a^{-1}$ reduces the pair to $(1,ba)$. With no fixed exterior frame there are two classes, each with the diagonal stabilizer $u=v$. Independent signs supplied at the two ends of one arc would falsely identify the classes: they need not extend over that arc. The sewings together with their re-markings form a \emph{groupoid}: its objects are the choices $(a,b)$, and its arrows are the pairs $(u,v)$ carrying one choice to another. Arrows compose by composing the regional maps and every arrow has an inverse. The set of two classes remembers possible values of $ba$, but not the nonidentity arrow fixing each object.

To compare values at \emph{different} points of a path in a vector bundle, choose a connection. Its covariant derivative is a linear operator taking sections to bundle-valued one-forms, with the Leibniz rule $\nabla(\lambda z)=\mathrm{d}\lambda\otimes z+\lambda\nabla z$ for a scalar function $\lambda$. In a frame it therefore has the form $\nabla z=\mathrm{d}z+A z$: apply the derivative to the frame vectors to obtain the columns of $A$. On a tangent vector $X$ its column is $X(z)+A(X)z$. For a matrix group in its defining representation the potential $A$ is a matrix of one-forms valued in the Lie algebra (the tangent space to $G$ at the identity). To construct a connection from these local potentials, require that the resulting derivatives represent the same fiber vector in either frame. Write $z_+=h z_-$ on a common base patch. Then
\begin{align*}
 \mathrm{d}(h z_-)+A_+h z_-
 &=h\,\mathrm{d}z_-+(\mathrm{d}h+A_+h)z_-
 =h(\mathrm{d}z_-+A_-z_-)
\end{align*}
for all columns precisely when
\begin{equation}
 A_+=h A_-h^{-1}-(\mathrm{d}h)h^{-1} .\label{eq:D-bundle-connection}
\end{equation}
The inhomogeneous term cancels the derivative of the frame change. The same calculation in fixed frames gives the action on a connection of an active gauge transformation $z\mapsto h z$. On a face the common-base notation on the plus side means $z_+\circ f$ and $f^*A_+$, with the potentials restricted to their faces: $(f^*A_+)_y(X)=A_{+,f(y)}(df_y X)$. This only compares tangential directions. For a comparison in the full signed collar, use the realized frame change and differentiate it in the collar, not just on the face.

To construct parallel transport, take a path $\gamma:[0,1]\to M$ in a frame patch and solve the matrix initial-value problem
\[
 \dot a(t)=-A(\dot\gamma(t))a(t),\qquad a(0)=1 .
\]
The right side is tangent to $G$ at $a(t)$: it is the Lie algebra element $-A(\dot\gamma(t))$ carried from the identity by right multiplication. The curve $e(\gamma(t))a(t)$ is the transported frame, called the \emph{horizontal lift} of $\gamma$ starting at $e(\gamma(0))$. To start at $e(\gamma(0))c$, use $a(t)c$; transport commutes with the principal right action. Thus an associated element $[e(\gamma(0)),w]$ is carried to $[e(\gamma(t)),\rho(a(t))w]$. On a longer path change frames at each crossing and compose the transports, with the later matrix on the left.

In the defining representation each transported column solves $\dot z=-A(\dot\gamma)z$. If the initial-to-final matrix is $U_\gamma$, covariance of this equation gives the new matrix $h(\gamma(1))U_\gamma h(\gamma(0))^{-1}$. The initial frame change is inverted because the new input column must first be put in the old frame. For a loop the two endpoints coincide, giving conjugation, as in Chapter~\ref{ch:09}. The resulting transport around a loop is the \emph{holonomy}. Even the \emph{product} complex line over $S^1$ with one global unit frame can have any $U(1)$ holonomy. Set $A=-i\alpha\,\mathrm{d}\theta$, $\alpha\in\mathbb R$; the angle is local but its differential agrees on overlaps. Parallel transport gives $z(2\pi)=e^{2\pi i\alpha}z(0)$. Changing components by $h=e^{i n\theta}$, $n\in\mathbb Z$, sends $\alpha$ to $\alpha+n$ without changing this phase. Here the nontrivial phase comes from the connection, not from any twist of the bundle.

A small loop also measures \emph{curvature}. In commuting coordinates $s,y$ define $F_{sy}$ by $[\nabla_s,\nabla_y]z=F_{sy}z$. Expanding the two derivatives cancels the terms containing derivatives of $z$ and leaves
\[
 F_{sy}=\partial_s A_y-\partial_y A_s+[A_s,A_y],\qquad [A_s,A_y]=A_sA_y-A_yA_s .
\]
These are the components of the curvature two-form $F=\mathrm{d}A+A\wedge A$, with matrix multiplication in the wedge product. For arbitrary vector fields $X,Y$ subtract $\nabla_{[X,Y]}$ from the commutator to remove the effect of their own failure to commute. The transformation of $\nabla$ implies $F_+=h F_-h^{-1}$ in a common base chart. For a direct check choose $A=-i B s\,\mathrm{d}y$ with constant $B$ in a product complex line. Around the rectangle $0\leq s\leq a$, $0\leq y\leq b$, oriented by $\mathrm{d}s\wedge\mathrm{d}y$, only the side $s=a$ contributes to $\oint A=-iBab$. The transported phase is $e^{-\oint A}=e^{iBab}$. The earlier circle connection had $F=0$ yet could have nontrivial holonomy: a loop encircling this rectangle bounds a surface in its patch, whereas a circuit generating the circle does not.

\begin{exercise}[Pull back before changing frame]
On a seam interval the plus coordinate is $x=2y$ in terms of the minus coordinate. Take
\[
 h(y)=\begin{pmatrix}1&y\\0&1\end{pmatrix},\qquad
 A_-=\begin{pmatrix}0&0\\1&0\end{pmatrix}\mathrm{d}y .
\]
Find $A_+$ in its own coordinate $x$, and check it on the minus-side parallel column $z_-=(1,-y)^{\mathsf T}$.

\emph{Answer.} The right side of \eqref{eq:D-bundle-connection} is $\left(\begin{smallmatrix}y&-1-y^2\\1&-y\end{smallmatrix}\right)\mathrm{d}y$. This is the \emph{pullback} of the desired potential. Since $\mathrm{d}x=2\,\mathrm{d}y$,
\[
 A_+(x)=\tfrac12\begin{pmatrix}x/2&-1-x^2/4\\1&-x/2\end{pmatrix}\mathrm{d}x .
\]
The plus column pulled back is $h(y)z_-(y)=(1-y^2,-y)^{\mathsf T}$. Its $y$ derivative is $(-2y,-1)^{\mathsf T}$, whereas $(f^*A_+)(\partial_y)$ applied to it gives $(2y,1)^{\mathsf T}$. The column is still parallel, although its ordinary derivative has changed.
\end{exercise}

% !TeX root = ../main.tex
\section{A jet is an equivalence class, not a Taylor series}\label{D-jets}
We now have a base manifold, a field bundle $\pi:E\to M$, and a way to compare fibers. How can a first-derivative Lagrangian take a field as input without depending on which coordinates we used to differentiate it? Fix $p\in M$. Consider all local smooth sections defined near $p$. Declare two sections equivalent to order $k$ if, in one common base-and-fiber chart, their values and all component derivatives of total order up to $k$ agree at $p$. The equivalence class of $u$ is its \emph{$k$-jet} $j^k_pu$. Its base point must be the same to compare it with another jet. The family of these classes is the \emph{jet bundle} $J^kE\to M$. For $k=0$ it is just $E$: only the value survives. Forgetting the derivatives of top order maps $J^kE\to J^{k-1}E$ for $k\geq1$. A Lagrangian on $J^1 E$ assigns a density to first jets, not a field history to test against an equation of motion.

Why does this definition not depend on the coordinates? Take a line as base and let two sections have real coefficients $u(x)=x$ and $v(x)=x+x^2$ in a product chart. At the origin their triples (value, first derivative, second derivative) are $(0,1,0)$ and $(0,1,2)$. In particular they have the same first jet but different second jets. Change both coordinates: use $y=2x$ on the base and $Z=\exp(u)$ in the fiber, so that a section with old coefficient $h$ now has coefficient $Z_h(y)=\exp(h(y/2))$. The chain rule gives
\[
 (Z_h)'=\tfrac12 e^h h',\qquad
 (Z_h)''=\tfrac14 e^h\big((h')^2+h''\big),
\]
where the right sides are evaluated at $y/2$. The new triples are $(1,1/2,1/4)$ and $(1,1/2,3/4)$. Neither a derivative's number nor the origin of the fiber coordinates was invariant. \emph{Equality to a specified order} was. In general an order-$j$ derivative of a composite uses only derivatives of the original section through order $j$ and derivatives of the coordinate changes at the corresponding arguments. Equal old jets therefore give equal new jets; applying the inverse changes proves the converse. This supplies smooth transitions on the jet coordinates (with symmetry in repeated differentiations), not a separate tensor transformation for each derivative.

A \emph{germ} remembers more: two sections give the same germ if they agree on some neighborhood of $p$. Even all jets at one point need not determine a smooth germ. The function zero and the function $e^{-1/x^2}$ for $x>0$, extended by zero for $x\leq0$, have the same jets at $0$ but different germs there. On the positive side every derivative is a polynomial in $1/x$ times the exponential and tends to zero. So a gluing proof based on matching derivatives must use the actual given sections, rather than summing a purported convergent Taylor series.

\subsection{Test two sections for smooth descent}
Work in the signed collar of Section~\ref{D-collars}. First rewrite the two local sections in the \emph{same} base coordinates and the \emph{same} fiber chart, using the specified collar realization of the lift. For instance a one-form $a(r,y)\,dr$ on the minus side has coefficient $-a(-s,y)$ in front of $ds$. For a nonconstant fiber change one must differentiate the change as well as the old field. Only after that operation compare ordinary derivatives. If $u_\pm$ denote the resulting scalar components in their inward variables and the common $y$, the test at the seam is
\begin{equation}
 (-1)^m \partial_y^\alpha\partial_{r_-}^m u_-(0,y)
   =\partial_y^\alpha\partial_{r_+}^m u_+(0,y).
 \label{D-field-test}
\end{equation}
Here $\alpha$ is a list of nonnegative integers, one for each tangential coordinate; $\partial_y^\alpha$ means differentiate that many times in each. For a $C^k$ test impose $m+|\alpha|\leq k$ on components already $C^k$ on closed halves; for smoothness impose all orders on smooth halves. The field value on the seam is its order-zero \emph{trace}. Equality of traces as functions of $y$ supplies all purely tangential equalities but says nothing about $\partial_s u$.

Here is a direct sufficiency proof. At an internal seam the two values paste into a continuous function. Paste the two proposed $s$-derivatives into a continuous function $g$ using the first derivative test. On a coordinate interval crossing zero the fundamental theorem of calculus on the two halves gives
\[
 u(h,y)-u(-a,y)=\int_{-a}^{0}g(s,y)\,ds+\int_0^h g(s,y)\,ds
 \qquad(a,h>0).
\]
The same identity holds on subintervals up to zero, so difference quotients give $\partial_s u(0,y)=g(0,y)$. Along tangential lines at the seam differentiate the common trace. All first partials are now continuous, so $u$ is $C^1$. Apply the same argument to the piecewise first partials, then the second partials, up to the required order. At a finite product junction the coordinate interval splits into finitely many sector intervals, and the same integral argument works for \emph{every mixed partial}. At a remaining exterior wall, smoothness in our convention means local extension across that wall: perform the argument on compatible local extensions and restrict back, as in Chapter~\ref{ch:03}. Necessity in each case is simply restriction of an actual smooth section. This is \emph{descent}: regional representatives satisfying the tests are the restrictions of one object, unique since the quotient points are all accounted for. For a gauge arrow apply the test to its group-valued functions in a common frame, not only to the fields on which it acts.

\begin{exercise}[Two different failures at a face and a vertex]
(i) On $s\leq0$ take $u=s$ and on $s\geq0$ take $u=s+bs^3$. Which jets match at $0$? (ii) On four closed quadrants put $u_\sigma=c_\sigma s_1s_2$. All values on both axes are zero. Which derivative at the common vertex detects differing $c_\sigma$?
\end{exercise}
\noindent\emph{Solution.} (i) Value, first, and second derivatives match. The third derivatives are $0$ and $6b$; if $b\ne0$ the field is $C^2$ but not $C^3$. (ii) Each pure derivative at the vertex is zero, but $\partial_{s_1}\partial_{s_2}u_\sigma=c_\sigma$. The full mixed test detects the discrepancy. These pieces need not even pass first-normal-derivative tests away from the vertex along an open axis; comparing only their restrictions to the axes did not test those derivatives.

\begin{exercise}[Can the metrics be glued in the stated charts?]
Use the inward collars $r_\pm\geq0$ and a common time $t$. Give each half the metric $g_\pm=dr_\pm^2-e^{2a_\pm r_\pm}dt^2$, with real constants $a_\pm$. Both induced metrics on the face are $-dt^2$. For which constants is there a smooth metric in the prescribed signed chart?
\end{exercise}
\noindent\emph{Solution.} The signed coefficients of $dt^2$ are $-e^{-2a_-s}$ and $-e^{2a_+s}$. Equality of the first derivatives at zero requires $2a_-=-2a_+$. In this family that condition is sufficient too: the two exponential formulas are then identical, the $ds^2$ coefficients are both one, and both choose $\partial_t$ as future if so marked. A general smooth metric need not be decided by its first jet. Here the special two-parameter family makes the test exact.

\section{Use a conormal to test an equation}\label{D-symbol}
Field equations can reduce the amount of data that must be exchanged. To see what is geometric about this assertion, let a linear second-order operator in coordinates have highest-derivative part $a^{\mu\nu}(x)\partial_\mu\partial_\nu$. Its \emph{principal symbol} at $p$ and covector $\xi$ is $a^{\mu\nu}(p)\xi_\mu\xi_\nu$ (we omit Fourier factors of $i$). For a system this is a linear map from field components to equation components. The second derivative of a change of coordinates only multiplies first derivatives of the unknown, by the chain rule. The highest-order coefficient therefore has exactly the two Jacobian factors canceled by the transformation of the covector slots. For a nonlinear equation one first linearizes at the solution jet under study. A hypersurface is \emph{characteristic} there if its symbol on a nonzero conormal is not invertible, for a square system of this type. Rescaling that conormal by a nonzero number only multiplies its symbol by the square of that number. This definition uses the equation; calling a hypersurface \emph{null} uses the metric.

In units where the metric light speed is one, put a field of speed $c>0$ in a specified inertial frame on the flat $(t,x)$ plane. Its density is $\frac12(u_t^2-c^2u_x^2-m^2u^2)\,|dt\,dx|$. If $c\ne1$ this is intentionally a propagation law distinct from the metric wave equation. Take the moving cut $s=x-vt=0$, $y=t$. The common coordinate field $U(y,s)=u(y,s+vy)$ has $u_t=U_y-vU_s$, $u_x=U_s$, and the absolute coordinate Jacobian is one. Differentiating the transformed density coefficient with respect to $U_s$ gives the cooriented momentum, and changing coordinates in the equation gives
\begin{align}
 p^s&=(v^2-c^2)U_s-vU_y,\nonumber\\
 0&=U_{yy}-2vU_{ys}+(v^2-c^2)U_{ss}+m^2U .\label{D-moving-symbol}
\end{align}
Suppose two actual half-solutions smooth up to the cut have the same trace $q(y)$ and the same $p^s(y)$. For $v^2\ne c^2$, one first obtains $U_s=(p^s+vq')/(v^2-c^2)$ on both sides. In \eqref{D-moving-symbol} the terms $U_{yy}$ and $U_{ys}$ are now obtained by tangential differentiation. Solving for $U_{ss}$ gives a common second normal derivative. Differentiate the equation in $s$ and repeat to reconstruct every higher normal derivative. Section~\ref{D-jets} turns the matched jets of the half-fields into a smooth solution. A noncharacteristic timelike face need not be a surface on which arbitrary imposed data give a well-posed evolution; this argument started with solutions, not arbitrary data.

For the metric $-dt^2+dx^2$, the squared conormal length is $1-v^2$, but the symbol of our operator on $ds$ is $v^2-c^2$. Thus if $c\ne1$, $v=1$ makes the cut metric-null without making it characteristic. At $v=c$ the momentum instead repeats tangential information: $p^s=-c q'$. For $m=0$ the halves $U_-=0$, $U_+=s$ are solutions with the same zero trace and momentum; their first normal derivatives differ. This failure cannot be repaired by a new choice of transverse coordinate. In Chapter~\ref{ch:07} a nonlinear normal momentum replaces this linear formula, and sharing an invertible branch becomes a separate requirement.

The conormal has one further use later in the book. Collect, for every $p\in\Gamma$, the covectors that vanish on $T_p\Gamma$: this gives the conormal bundle $N^*\Gamma\subset T^*M|_\Gamma$. For a single defining function they are just the multiples of $ds|_p$. For a smooth field the pullback to $\Gamma$ is always defined. For a distribution, the question depends on whether its singular directions include a nonzero conormal; Appendix~\ref{app:C} states the precise sufficient pullback test. That is a third use of $ds$, not another definition of characteristicity.

\section{Pull back a two-form before dividing out a direction}\label{D-reduction}
The preceding tangent spaces belonged to spacetime. On a space of fields a tangent vector $Y$ instead differentiates a family $\phi_\epsilon$ of fields at $\epsilon=0$; on a family of solutions it respects the linearized equation and boundary conditions. We use $\delta$ for exterior differentiation in these variables. A \emph{presymplectic form} is a closed two-form $\omega$ ($\delta\omega=0$); at a point its kernel consists of the $Y$ for which $\omega(Y,Z)=0$ for \emph{every} $Z$. A closed nondegenerate two-form is symplectic in finite dimensions. For the covariant potential of a field the space and the variations have to be specified; a symmetry with a boundary charge need not be in the kernel. We can do a whole reduction with ordinary coordinates to see what the assertion about a kernel means.

In the circle Maxwell example of Chapter~\ref{ch:09}, use arc angles $\alpha_1,\alpha_2$ modulo $2\pi$ for the two open transports $e^{i\alpha_i}$, and real electric momenta $e_1,e_2$. Angle differences by $2\pi$ on local charts are constant, so the one-forms $\delta\alpha_i$ agree globally. The unreduced potential is $\theta=e_1\,\delta\alpha_1+e_2\,\delta\alpha_2$. Impose flux matching first, by the inclusion
\[
 j:(\alpha_1,\alpha_2,e)\longmapsto(\alpha_1,e,\alpha_2,e)
\]
with output ordered as $(\alpha_1,e_1,\alpha_2,e_2)$. Restriction of a form is pullback, so
\[
 j^*\omega=j^*\delta\theta=\delta e\wedge(\delta\alpha_1+\delta\alpha_2).
\]
Contract this with $Y=A\partial_{\alpha_1}+B\partial_{\alpha_2}+C\partial_e$. The result is $C(\delta\alpha_1+\delta\alpha_2)-(A+B)\delta e$. It vanishes exactly for $C=0$ and $A+B=0$. The shifts $(\alpha_1,\alpha_2)\mapsto(\alpha_1+\lambda,\alpha_2-\lambda)$ move along precisely these directions. They are the effective relative endpoint re-marking. The map
\[
 (\alpha_1,\alpha_2,e)\longmapsto(\alpha=\alpha_1+\alpha_2\ \mathrm{mod}\ 2\pi,\ e)
\]
has exactly the shift orbits as fibers; every output has a representative $(\alpha,0,e)$. The resulting cylinder carries $\omega_{\rm red}=\delta e\wedge\delta\alpha$, nondegenerate and pulling back to $j^*\omega$. This is \emph{reduction} in this example: restrict, identify the null directions, and exhibit the form on the quotient. With $\delta H(Y)=\omega(Y,X_H)$ one gets $X_e=\partial_\alpha$, hence $\{e^{i\alpha},e\}=i e^{i\alpha}$. The relative angle was redundant; the sum, the phase of the closed loop, is measurable. The calculation would lose it if both open phases had been discarded at the outset.

\section{Build a field on the circle from two strips}\label{D-capstone}
Here is a problem to try without first guessing a field on an uncut circle. Let $R_i=\mathbb R_t\times[0,d_i]_{x_i}$, $d_i>0$, $i=1,2$. Identify $(t,d_1)_1$ with $(t,0)_2$, and $(t,d_2)_2$ with $(t,0)_1$. Each strip has the flat metric $-dt^2+dx_i^2$ and increasing future time. On each take a real line with its fixed fiber norm, so its orthonormal frame group is $O(1)=\{+1,-1\}$. A field is a section of this line bundle. When going from the end of strip 1 to the start of strip 2, multiply field components by $a\in\{+1,-1\}$; at the return seam multiply by $b\in\{+1,-1\}$. Use these constant identifications throughout the respective collars. On each strip take the massive wave density $\tfrac12[(\partial_t u_i)^2-(\partial_{x_i}u_i)^2-m^2u_i^2]$, $m\geq0$. There is no sheet action. From these data construct the space, identify the field bundle, and find the allowed frequencies of real separated modes. Does the answer depend on which strip end carries a minus sign?

\subsection{Solution: build the base, then the field and its derivative}
Put $L=d_1+d_2$. The circular coordinate is $X=x_1$ on strip 1 and $X=d_1+x_2$ on strip 2, read modulo $L$. The induced map is a homeomorphism onto $\mathbb R\times S^1$; one can write the inverse locally even for a time interval of infinite duration. At the first seam take the signed coordinate $s=x_1-d_1$ on strip 1 and $s=x_2$ on strip 2. At the return seam take $s=x_2-d_2$ on strip 2 and $s=x_1$ on strip 1. The inverse uses the old strip for $s\leq0$ and the new one for $s\geq0$. All transition derivatives for the spatial variable are $+1$ in these oriented charts. The local inverses and direct maps are smooth, both metrics become $-dt^2+ds^2$, and $\partial_t$ agrees. The time slices are Cauchy circles, as in the zero-shift case of Chapter~\ref{ch:01}.

At the first seam a common fiber coordinate is the coefficient on strip 2: it equals $a u_1$ on the old side and $u_2$ on the new. On the covering line, changing the strip-2 frame by $a$ makes the first transition $+1$ and the return transition $H=ba$. Thus the pulled-back component obeys $u(t,X+L)=H u(t,X)$. This is a statement about the chosen field bundle, not an extra force at either seam. The product $H$ is unchanged by admissible constant re-markings of either connected strip. A first derivative of the field in the same common frame also carries the sign $a$, not a second, independently chosen transition.

We can derive the derivative matching without first assuming a smooth field. Vary the regional actions under a common value at the first seam, so $\delta u_2(t,0)=a\,\delta u_1(t,d_1)$. Spatial integration by parts leaves
\[
 \int dt\,\big[-\partial_{x_1}u_1(t,d_1)+a\,\partial_{x_2}u_2(t,0)\big]\,\delta u_1(t,d_1).
\]
Time variations may be compactly supported. Since the face variation is free and $a^2=1$, this gives $\partial_{x_2}u_2(t,0)=a\,\partial_{x_1}u_1(t,d_1)$. The same argument at the second seam gives the return sign $b$. It is the stationarity of this specified action that yields the first derivative, while the fiber identification supplied the common field value. The bulk equation $u_{x_i x_i}=u_{tt}+m^2u$ then recursively supplies all higher spatial derivatives, with time derivatives obtained along the seams. The pieces consequently pass the jet test, not merely a point-value test.

\subsection{Solution: carry two numbers around the assembled space}
For a real mode write $u_i(t,x_i)=\cos(\omega t) f_i(x_i)$. Then $-f_i''=k^2 f_i$ with $k^2=\omega^2-m^2$. The symbol $k^2$ is initially just an eigenvalue; why is it nonnegative? Multiply by $f_i$ and integrate by parts on both intervals. The first seam gives $-f_1(d_1)f'_1(d_1)+f_2(0)f'_2(0)=0$, because both factors acquire $a$ and $a^2=1$. The return seam cancels in the same way. Thus
\[
 k^2\sum_i\int_0^{d_i} f_i^2\,dx_i=\sum_i\int_0^{d_i}(f_i')^2\,dx_i .
\]
For a nonzero mode we may take $k\geq0$. For $k>0$, solving the ordinary initial value problem at either strip's start gives
\[
 Y_i(d)=T_k(d)Y_i(0),\qquad Y_i=\binom{f_i}{f_i'},\qquad
 T_k(d)=\begin{pmatrix}\cos(kd)&\sin(kd)/k\\-k\sin(kd)&\cos(kd)\end{pmatrix}.
\]
For example the two columns are the solutions with initial data $(1,0)$ and $(0,1)$, together with their derivatives. The first join gives $Y_2(0)=a T_k(d_1)Y_1(0)$; the second gives $Y_1(0)=b T_k(d_2)Y_2(0)$. Addition of angles gives $T_k(d_2)T_k(d_1)=T_k(L)$. A nonzero initial vector returns consistently if and only if
\[
 0=\det(\mathbf{1}_2-H T_k(L))=1-H\operatorname{tr}T_k(L)+H^2\det T_k(L)
   =2(1-H\cos(kL)).
\]
Thus for $H=+1$ the positive wave numbers are $2\pi n/L$, $n=1,2,\ldots$; for $H=-1$ they are $(2n+1)\pi/L$, $n=0,1,\ldots$. At these roots $H T_k(L)=\mathbf{1}_2$, the identity on value and slope, so any nonzero initial vector gives a mode with $\omega^2=m^2+k^2$. Check the excluded point $k=0$ separately: $T_0(L)=\left(\begin{smallmatrix}1&L\\0&1\end{smallmatrix}\right)$. Only a constant spatial mode survives for $H=+1$, and none survives for $H=-1$. If both $k=m=0$, the time equation for this one spatial mode also allows a linear function of time. No completeness theorem for mode expansions is needed to decide which separated modes exist.

\begin{exercise}[Hand the data to the next strip]
Let $d_1=d_2=L/2$, $a=-1$, $b=+1$. Start with $Y_1(0)=(1,k)^{\mathsf T}$ at $k=\pi/L$. Find $f_1$ and $f_2$ rather than inferring them from a known antiperiodic solution. What happens if the first join transmits the value but you forget to transmit the slope with the same sign?
\end{exercise}
\noindent\emph{Solution.} The first interval gives $f_1=\cos(kx_1)+\sin(kx_1)$ and $Y_1(d_1)=(1,-k)^{\mathsf T}$. Multiplying both entries by $a$ gives $Y_2(0)=(-1,k)^{\mathsf T}$. Hence $f_2=-\cos(kx_2)+\sin(kx_2)$ and $Y_2(d_2)=(1,k)^{\mathsf T}$, the original data since $b=+1$. Transmitting $(-1,-k)^{\mathsf T}$ instead already violates the stationary interface equation and returns $(-1,k)^{\mathsf T}$ at the final seam. The measurable distinction here is between two products $ba$, not between two positions for the same sign. The local charts, the metric, the incidence of the faces and the regional field equations did not change. What changed was the geometric instruction for comparing a field value after a complete circuit.


\backmatter
\chapter*{Guide to notation and constructions}
\addcontentsline{toc}{chapter}{Guide to notation and constructions}\markboth{Notation and constructions}{}
\begin{description}[font=\normalfont\bfseries,leftmargin=1.2em]
\item[Collar and coorientation.] A collar extends a face into a piece. The coordinates $r_-,r_+\geq0$ both point inward; one smooth transverse coordinate is $s=-r_-$ on the minus side and $s=r_+$ on the plus side. See Chapter~\ref{ch:01}, p.~\pageref{ch:01}.
\item[Seam identification.] $D$ identifies actual fibers or coordinates at corresponding points; it is not a coupling constant. Whether $D$ or $D^{-1}$ appears depends on the indicated source and target fibers. The three notions of lift, equivalence and stabilizer are distinguished in Chapter~\ref{ch:02}, p.~\pageref{ch:02}.
\item[Jets and traces.] $j^k_p\phi$ specifies derivatives through total order $k$ at $p$. The variational trace $q$ is the part actually varied on the interface. An infinite jet test off shell (p.~\pageref{thm:off-shell-descent}) and a finite trace test on shell (p.~\pageref{thm:branch}) have different hypotheses.
\item[Potential, form, charge.] $\Theta$ is a spacetime potential current, $\theta_\Sigma$ its slice integral with its corner correction, and $\omega_\Sigma=\delta\theta_\Sigma$. A Hamiltonian $H_X$ satisfies $\delta H_X(Y)=\omega(Y,X)$ on its domain; see p.~\pageref{eq:05-noether}. A proper gauge transformation is a redundancy, not merely one with a gauge parameter.
\item[Trace and response.] $\Pi_a$ is an \emph{outward} variational response density. At a transparent seam the two outward responses sum to zero after transport. Relative to a \emph{single} transverse coordinate this is equality of cooriented momenta; see p.~\pageref{eq:summed-momentum}.
\item[Wave inverses.] $J$ denotes the Hessian, not the causal future $J^+(K)$ of a source set. For a specified operator $P$, $G^{R/A}$ solves its bulk inhomogeneous problem and $H^{R/A}$ lifts a boundary history with no bulk source. $M$ changes a mismatch and $L$ solves for the compensating history (p.~\pageref{prop:causal}). The $H$ here is not a Hamiltonian.
\item[Bracket and state.] $E=G^R-G^A$ fixes a commutator. The symmetric covariance $\mu$ is state data; it has to pass positivity and, in the smooth-bulk field theory, the chosen short-distance condition. See p.~\pageref{eq:11-covariance}. A background solution for perturbation theory is neither of these objects.
\item[Coincidence and contact terms.] A time-ordered product at separated points does not prescribe its extension to coincident points. Equation~\eqref{eq:13-fish-extension} on p.~\pageref{eq:13-fish-extension} gives an explicit extension; Appendix~\ref{app:C} explains what its scaling degree controls.\end{description}

\chapter*{Further reading}\addcontentsline{toc}{chapter}{Further reading}\markboth{Further reading}{}
For the geometry of local fields, start with the early chapters of Giachetta, Mangiarotti and Sardanashvily~\cite{GMS}: the reason to learn a jet bundle is to give a coordinate-independent domain to a differential equation or Lagrangian. Lee~\cite{Lee} and Steenrod~\cite{Steenrod} supply background on smooth manifolds and bundle clutching. At a null face a Lorentzian reference such as O'Neill~\cite{ON} is useful precisely where Riemannian intuitions about normal vectors fail.

For boundary variations the convention used here follows Harlow and Wu~\cite{HW}. Khavkine~\cite{Khavkine} explains the dual descriptions of brackets and the conditions needed in the presence of constraints. Cattaneo, Mnev and Reshetikhin~\cite{CMR} explain why a boundary in BV theory has its own structure, while Donnelly and Freidel~\cite{DF} develop an extended description of subsystems. Neither reference licenses treating an arbitrary boundary gauge transformation as a proper gauge degeneracy. For wave operators, B\"ar, Ginoux and Pf\"affle~\cite{BGP} treat the globally hyperbolic Cauchy problem; B\"ar and Tagne Wafo~\cite{BW} treat characteristic data under additional geometric and support hypotheses. Neither is a universal timelike transmission theorem.

For the distinction between a field algebra and a choice of state, compare Wald~\cite{WaldQFT} with Haag~\cite{Haag}. After the elementary scaling calculation in Chapter~\ref{ch:13}, Brunetti and Fredenhagen~\cite{BF} and Hollands and Wald~\cite{HWick,HWT} are natural points of entry into causal renormalization and its local freedom. Their results on smooth curved spacetimes are not existence theorems for all possible physical boundaries, corners or defects.

\renewcommand{\bibname}{References}
\begin{thebibliography}{99}
\addcontentsline{toc}{chapter}{References}
\bibitem{GMS} G. Giachetta, L. Mangiarotti and G. Sardanashvily, \emph{Advanced Classical Field Theory}, World Scientific (2009).
\bibitem{Lee} J. M. Lee, \emph{Introduction to Smooth Manifolds}, 2nd ed., Springer (2013).
\bibitem{Steenrod} N. Steenrod, \emph{The Topology of Fibre Bundles}, Princeton University Press (1951).
\bibitem{ON} B. O'Neill, \emph{Semi-Riemannian Geometry with Applications to Relativity}, Academic Press (1983).
\bibitem{HW} D. Harlow and J.-q. Wu, ``Covariant phase space with boundaries,'' JHEP \textbf{10} (2020) 146, \href{https://arxiv.org/abs/1906.08616}{arXiv:1906.08616}.
\bibitem{Khavkine} I. Khavkine, ``Covariant phase space, constraints, gauge and the Peierls formula,'' Int. J. Mod. Phys. A \textbf{29} (2014) 1430009, \href{https://arxiv.org/abs/1402.1282}{arXiv:1402.1282}.
\bibitem{CMR} A. S. Cattaneo, P. Mnev and N. Reshetikhin, ``Classical BV theories on manifolds with boundary,'' Commun. Math. Phys. \textbf{332} (2014) 535--603, \href{https://arxiv.org/abs/1201.0290}{arXiv:1201.0290}.
\bibitem{DF} W. Donnelly and L. Freidel, ``Local subsystems in gauge theory and gravity,'' JHEP \textbf{09} (2016) 102, \href{https://arxiv.org/abs/1601.04744}{arXiv:1601.04744}.
\bibitem{BGP} C. B\"ar, N. Ginoux and F. Pf\"affle, \emph{Wave Equations on Lorentzian Manifolds and Quantization}, European Mathematical Society (2007).
\bibitem{BW} C. B\"ar and R. Tagne Wafo, ``Initial value problems for wave equations on manifolds,'' \emph{Math. Phys. Anal. Geom.} \textbf{18} (2015), article 7, \href{https://arxiv.org/abs/1408.4995}{arXiv:1408.4995}.
\bibitem{WaldQFT} R. M. Wald, \emph{Quantum Field Theory in Curved Spacetime and Black Hole Thermodynamics}, University of Chicago Press (1994).
\bibitem{Haag} R. Haag, \emph{Local Quantum Physics: Fields, Particles, Algebras}, 2nd ed., Springer (1996).
\bibitem{Rad} M. J. Radzikowski, ``Micro-local approach to the Hadamard condition in quantum field theory on curved space-time,'' Commun. Math. Phys. \textbf{179} (1996) 529--553.
\bibitem{BF} R. Brunetti and K. Fredenhagen, ``Microlocal analysis and interacting quantum field theories: renormalization on physical backgrounds,'' Commun. Math. Phys. \textbf{208} (2000) 623--661, \href{https://arxiv.org/abs/math-ph/9903028}{arXiv:math-ph/9903028}.
\bibitem{HWick} S. Hollands and R. M. Wald, ``Local Wick polynomials and time ordered products of quantum fields in curved spacetime,'' Commun. Math. Phys. \textbf{223} (2001) 289--326, \href{https://arxiv.org/abs/gr-qc/0103074}{arXiv:gr-qc/0103074}.
\bibitem{HWT} S. Hollands and R. M. Wald, ``Existence of local covariant time ordered products of quantum fields in curved spacetime,'' Commun. Math. Phys. \textbf{231} (2002) 309--345, \href{https://arxiv.org/abs/gr-qc/0111108}{arXiv:gr-qc/0111108}.
\bibitem{FR} K. Fredenhagen and K. Rejzner, ``Perturbative algebraic quantum field theory,'' \href{https://arxiv.org/abs/1208.1428}{arXiv:1208.1428}; see also K. Rejzner, \emph{Perturbative Algebraic Quantum Field Theory}, Springer (2016).
\end{thebibliography}

\end{document}
